%%%%%%%%%%%%%%%%%%%%%%%%%%%%%%%%%%%%%%%%%%%%%%%%%%%%%%%%%%%%%%%%%%%
% Machine Learning - Lecture Summaries
%%%%%%%%%%%%%%%%%%%%%%%%%%%%%%%%%%%%%%%%%%%%%%%%%%%%%%%%%%%%%%%%%%%

\documentclass[11pt,a4paper]{scrartcl}

% Layout
\usepackage[margin=2.5cm]{geometry}
\usepackage[onehalfspacing]{setspace}

% General packages
\usepackage{graphicx}
\usepackage{enumitem}
\usepackage{tcolorbox}
\tcbuselibrary{breakable}
\usepackage[breaklinks=true,colorlinks=true,linkcolor=blue,urlcolor=blue,citecolor=blue]{hyperref}
\usepackage{amsmath,amsthm,amssymb,mathtools}
\usepackage{icomma}
\usepackage[T1]{fontenc}
\usepackage[utf8]{inputenc}
\usepackage[ngerman]{babel}
\usepackage{tikz}

% Units / tables / floats / references
\usepackage[locale = DE,space-before-unit=true,per-mode = symbol]{siunitx}
\usepackage{booktabs,multirow}
\usepackage{placeins}
\usepackage[natbib,abbreviate=true,doi=false,style=numeric-comp,giveninits=true,sorting=none]{biblatex}
\usepackage{csquotes}
\usepackage{pgfplots}
\pgfplotsset{compat=1.18}

% Graphics paths
\graphicspath{{Bilder/} {images/} {figures/} {chapters/}}

% Optional math shortcuts
\newcommand{\R}{\mathbb{R}}
\newcommand{\N}{\mathbb{N}}
\newcommand{\Q}{\mathbb{Q}}
\newcommand{\set}[1]{\left\{ #1 \right\}}
\newcommand{\abs}[1]{\left| #1 \right|}
\newcommand{\norm}[1]{\left\lVert #1 \right\rVert}
\newcommand{\ol}[1]{\overline{#1}}
\newcommand{\half}{\frac{1}{2}}
\newcommand{\dd}{\,\mathrm{d}}

% Machine Learning shortcuts
\newcommand{\argmax}{\operatorname*{arg\,max}}
\newcommand{\E}{\mathbb{E}}
\newcommand{\Var}{\operatorname{Var}}
\newcommand{\Cov}{\operatorname{Cov}}

% Box styles
\newtcolorbox{calculationbox}[1][]{
    title=Calculation,
    breakable,
    #1
}

\newtcolorbox{solutionbox}[1][]{
    title=Solution,
    breakable,
    #1
}

\newtcolorbox{answerbox}[1][]{
    title=Answer,
    breakable,
    #1
}

\newtcolorbox{notebox}[1][]{
    title=Notes,
    breakable,
    #1
}

\newtcolorbox{sidecalcbox}[1][]{
    title=Side Calculation,
    breakable,
    #1
}

%%%%%%%%%%%%%%%%%%%%%%%%%%%%%%%%%%%%%%%%%%%%%%%%%%%%%%%%%%%%%%%%%%%
\begin{document}

    \title{Machine Learning -- Lecture Summaries}
    \subtitle{Summer Semester 2026}
    \author{Tobias Laser}
    \date{\today}
    \maketitle
    \vfill
    \thispagestyle{empty}
    \newpage

    \tableofcontents
    \newpage

    \pagenumbering{arabic}

%%%%%%%%%%%%%%%%%%%%%%%%%%%%%%%%%%%%%%%%%%%%%%%%%%%%%%%%%%%%%%%%%%%
\section{Lecture 1: Introduction and Elements of Probabilistic Modelling}
%%%%%%%%%%%%%%%%%%%%%%%%%%%%%%%%%%%%%%%%%%%%%%%%%%%%%%%%%%%%%%%%%%%

\subsection{Introduction: What should an AI system do for us?}

\begin{notebox}
    Die zentrale Frage der ersten Vorlesung ist:
    \[
        \text{Was soll ein AI-System eigentlich für uns tun?}
    \]

    Ein AI-System soll aus Daten sinnvolle Entscheidungen, Vorhersagen oder
    Schlussfolgerungen ableiten. Dabei soll es nicht nur feste Regeln ausführen,
    sondern auch mit Unsicherheit umgehen können.

    Typische Aufgaben sind zum Beispiel:
    \begin{itemize}
        \item Objekte klassifizieren,
        \item zukünftige Werte vorhersagen,
        \item Muster in Daten erkennen,
        \item Entscheidungen unter Unsicherheit treffen.
    \end{itemize}
\end{notebox}

\begin{solutionbox}
    In Machine Learning betrachtet man häufig Situationen, in denen Daten nicht
    eindeutig sind. Ein Bild, ein Messwert oder ein Text kann verrauscht,
    unvollständig oder mehrdeutig sein.

    Deshalb ist es oft nicht sinnvoll, nur eine harte Entscheidung auszugeben,
    sondern eine probabilistische Antwort:
    \[
        P(\text{Klasse} \mid \text{Daten}).
    \]

    Das bedeutet: Das System gibt an, wie wahrscheinlich verschiedene mögliche
    Erklärungen oder Klassen sind, nachdem die Daten beobachtet wurden.
\end{solutionbox}

\begin{answerbox}
    Ein AI-System soll aus Daten nützliche Entscheidungen oder Vorhersagen
    ableiten. Da reale Daten oft unsicher oder mehrdeutig sind, ist eine
    probabilistische Beschreibung sinnvoll:
    \[
        \boxed{
            \text{AI-Systeme sollen unter Unsicherheit sinnvoll entscheiden.}
        }
    \]
\end{answerbox}

%%%%%%%%%%%%%%%%%%%%%%%%%%%%%%%%%%%%%%%%%%%%%%%%%%%%%%%%%%%%%%%%%%%
\subsection{Example: Classification}
%%%%%%%%%%%%%%%%%%%%%%%%%%%%%%%%%%%%%%%%%%%%%%%%%%%%%%%%%%%%%%%%%%%

\begin{notebox}
    Bei der Klassifikation soll ein Objekt einer Klasse zugeordnet werden.

    Wir bezeichnen:
    \[
        x = \text{Daten bzw. Beobachtung}
    \]
    und
    \[
        y = \text{Klasse bzw. Label}.
    \]

    Gesucht ist die Wahrscheinlichkeit:
    \[
        P(y \mid x).
    \]
\end{notebox}

\begin{calculationbox}
    Die Größe
    \[
        P(y \mid x)
    \]
    bedeutet:

    \begin{quote}
        Wie wahrscheinlich ist die Klasse \(y\), wenn die Daten \(x\)
        beobachtet wurden?
    \end{quote}

    Ein Klassifikator entscheidet sich meistens für die Klasse mit der größten
    Wahrscheinlichkeit:
    \[
        \hat y = \argmax_y P(y \mid x).
    \]

    Dabei bedeutet \(\hat y\) die vorhergesagte Klasse.
\end{calculationbox}

\begin{solutionbox}
    Beispielhaft kann man sich eine Klassifikation von Bildern vorstellen.
    Das System erhält ein Bild \(x\) und soll entscheiden, ob darauf zum Beispiel
    eine Katze, ein Hund oder ein Auto zu sehen ist.

    Statt nur eine einzelne Antwort auszugeben, kann das System Wahrscheinlichkeiten
    für alle Klassen angeben:
    \[
        P(\text{Katze} \mid x),
        \qquad
        P(\text{Hund} \mid x),
        \qquad
        P(\text{Auto} \mid x).
    \]

    Wenn
    \[
        P(\text{Katze} \mid x)
    \]
    am größten ist, wird das Bild als Katze klassifiziert.
\end{solutionbox}

\begin{answerbox}
    Klassifikation bedeutet:
    \[
        \boxed{
            \text{Aus Daten } x \text{ wird eine Klasse } y \text{ vorhergesagt.}
        }
    \]

    Probabilistisch formuliert sucht man:
    \[
        \boxed{
            P(y \mid x)
        }
    \]

    Die Entscheidung ist:
    \[
        \boxed{
            \hat y = \argmax_y P(y \mid x)
        }
    \]
\end{answerbox}

%%%%%%%%%%%%%%%%%%%%%%%%%%%%%%%%%%%%%%%%%%%%%%%%%%%%%%%%%%%%%%%%%%%
\subsection{Different Types of Learning}
%%%%%%%%%%%%%%%%%%%%%%%%%%%%%%%%%%%%%%%%%%%%%%%%%%%%%%%%%%%%%%%%%%%

\begin{notebox}
    In Machine Learning unterscheidet man verschiedene Arten des Lernens.
    Wichtig sind vor allem:

    \begin{itemize}
        \item supervised learning,
        \item semi-supervised learning,
        \item unsupervised learning,
        \item reinforcement learning.
    \end{itemize}
\end{notebox}

\begin{solutionbox}
    Beim \emph{supervised learning} sind Trainingsdaten mit Labels gegeben.
    Das bedeutet, man hat Beispiele der Form
    \[
        (x_i,y_i),
    \]
    wobei \(x_i\) die Daten und \(y_i\) das zugehörige richtige Label ist.

    Ziel ist es, aus diesen Beispielen ein Modell zu lernen, das auch für neue
    Daten gute Vorhersagen macht.
\end{solutionbox}

\begin{solutionbox}
    Beim \emph{unsupervised learning} sind keine Labels gegeben.
    Man beobachtet nur Datenpunkte
    \[
        x_1, x_2, \dots, x_n.
    \]

    Das Ziel ist dann nicht, bekannte Labels vorherzusagen, sondern Strukturen
    in den Daten zu finden.

    Typische Aufgaben sind:
    \begin{itemize}
        \item Clustering,
        \item Dimensionsreduktion,
        \item Dichteschätzung,
        \item Finden verborgener Muster.
    \end{itemize}
\end{solutionbox}

\begin{solutionbox}
    Beim \emph{semi-supervised learning} gibt es wenige gelabelte Daten und viele
    ungelabelte Daten.

    Man hat also zum Beispiel:
    \[
        (x_1,y_1),\dots,(x_m,y_m)
    \]
    mit Labels, aber zusätzlich viele Daten
    \[
        x_{m+1},\dots,x_n
    \]
    ohne Labels.

    Die Idee ist, beide Informationsquellen zu nutzen.
\end{solutionbox}

\begin{solutionbox}
    Beim \emph{reinforcement learning} lernt ein Agent durch Interaktion mit
    einer Umgebung.

    Der Agent beobachtet einen Zustand, führt eine Aktion aus und erhält eine
    Belohnung. Ziel ist es, langfristig möglichst viel Belohnung zu sammeln.

    Man kann sich das schematisch so vorstellen:
    \[
        \text{Zustand}
        \longrightarrow
        \text{Aktion}
        \longrightarrow
        \text{Belohnung}.
    \]
\end{solutionbox}

\begin{answerbox}
    Die wichtigsten Lernarten sind:

    \begin{itemize}
        \item \textbf{Supervised learning:} Lernen aus gelabelten Daten.
        \item \textbf{Semi-supervised learning:} Lernen aus wenigen gelabelten und vielen ungelabelten Daten.
        \item \textbf{Unsupervised learning:} Finden von Struktur in ungelabelten Daten.
        \item \textbf{Reinforcement learning:} Lernen durch Belohnung und Interaktion.
    \end{itemize}
\end{answerbox}

%%%%%%%%%%%%%%%%%%%%%%%%%%%%%%%%%%%%%%%%%%%%%%%%%%%%%%%%%%%%%%%%%%%
\subsection{Probabilistic Answers as Optimal Inference}
%%%%%%%%%%%%%%%%%%%%%%%%%%%%%%%%%%%%%%%%%%%%%%%%%%%%%%%%%%%%%%%%%%%

\begin{notebox}
    Ein wichtiger Gedanke der Vorlesung ist:

    \begin{quote}
        Probabilistische Antworten können als optimale Inferenz verstanden werden.
    \end{quote}

    Inferenz bedeutet, aus Beobachtungen auf unbekannte Größen zu schließen.
\end{notebox}

\begin{solutionbox}
    Wenn ein AI-System Daten \(x\) beobachtet, möchte es auf eine unbekannte
    Größe \(y\) schließen. Das kann zum Beispiel eine Klasse, ein zukünftiger
    Wert oder eine verborgene Ursache sein.

    Statt nur einen einzelnen Wert auszugeben, beschreibt ein probabilistisches
    Modell die Unsicherheit über \(y\) durch eine Wahrscheinlichkeitsverteilung:
    \[
        P(y \mid x).
    \]

    Diese Verteilung fasst zusammen, welche Werte von \(y\) nach Beobachtung von
    \(x\) plausibel sind.
\end{solutionbox}

\begin{answerbox}
    Probabilistische Inferenz bedeutet:
    \[
        \boxed{
            \text{Aus beobachteten Daten } x
            \text{ wird auf unbekannte Größen } y \text{ geschlossen.}
        }
    \]

    Das zentrale Objekt ist:
    \[
        \boxed{
            P(y \mid x)
        }
    \]
\end{answerbox}

%%%%%%%%%%%%%%%%%%%%%%%%%%%%%%%%%%%%%%%%%%%%%%%%%%%%%%%%%%%%%%%%%%%
\subsection{Joint Probabilities and Conditional Probabilities}
%%%%%%%%%%%%%%%%%%%%%%%%%%%%%%%%%%%%%%%%%%%%%%%%%%%%%%%%%%%%%%%%%%%

\begin{notebox}
    Die gemeinsame Wahrscheinlichkeit beschreibt, wie wahrscheinlich es ist,
    dass zwei Ereignisse gleichzeitig eintreten:
    \[
        P(A,B).
    \]

    Beispiel:
    \[
        P(\text{krank}, \text{positiver Test}).
    \]

    Das ist die Wahrscheinlichkeit, dass eine Person krank ist und gleichzeitig
    einen positiven Test hat.
\end{notebox}

\begin{notebox}
    Die bedingte Wahrscheinlichkeit beschreibt, wie wahrscheinlich ein Ereignis
    ist, wenn ein anderes Ereignis bereits bekannt ist:
    \[
        P(A \mid B).
    \]

    Beispiel:
    \[
        P(\text{krank} \mid \text{positiver Test}).
    \]

    Das bedeutet:
    \begin{quote}
        Wie wahrscheinlich ist es, krank zu sein, wenn der Test positiv ist?
    \end{quote}
\end{notebox}

\begin{calculationbox}
    Die bedingte Wahrscheinlichkeit ist definiert durch:
    \[
        P(A \mid B)
        =
        \frac{P(A,B)}{P(B)}
    \]
    für \(P(B) > 0\).

    Umgestellt erhält man:
    \[
        P(A,B)
        =
        P(A \mid B)P(B).
    \]

    Diese Gleichung ist bereits die Produktregel.
\end{calculationbox}

\begin{answerbox}
    Wichtige Begriffe:

    \[
        \boxed{
            P(A,B)
            =
            \text{gemeinsame Wahrscheinlichkeit von } A \text{ und } B
        }
    \]

    \[
        \boxed{
            P(A \mid B)
            =
            \text{Wahrscheinlichkeit von } A \text{ gegeben } B
        }
    \]
\end{answerbox}

%%%%%%%%%%%%%%%%%%%%%%%%%%%%%%%%%%%%%%%%%%%%%%%%%%%%%%%%%%%%%%%%%%%
\subsection{Sum Rule and Marginal Probabilities}
%%%%%%%%%%%%%%%%%%%%%%%%%%%%%%%%%%%%%%%%%%%%%%%%%%%%%%%%%%%%%%%%%%%

\begin{calculationbox}
    Die Summenregel erlaubt es, eine Variable aus einer gemeinsamen
    Wahrscheinlichkeit zu eliminieren.

    Für diskrete Zufallsvariablen gilt:
    \[
        P(A) = \sum_B P(A,B).
    \]

    Das nennt man \emph{Marginalisierung}.
\end{calculationbox}

\begin{solutionbox}
    Die Idee ist: Wenn \(B\) verschiedene mögliche Werte annehmen kann, summiert
    man über alle Möglichkeiten von \(B\), um nur noch die Wahrscheinlichkeit von
    \(A\) zu erhalten.

    Hat \(B\) zum Beispiel die möglichen Werte \(b_1,b_2,\dots,b_n\), dann gilt:
    \[
        P(A)
        =
        \sum_{i=1}^{n} P(A,b_i).
    \]
\end{solutionbox}

\begin{answerbox}
    Die Summenregel lautet:
    \[
        \boxed{
            P(A) = \sum_B P(A,B)
        }
    \]

    Sie wird verwendet, um aus einer gemeinsamen Wahrscheinlichkeit eine
    Randwahrscheinlichkeit zu erhalten.
\end{answerbox}

%%%%%%%%%%%%%%%%%%%%%%%%%%%%%%%%%%%%%%%%%%%%%%%%%%%%%%%%%%%%%%%%%%%
\subsection{Product Rule}
%%%%%%%%%%%%%%%%%%%%%%%%%%%%%%%%%%%%%%%%%%%%%%%%%%%%%%%%%%%%%%%%%%%

\begin{calculationbox}
    Die Produktregel verbindet gemeinsame und bedingte Wahrscheinlichkeiten:
    \[
        P(A,B) = P(A \mid B)P(B).
    \]

    Genauso gilt:
    \[
        P(A,B) = P(B \mid A)P(A).
    \]

    Beide Ausdrücke beschreiben dieselbe gemeinsame Wahrscheinlichkeit.
\end{calculationbox}

\begin{sidecalcbox}
    Man erhält die Produktregel direkt aus der Definition der bedingten
    Wahrscheinlichkeit:
    \[
        P(A \mid B)
        =
        \frac{P(A,B)}{P(B)}.
    \]

    Multiplikation mit \(P(B)\) liefert:
    \[
        P(A \mid B)P(B)
        =
        P(A,B).
    \]

    Also:
    \[
        P(A,B)
        =
        P(A \mid B)P(B).
    \]
\end{sidecalcbox}

\begin{answerbox}
    Die Produktregel lautet:
    \[
        \boxed{
            P(A,B) = P(A \mid B)P(B)
        }
    \]

    und äquivalent:
    \[
        \boxed{
            P(A,B) = P(B \mid A)P(A)
        }
    \]
\end{answerbox}

%%%%%%%%%%%%%%%%%%%%%%%%%%%%%%%%%%%%%%%%%%%%%%%%%%%%%%%%%%%%%%%%%%%
\subsection{Bayes' Rule}
%%%%%%%%%%%%%%%%%%%%%%%%%%%%%%%%%%%%%%%%%%%%%%%%%%%%%%%%%%%%%%%%%%%

\begin{calculationbox}
    Aus der Produktregel erhält man:
    \[
        P(A,B) = P(A \mid B)P(B)
    \]
    und auch:
    \[
        P(A,B) = P(B \mid A)P(A).
    \]

    Da beide Seiten dieselbe gemeinsame Wahrscheinlichkeit beschreiben, gilt:
    \[
        P(A \mid B)P(B) = P(B \mid A)P(A).
    \]

    Durch Umformen nach \(P(A \mid B)\) erhält man Bayes' Regel:
    \[
        P(A \mid B)
        =
        \frac{P(B \mid A)P(A)}{P(B)}.
    \]
\end{calculationbox}

\begin{calculationbox}
    Für Klassifikationsprobleme schreibt man häufig:
    \[
        P(y \mid x)
        =
        \frac{P(x \mid y)P(y)}{P(x)}.
    \]

    Dabei ist:
    \[
        P(y \mid x)
        \quad \text{Posterior,}
    \]
    also die Wahrscheinlichkeit der Klasse \(y\) nach Beobachtung der Daten \(x\).

    \[
        P(x \mid y)
        \quad \text{Likelihood,}
    \]
    also die Wahrscheinlichkeit, die Daten \(x\) zu beobachten, wenn die Klasse
    \(y\) gilt.

    \[
        P(y)
        \quad \text{Prior,}
    \]
    also das Vorwissen über die Klasse \(y\).

    \[
        P(x)
        \quad \text{Evidence,}
    \]
    also der Normalisierungsfaktor.
\end{calculationbox}

\begin{answerbox}
    Bayes' Regel lautet:
    \[
        \boxed{
            P(A \mid B)
            =
            \frac{P(B \mid A)P(A)}{P(B)}
        }
    \]

    Für Klassifikation:
    \[
        \boxed{
            P(y \mid x)
            =
            \frac{P(x \mid y)P(y)}{P(x)}
        }
    \]

    Bayes' Regel erlaubt es, aus beobachteten Daten auf die wahrscheinlichste
    Ursache oder Klasse zu schließen.
\end{answerbox}

%%%%%%%%%%%%%%%%%%%%%%%%%%%%%%%%%%%%%%%%%%%%%%%%%%%%%%%%%%%%%%%%%%%
\subsection{Introductory Example: Pick-One-Berry-From-Two-Bowls}
%%%%%%%%%%%%%%%%%%%%%%%%%%%%%%%%%%%%%%%%%%%%%%%%%%%%%%%%%%%%%%%%%%%

\begin{solutionbox}
    Ein typisches Einstiegsbeispiel ist folgendes:

    Es gibt zwei Schüsseln mit unterschiedlichen Mischungen aus Beeren.
    Man zieht eine Beere und möchte daraus schließen, aus welcher Schüssel sie
    wahrscheinlich stammt.

    Die relevante Wahrscheinlichkeit ist:
    \[
        P(\text{Schüssel} \mid \text{Beere}).
    \]

    Man beobachtet also die Beere und schließt rückwärts auf die wahrscheinlichere
    Ursache, nämlich die Schüssel.
\end{solutionbox}

\begin{calculationbox}
    Mit Bayes' Regel erhält man:
    \[
        P(\text{Schüssel} \mid \text{Beere})
        =
        \frac{
            P(\text{Beere} \mid \text{Schüssel})
            P(\text{Schüssel})
        }{
            P(\text{Beere})
        }.
    \]

    Die Interpretation ist:
    \begin{itemize}
        \item \(P(\text{Schüssel})\) beschreibt das Vorwissen darüber, wie
        wahrscheinlich die jeweilige Schüssel ist.
        \item \(P(\text{Beere} \mid \text{Schüssel})\) beschreibt die
        Wahrscheinlichkeit, diese Beere aus dieser Schüssel zu ziehen.
        \item \(P(\text{Schüssel} \mid \text{Beere})\) ist die gesuchte
        Wahrscheinlichkeit nach Beobachtung der Beere.
    \end{itemize}
\end{calculationbox}

\begin{answerbox}
    Das Beeren-Schüssel-Beispiel ist ein Bayes-Beispiel:

    \[
        \boxed{
            \text{Beobachtung: Beere}
            \quad \Longrightarrow \quad
            \text{gesucht: wahrscheinliche Schüssel}
        }
    \]

    Mathematisch:
    \[
        \boxed{
            P(\text{Schüssel} \mid \text{Beere})
            =
            \frac{
                P(\text{Beere} \mid \text{Schüssel})
                P(\text{Schüssel})
            }{
                P(\text{Beere})
            }
        }
    \]
\end{answerbox}

%%%%%%%%%%%%%%%%%%%%%%%%%%%%%%%%%%%%%%%%%%%%%%%%%%%%%%%%%%%%%%%%%%%
\subsection{Example Algorithms: Outlook}
%%%%%%%%%%%%%%%%%%%%%%%%%%%%%%%%%%%%%%%%%%%%%%%%%%%%%%%%%%%%%%%%%%%

\begin{notebox}
    Die erste Vorlesung gibt außerdem einen Ausblick auf Algorithmen, die später
    in der Lehrveranstaltung behandelt werden können.

    Die Grundidee ist: Viele Machine-Learning-Algorithmen lassen sich als
    Verfahren verstehen, die Wahrscheinlichkeiten, Parameter oder Strukturen aus
    Daten lernen.
\end{notebox}

\begin{solutionbox}
    Beispiele für spätere Modellklassen und Algorithmen sind:

    \begin{itemize}
        \item Regressionsmodelle,
        \item logistische Regression,
        \item neuronale Netze,
        \item PCA zur Dimensionsreduktion,
        \item Clustering-Verfahren,
        \item Gaussian Mixture Models,
        \item Expectation-Maximization.
    \end{itemize}

    Die probabilistische Denkweise aus Lecture 1 bildet eine wichtige Grundlage
    für viele dieser späteren Methoden.
\end{solutionbox}

%%%%%%%%%%%%%%%%%%%%%%%%%%%%%%%%%%%%%%%%%%%%%%%%%%%%%%%%%%%%%%%%%%%
\subsection{Evidence of Strong Models in Biological Vision}
%%%%%%%%%%%%%%%%%%%%%%%%%%%%%%%%%%%%%%%%%%%%%%%%%%%%%%%%%%%%%%%%%%%

\begin{solutionbox}
    Die Vorlesung erwähnt auch die Adelson's Checkerboard Illusion.
    Diese optische Täuschung zeigt, dass Wahrnehmung nicht einfach nur rohe
    Bilddaten verarbeitet.

    Das Gehirn interpretiert Sinneseindrücke unter Annahmen über Licht, Schatten
    und Objekte. Zwei Felder können physikalisch dieselbe Helligkeit haben, aber
    unterschiedlich wahrgenommen werden, weil das Gehirn Kontextinformationen
    einbezieht.
\end{solutionbox}

\begin{answerbox}
    Die Botschaft für Machine Learning ist:

    \begin{quote}
        Auch biologische Intelligenz arbeitet mit starken Modellen der Welt und
        interpretiert Beobachtungen nicht rein passiv.
    \end{quote}

    Das motiviert probabilistische Modelle: Ein gutes System nutzt Daten und
    Annahmen, um plausible Schlussfolgerungen zu ziehen.
\end{answerbox}

%%%%%%%%%%%%%%%%%%%%%%%%%%%%%%%%%%%%%%%%%%%%%%%%%%%%%%%%%%%%%%%%%%%
\subsection{Prüfungsrelevante Kurzfassung}
%%%%%%%%%%%%%%%%%%%%%%%%%%%%%%%%%%%%%%%%%%%%%%%%%%%%%%%%%%%%%%%%%%%

\begin{answerbox}
    Lecture 1 führt Machine Learning als Entscheidungsfindung unter Unsicherheit
    ein.

    Besonders wichtig ist Klassifikation:
    \[
        x \longmapsto y.
    \]

    Probabilistisch formuliert sucht man:
    \[
        P(y \mid x).
    \]

    Dafür braucht man:
    \begin{itemize}
        \item gemeinsame Wahrscheinlichkeiten \(P(A,B)\),
        \item bedingte Wahrscheinlichkeiten \(P(A \mid B)\),
        \item die Summenregel,
        \item die Produktregel,
        \item Bayes' Regel.
    \end{itemize}

    Die wichtigste Formel ist:
    \[
        \boxed{
            P(y \mid x)
            =
            \frac{P(x \mid y)P(y)}{P(x)}
        }
    \]

    Bayes' Regel ist zentral, weil sie erlaubt, aus beobachteten Daten auf
    wahrscheinliche Ursachen oder Klassen zu schließen.
\end{answerbox}

%%%%%%%%%%%%%%%%%%%%%%%%%%%%%%%%%%%%%%%%%%%%%%%%%%%%%%%%%%%%%%%%%%%
\section{Lecture 2: Probabilistic Calculus}
%%%%%%%%%%%%%%%%%%%%%%%%%%%%%%%%%%%%%%%%%%%%%%%%%%%%%%%%%%%%%%%%%%%

\subsection{Overview}

\begin{notebox}
    Lecture 2 erweitert die probabilistischen Grundlagen aus Lecture 1.

    Die zentralen Themen sind:

    \begin{itemize}
        \item allgemeine Form von Produktregel, Summenregel und Bayes' Regel,
        \item kontinuierliche Zufallsvariablen und Wahrscheinlichkeitsdichten,
        \item Produktregel, Summenregel und Bayes' Regel für kontinuierliche Variablen,
        \item Dichteschätzung mit Histogrammen,
        \item Erwartungswerte,
        \item Approximation von Erwartungswerten durch Sampling,
        \item Mittelwert, Varianz, Kovarianz und Korrelation,
        \item wichtige Wahrscheinlichkeitsverteilungen.
    \end{itemize}
\end{notebox}

\begin{answerbox}
    Die Kernaussage von Lecture 2 ist:

    \[
        \boxed{
            \text{Wahrscheinlichkeitsrechnung liefert die Rechenregeln,
            mit denen Unsicherheit mathematisch verarbeitet wird.}
        }
    \]

    Diese Regeln bilden die Grundlage für probabilistische Modelle im Machine
    Learning.
\end{answerbox}

%%%%%%%%%%%%%%%%%%%%%%%%%%%%%%%%%%%%%%%%%%%%%%%%%%%%%%%%%%%%%%%%%%%
\subsection{Generality of the Product Rule, Sum Rule and Bayes' Rule}
%%%%%%%%%%%%%%%%%%%%%%%%%%%%%%%%%%%%%%%%%%%%%%%%%%%%%%%%%%%%%%%%%%%

\begin{notebox}
    Die drei wichtigsten Regeln aus Lecture 1 gelten nicht nur für das
    Beeren-Beispiel, sondern allgemein für Zufallsvariablen und Ereignisse.

    Die drei Grundregeln sind:

    \begin{itemize}
        \item Produktregel,
        \item Summenregel,
        \item Bayes' Regel.
    \end{itemize}
\end{notebox}

\begin{calculationbox}
    Für zwei Zufallsvariablen \(X\) und \(Y\) beschreibt die gemeinsame
    Wahrscheinlichkeit
    \[
        P(X,Y)
    \]
    die Wahrscheinlichkeit, dass \(X\) und \(Y\) gemeinsam bestimmte Werte
    annehmen.

    Die bedingte Wahrscheinlichkeit ist definiert durch:
    \[
        P(X \mid Y)
        =
        \frac{P(X,Y)}{P(Y)}
    \]
    für \(P(Y)>0\).

    Daraus folgt durch Umstellen die Produktregel:
    \[
        P(X,Y)
        =
        P(X \mid Y)P(Y).
    \]

    Genauso gilt:
    \[
        P(X,Y)
        =
        P(Y \mid X)P(X).
    \]
\end{calculationbox}

\begin{calculationbox}
    Da beide Ausdrücke dieselbe gemeinsame Wahrscheinlichkeit beschreiben,
    gilt:
    \[
        P(X \mid Y)P(Y)
        =
        P(Y \mid X)P(X).
    \]

    Auflösen nach \(P(X \mid Y)\) ergibt Bayes' Regel:
    \[
        P(X \mid Y)
        =
        \frac{P(Y \mid X)P(X)}{P(Y)}.
    \]
\end{calculationbox}

\begin{calculationbox}
    Die Summenregel erlaubt es, eine Variable aus einer gemeinsamen
    Wahrscheinlichkeit zu entfernen.

    Für diskrete Zufallsvariablen gilt:
    \[
        P(X)
        =
        \sum_Y P(X,Y).
    \]

    Das nennt man Marginalisierung.
\end{calculationbox}

\begin{answerbox}
    Die wichtigsten diskreten Regeln sind:

    \[
        \boxed{
            P(X,Y)
            =
            P(X \mid Y)P(Y)
        }
    \]

    \[
        \boxed{
            P(X)
            =
            \sum_Y P(X,Y)
        }
    \]

    \[
        \boxed{
            P(X \mid Y)
            =
            \frac{P(Y \mid X)P(X)}{P(Y)}
        }
    \]
\end{answerbox}

%%%%%%%%%%%%%%%%%%%%%%%%%%%%%%%%%%%%%%%%%%%%%%%%%%%%%%%%%%%%%%%%%%%
\subsection{Continuous Random Variables and Probability Density Functions}
%%%%%%%%%%%%%%%%%%%%%%%%%%%%%%%%%%%%%%%%%%%%%%%%%%%%%%%%%%%%%%%%%%%

\begin{notebox}
    Viele Größen in Machine Learning sind nicht diskret, sondern kontinuierlich.

    Beispiele:

    \begin{itemize}
        \item Körpergröße,
        \item Temperatur,
        \item Gewicht,
        \item Zeit,
        \item Messwerte,
        \item Positionen in einem Bild,
        \item Parameter eines Modells.
    \end{itemize}

    Eine kontinuierliche Zufallsvariable kann unendlich viele Werte in einem
    Intervall annehmen.
\end{notebox}

\begin{solutionbox}
    Bei diskreten Zufallsvariablen kann man einzelnen Werten direkte
    Wahrscheinlichkeiten zuordnen:
    \[
        P(X=x).
    \]

    Bei kontinuierlichen Zufallsvariablen ist die Wahrscheinlichkeit für einen
    exakt einzelnen Wert normalerweise:
    \[
        P(X=x)=0.
    \]

    Stattdessen betrachtet man Wahrscheinlichkeiten für Intervalle:
    \[
        P(a \leq X \leq b).
    \]

    Diese werden über eine Wahrscheinlichkeitsdichte \(p(x)\) berechnet:
    \[
        P(a \leq X \leq b)
        =
        \int_a^b p(x)\,dx.
    \]
\end{solutionbox}

\begin{calculationbox}
    Eine Wahrscheinlichkeitsdichte \(p(x)\) muss zwei Eigenschaften erfüllen:

    \[
        p(x) \geq 0
    \]
    für alle \(x\), und

    \[
        \int_{-\infty}^{\infty} p(x)\,dx
        =
        1.
    \]

    Die Gesamtfläche unter der Dichtefunktion ist also \(1\).
\end{calculationbox}

\begin{answerbox}
    Für kontinuierliche Zufallsvariablen gilt:

    \[
        \boxed{
            P(a \leq X \leq b)
            =
            \int_a^b p(x)\,dx
        }
    \]

    Die Funktion \(p(x)\) ist keine Wahrscheinlichkeit an einem Punkt, sondern
    eine Dichte.

    Wichtig:
    \[
        \boxed{
            P(X=x)=0
            \quad \text{für kontinuierliche } X.
        }
    \]
\end{answerbox}

%%%%%%%%%%%%%%%%%%%%%%%%%%%%%%%%%%%%%%%%%%%%%%%%%%%%%%%%%%%%%%%%%%%
\subsection{Product Rule, Sum Rule and Bayes' Rule for Continuous Random Variables}
%%%%%%%%%%%%%%%%%%%%%%%%%%%%%%%%%%%%%%%%%%%%%%%%%%%%%%%%%%%%%%%%%%%

\begin{notebox}
    Die bekannten Rechenregeln gelten auch für kontinuierliche Zufallsvariablen.

    Der Unterschied ist:

    \begin{itemize}
        \item Summen werden durch Integrale ersetzt.
        \item Wahrscheinlichkeiten werden durch Dichten ersetzt.
    \end{itemize}
\end{notebox}

\begin{calculationbox}
    Für kontinuierliche Zufallsvariablen \(X\) und \(Y\) verwendet man die
    gemeinsame Dichte
    \[
        p(x,y).
    \]

    Die bedingte Dichte ist:
    \[
        p(x \mid y)
        =
        \frac{p(x,y)}{p(y)}.
    \]

    Daraus folgt die Produktregel:
    \[
        p(x,y)
        =
        p(x \mid y)p(y).
    \]

    Ebenso gilt:
    \[
        p(x,y)
        =
        p(y \mid x)p(x).
    \]
\end{calculationbox}

\begin{calculationbox}
    Die kontinuierliche Summenregel lautet:
    \[
        p(x)
        =
        \int p(x,y)\,dy.
    \]

    Mit der Produktregel:
    \[
        p(x)
        =
        \int p(x \mid y)p(y)\,dy.
    \]

    Das Integral marginalisiert die Variable \(Y\) heraus.
\end{calculationbox}

\begin{calculationbox}
    Bayes' Regel für kontinuierliche Zufallsvariablen lautet:
    \[
        p(x \mid y)
        =
        \frac{p(y \mid x)p(x)}{p(y)}.
    \]

    Der Nenner kann wieder durch Marginalisierung berechnet werden:
    \[
        p(y)
        =
        \int p(y \mid x)p(x)\,dx.
    \]

    Damit:
    \[
        p(x \mid y)
        =
        \frac{p(y \mid x)p(x)}
        {\int p(y \mid x)p(x)\,dx}.
    \]
\end{calculationbox}

\begin{answerbox}
    Für kontinuierliche Zufallsvariablen gelten die analogen Regeln:

    \[
        \boxed{
            p(x,y)
            =
            p(x \mid y)p(y)
        }
    \]

    \[
        \boxed{
            p(x)
            =
            \int p(x,y)\,dy
        }
    \]

    \[
        \boxed{
            p(x \mid y)
            =
            \frac{p(y \mid x)p(x)}{p(y)}
        }
    \]

    Der wichtigste Unterschied zur diskreten Version ist:
    \[
        \boxed{
            \sum \quad \longrightarrow \quad \int
        }
    \]
\end{answerbox}

%%%%%%%%%%%%%%%%%%%%%%%%%%%%%%%%%%%%%%%%%%%%%%%%%%%%%%%%%%%%%%%%%%%
\subsection{Density Estimation Using Histograms}
%%%%%%%%%%%%%%%%%%%%%%%%%%%%%%%%%%%%%%%%%%%%%%%%%%%%%%%%%%%%%%%%%%%

\begin{notebox}
    In vielen Anwendungen kennt man die wahre Dichte \(p(x)\) nicht.

    Man hat nur Daten:
    \[
        x^{(1)},x^{(2)},\dots,x^{(N)}.
    \]

    Aus diesen Daten möchte man eine Schätzung für die zugrundeliegende Dichte
    konstruieren. Das nennt man \emph{density estimation}.
\end{notebox}

\begin{solutionbox}
    Eine einfache Methode zur Dichteschätzung ist ein Histogramm.

    Man teilt den Wertebereich in Intervalle, sogenannte Bins, ein.

    Für jeden Bin zählt man, wie viele Datenpunkte in diesem Intervall liegen.

    Wenn ein Bin die Breite \(\Delta\) hat und \(n_i\) Datenpunkte enthält,
    dann kann man die Dichte in diesem Bin approximieren durch:
    \[
        \hat p_i
        =
        \frac{n_i}{N\Delta}.
    \]

    Dabei ist:

    \begin{itemize}
        \item \(n_i\) die Anzahl der Datenpunkte im Bin,
        \item \(N\) die Gesamtzahl der Datenpunkte,
        \item \(\Delta\) die Breite des Bins.
    \end{itemize}
\end{solutionbox}

\begin{calculationbox}
    Die Normierung funktioniert, weil:
    \[
        \sum_i \hat p_i \Delta
        =
        \sum_i \frac{n_i}{N\Delta}\Delta.
    \]

    Kürzen von \(\Delta\) ergibt:
    \[
        \sum_i \hat p_i \Delta
        =
        \frac{1}{N}\sum_i n_i.
    \]

    Da alle Datenpunkte auf die Bins verteilt werden, gilt:
    \[
        \sum_i n_i = N.
    \]

    Also:
    \[
        \sum_i \hat p_i \Delta
        =
        1.
    \]
\end{calculationbox}

\begin{answerbox}
    Histogramme schätzen eine Dichte durch Zählen von Datenpunkten in Intervallen:

    \[
        \boxed{
            \hat p_i
            =
            \frac{n_i}{N\Delta}
        }
    \]

    Kleine Bin-Breiten liefern detaillierte, aber oft verrauschte Schätzungen.

    Große Bin-Breiten liefern glattere, aber gröbere Schätzungen.
\end{answerbox}

%%%%%%%%%%%%%%%%%%%%%%%%%%%%%%%%%%%%%%%%%%%%%%%%%%%%%%%%%%%%%%%%%%%
\subsection{Expectation Values}
%%%%%%%%%%%%%%%%%%%%%%%%%%%%%%%%%%%%%%%%%%%%%%%%%%%%%%%%%%%%%%%%%%%

\begin{notebox}
    Der Erwartungswert beschreibt den durchschnittlichen Wert einer Zufallsvariable
    unter ihrer Wahrscheinlichkeitsverteilung.

    Man schreibt:
    \[
        \mathbb{E}[X].
    \]
\end{notebox}

\begin{calculationbox}
    Für eine diskrete Zufallsvariable \(X\) mit möglichen Werten \(x_i\) gilt:
    \[
        \mathbb{E}[X]
        =
        \sum_i x_i P(X=x_i).
    \]

    Allgemeiner für eine Funktion \(f(X)\):
    \[
        \mathbb{E}[f(X)]
        =
        \sum_i f(x_i)P(X=x_i).
    \]
\end{calculationbox}

\begin{calculationbox}
    Für eine kontinuierliche Zufallsvariable mit Dichte \(p(x)\) gilt:
    \[
        \mathbb{E}[X]
        =
        \int x p(x)\,dx.
    \]

    Allgemeiner:
    \[
        \mathbb{E}[f(X)]
        =
        \int f(x)p(x)\,dx.
    \]
\end{calculationbox}

\begin{answerbox}
    Der Erwartungswert ist ein gewichteter Mittelwert.

    Diskret:
    \[
        \boxed{
            \mathbb{E}[f(X)]
            =
            \sum_i f(x_i)P(X=x_i)
        }
    \]

    Kontinuierlich:
    \[
        \boxed{
            \mathbb{E}[f(X)]
            =
            \int f(x)p(x)\,dx
        }
    \]
\end{answerbox}

%%%%%%%%%%%%%%%%%%%%%%%%%%%%%%%%%%%%%%%%%%%%%%%%%%%%%%%%%%%%%%%%%%%
\subsection{Sampling Approximation of Expectation Values}
%%%%%%%%%%%%%%%%%%%%%%%%%%%%%%%%%%%%%%%%%%%%%%%%%%%%%%%%%%%%%%%%%%%

\begin{notebox}
    Erwartungswerte können oft nicht exakt berechnet werden.

    Wenn man aber Samples
    \[
        x^{(1)},x^{(2)},\dots,x^{(N)}
    \]
    aus der Verteilung \(p(x)\) hat, kann man Erwartungswerte durch Mittelwerte
    über die Samples approximieren.
\end{notebox}

\begin{calculationbox}
    Der exakte Erwartungswert einer Funktion \(f(X)\) ist:
    \[
        \mathbb{E}[f(X)]
        =
        \int f(x)p(x)\,dx.
    \]

    Mit Samples \(x^{(1)},\dots,x^{(N)}\) approximiert man:
    \[
        \mathbb{E}[f(X)]
        \approx
        \frac{1}{N}
        \sum_{n=1}^{N}
        f(x^{(n)}).
    \]

    Diese Formel ist die Sampling-Approximation des Erwartungswertes.
\end{calculationbox}

\begin{solutionbox}
    Die Idee ist:

    Werte, die unter \(p(x)\) häufig auftreten, kommen auch häufiger in den
    Samples vor.

    Dadurch ersetzt die Summe über Samples näherungsweise das Integral über die
    Dichte.

    Für große \(N\) wird die Approximation typischerweise besser.
\end{solutionbox}

\begin{answerbox}
    Die Sampling-Approximation lautet:

    \[
        \boxed{
            \mathbb{E}[f(X)]
            \approx
            \frac{1}{N}
            \sum_{n=1}^{N}
            f(x^{(n)})
        }
    \]

    Sie ist zentral für Monte-Carlo-Methoden und viele numerische Verfahren im
    Machine Learning.
\end{answerbox}

%%%%%%%%%%%%%%%%%%%%%%%%%%%%%%%%%%%%%%%%%%%%%%%%%%%%%%%%%%%%%%%%%%%
\subsection{Mean, Variance, Covariance and Correlation}
%%%%%%%%%%%%%%%%%%%%%%%%%%%%%%%%%%%%%%%%%%%%%%%%%%%%%%%%%%%%%%%%%%%

\begin{notebox}
    Erwartungswerte werden verwendet, um wichtige Kenngrößen von Verteilungen zu
    definieren:

    \begin{itemize}
        \item mean,
        \item variance,
        \item covariance,
        \item correlation.
    \end{itemize}
\end{notebox}

\begin{calculationbox}
    Der Mittelwert oder Erwartungswert einer Zufallsvariable \(X\) ist:
    \[
        \mu
        =
        \mathbb{E}[X].
    \]

    Für kontinuierliche Variablen:
    \[
        \mu
        =
        \int x p(x)\,dx.
    \]
\end{calculationbox}

\begin{calculationbox}
    Die Varianz misst die mittlere quadratische Abweichung vom Mittelwert:
    \[
        \operatorname{Var}(X)
        =
        \mathbb{E}\left[(X-\mu)^2\right].
    \]

    Äquivalent gilt:
    \[
        \operatorname{Var}(X)
        =
        \mathbb{E}[X^2]
        -
        \mathbb{E}[X]^2.
    \]

    Die Standardabweichung ist:
    \[
        \sigma
        =
        \sqrt{\operatorname{Var}(X)}.
    \]
\end{calculationbox}

\begin{calculationbox}
    Die Kovarianz zweier Zufallsvariablen \(X\) und \(Y\) ist:
    \[
        \operatorname{Cov}(X,Y)
        =
        \mathbb{E}\left[(X-\mu_X)(Y-\mu_Y)\right].
    \]

    Äquivalent:
    \[
        \operatorname{Cov}(X,Y)
        =
        \mathbb{E}[XY]
        -
        \mathbb{E}[X]\mathbb{E}[Y].
    \]

    Sie misst, ob zwei Variablen gemeinsam größer oder kleiner werden.
\end{calculationbox}

\begin{calculationbox}
    Die Korrelation ist die normierte Kovarianz:
    \[
        \rho_{X,Y}
        =
        \frac{\operatorname{Cov}(X,Y)}
        {\sigma_X\sigma_Y}.
    \]

    Dabei gilt:
    \[
        -1 \leq \rho_{X,Y} \leq 1.
    \]

    Interpretation:
    \begin{itemize}
        \item \(\rho_{X,Y} > 0\): positive lineare Beziehung,
        \item \(\rho_{X,Y} < 0\): negative lineare Beziehung,
        \item \(\rho_{X,Y} = 0\): keine lineare Korrelation.
    \end{itemize}
\end{calculationbox}

\begin{answerbox}
    Die wichtigsten Größen sind:

    \[
        \boxed{
            \mu
            =
            \mathbb{E}[X]
        }
    \]

    \[
        \boxed{
            \operatorname{Var}(X)
            =
            \mathbb{E}\left[(X-\mu)^2\right]
            =
            \mathbb{E}[X^2] - \mathbb{E}[X]^2
        }
    \]

    \[
        \boxed{
            \operatorname{Cov}(X,Y)
            =
            \mathbb{E}[XY]
            -
            \mathbb{E}[X]\mathbb{E}[Y]
        }
    \]

    \[
        \boxed{
            \rho_{X,Y}
            =
            \frac{\operatorname{Cov}(X,Y)}
            {\sigma_X\sigma_Y}
        }
    \]
\end{answerbox}

%%%%%%%%%%%%%%%%%%%%%%%%%%%%%%%%%%%%%%%%%%%%%%%%%%%%%%%%%%%%%%%%%%%
\subsection{Important Probability Distributions}
%%%%%%%%%%%%%%%%%%%%%%%%%%%%%%%%%%%%%%%%%%%%%%%%%%%%%%%%%%%%%%%%%%%

\begin{notebox}
    In Machine Learning treten bestimmte Wahrscheinlichkeitsverteilungen immer
    wieder auf.

    Sie dienen als Modelle für Daten, Unsicherheit, Fehler oder Parameter.
\end{notebox}

\begin{solutionbox}
    Die \textbf{Bernoulli-Verteilung} modelliert ein binäres Ereignis:
    \[
        X \in \{0,1\}.
    \]

    Mit Parameter \(p\) gilt:
    \[
        P(X=1)=p,
        \qquad
        P(X=0)=1-p.
    \]

    Kompakter:
    \[
        P(X=x)
        =
        p^x(1-p)^{1-x},
        \qquad x \in \{0,1\}.
    \]
\end{solutionbox}

\begin{solutionbox}
    Die \textbf{kategoriale Verteilung} modelliert eine Zufallsvariable mit
    mehreren möglichen Klassen:
    \[
        X \in \{1,\dots,K\}.
    \]

    Für Klassenwahrscheinlichkeiten \(\pi_1,\dots,\pi_K\) gilt:
    \[
        P(X=k)
        =
        \pi_k
    \]
    mit
    \[
        \sum_{k=1}^K \pi_k = 1.
    \]
\end{solutionbox}

\begin{solutionbox}
    Die \textbf{Binomialverteilung} modelliert die Anzahl der Erfolge in \(N\)
    unabhängigen Bernoulli-Experimenten.

    Für \(k\) Erfolge gilt:
    \[
        P(X=k)
        =
        \binom{N}{k}
        p^k(1-p)^{N-k}.
    \]
\end{solutionbox}

\begin{solutionbox}
    Die \textbf{Normalverteilung} oder Gauß-Verteilung ist eine der wichtigsten
    kontinuierlichen Verteilungen.

    Für Mittelwert \(\mu\) und Varianz \(\sigma^2\) gilt:
    \[
        \mathcal{N}(x \mid \mu,\sigma^2)
        =
        \frac{1}{\sqrt{2\pi\sigma^2}}
        \exp\left(
            -\frac{(x-\mu)^2}{2\sigma^2}
        \right).
    \]

    Sie ist wichtig für Messfehler, natürliche Streuung und viele approximative
    Methoden.
\end{solutionbox}

\begin{solutionbox}
    Die \textbf{multivariate Normalverteilung} verallgemeinert die
    Normalverteilung auf Vektoren.

    Für \(x \in \mathbb{R}^D\), Mittelwertvektor \(\mu\) und Kovarianzmatrix
    \(\Sigma\) gilt:
    \[
        \mathcal{N}(x \mid \mu,\Sigma)
        =
        \frac{1}{(2\pi)^{D/2}|\Sigma|^{1/2}}
        \exp\left(
            -\frac{1}{2}
            (x-\mu)^T
            \Sigma^{-1}
            (x-\mu)
        \right).
    \]

    Sie modelliert kontinuierliche, mehrdimensionale Daten.
\end{solutionbox}

\begin{answerbox}
    Wichtige Wahrscheinlichkeitsverteilungen sind:

    \begin{itemize}
        \item \textbf{Bernoulli-Verteilung:} binäre Ereignisse,
        \item \textbf{kategoriale Verteilung:} mehrere Klassen,
        \item \textbf{Binomialverteilung:} Anzahl von Erfolgen,
        \item \textbf{Normalverteilung:} kontinuierliche Streuung um einen Mittelwert,
        \item \textbf{multivariate Normalverteilung:} kontinuierliche Vektordaten.
    \end{itemize}

    Diese Verteilungen sind Grundbausteine probabilistischer Modelle.
\end{answerbox}

%%%%%%%%%%%%%%%%%%%%%%%%%%%%%%%%%%%%%%%%%%%%%%%%%%%%%%%%%%%%%%%%%%%
\subsection{Prüfungsrelevante Kurzfassung}
%%%%%%%%%%%%%%%%%%%%%%%%%%%%%%%%%%%%%%%%%%%%%%%%%%%%%%%%%%%%%%%%%%%

\begin{answerbox}
    Lecture 2 behandelt den probabilistischen Kalkül, also die mathematischen
    Rechenregeln für Wahrscheinlichkeiten.

    Die wichtigsten Punkte sind:

    \begin{itemize}
        \item Produktregel:
        \[
            P(X,Y)
            =
            P(X \mid Y)P(Y).
        \]

        \item Summenregel:
        \[
            P(X)
            =
            \sum_Y P(X,Y).
        \]

        \item Bayes' Regel:
        \[
            P(X \mid Y)
            =
            \frac{P(Y \mid X)P(X)}{P(Y)}.
        \]

        \item Für kontinuierliche Variablen ersetzt man Summen durch Integrale:
        \[
            p(x)
            =
            \int p(x,y)\,dy.
        \]

        \item Wahrscheinlichkeiten für kontinuierliche Variablen werden über
        Dichten berechnet:
        \[
            P(a \leq X \leq b)
            =
            \int_a^b p(x)\,dx.
        \]

        \item Erwartungswert:
        \[
            \mathbb{E}[f(X)]
            =
            \int f(x)p(x)\,dx.
        \]

        \item Sampling-Approximation:
        \[
            \mathbb{E}[f(X)]
            \approx
            \frac{1}{N}
            \sum_{n=1}^{N} f(x^{(n)}).
        \]

        \item Varianz:
        \[
            \operatorname{Var}(X)
            =
            \mathbb{E}[X^2] - \mathbb{E}[X]^2.
        \]

        \item Kovarianz:
        \[
            \operatorname{Cov}(X,Y)
            =
            \mathbb{E}[XY]
            -
            \mathbb{E}[X]\mathbb{E}[Y].
        \]

        \item Korrelation:
        \[
            \rho_{X,Y}
            =
            \frac{\operatorname{Cov}(X,Y)}
            {\sigma_X\sigma_Y}.
        \]
    \end{itemize}

    Die zentrale Botschaft lautet:
    \[
        \boxed{
            \text{Probabilistic calculus ist das Rechenwerkzeug für Machine Learning unter Unsicherheit.}
        }
    \]
\end{answerbox}

%%%%%%%%%%%%%%%%%%%%%%%%%%%%%%%%%%%%%%%%%%%%%%%%%%%%%%%%%%%%%%%%%%%
\section{Lecture 3: Probabilities and Estimators}
%%%%%%%%%%%%%%%%%%%%%%%%%%%%%%%%%%%%%%%%%%%%%%%%%%%%%%%%%%%%%%%%%%%

\subsection{Overview}

\begin{notebox}
    Lecture 3 behandelt den Übergang von Wahrscheinlichkeitsrechnung zu
    statistischer Schätzung.

    Die zentrale Frage lautet:

    \[
        \text{Wie kann man unbekannte Parameter einer Verteilung aus Daten schätzen?}
    \]

    Wichtige Themen sind:

    \begin{itemize}
        \item Parameterschätzung für Gleichverteilungen,
        \item Methode der Momente,
        \item Funktionen von Zufallsvariablen,
        \item i.i.d. Zufallsvariablen,
        \item statistische Schätzer,
        \item Bias und Varianz eines Schätzers,
        \item Bias und Varianz des Stichprobenmittelwerts,
        \item unverzerrte Stichprobenvarianz,
        \item Maximum-Likelihood-Schätzer,
        \item Maximum-a-posteriori-Schätzer,
        \item Bayes-Schätzer,
        \item Standardableitungen für ML- und MAP-Schätzer.
    \end{itemize}
\end{notebox}

\begin{answerbox}
    Die Kernaussage von Lecture 3 ist:

    \[
        \boxed{
            \text{Ein Schätzer ist eine Funktion der Daten, die einen unbekannten Parameter approximiert.}
        }
    \]

    In Machine Learning bedeutet Lernen häufig genau das:
    Aus Trainingsdaten werden Modellparameter geschätzt.
\end{answerbox}

%%%%%%%%%%%%%%%%%%%%%%%%%%%%%%%%%%%%%%%%%%%%%%%%%%%%%%%%%%%%%%%%%%%
\subsection{Parameter Estimation for Uniform Distributions}
%%%%%%%%%%%%%%%%%%%%%%%%%%%%%%%%%%%%%%%%%%%%%%%%%%%%%%%%%%%%%%%%%%%

\begin{notebox}
    Als einfaches Beispiel betrachten wir eine Gleichverteilung.

    Sei
    \[
        X \sim \operatorname{Unif}(a,b).
    \]

    Dann ist \(X\) gleichverteilt auf dem Intervall \([a,b]\).

    Die Dichte lautet:
    \[
        p(x \mid a,b)
        =
        \begin{cases}
            \dfrac{1}{b-a}, & a \leq x \leq b,\\[1em]
            0, & \text{sonst}.
        \end{cases}
    \]

    Die unbekannten Parameter sind:
    \[
        a
        \qquad \text{und} \qquad
        b.
    \]
\end{notebox}

\begin{solutionbox}
    Angenommen, wir beobachten Daten
    \[
        x^{(1)},x^{(2)},\dots,x^{(N)}
    \]
    und nehmen an, dass diese Daten aus einer Gleichverteilung stammen.

    Dann ist die Aufgabe:

    \[
        \text{Schätze die Intervallgrenzen } a \text{ und } b.
    \]

    Dafür gibt es verschiedene Ansätze. Zwei einfache Methoden sind:

    \begin{itemize}
        \item heuristische Schätzung mit Minimum und Maximum der Daten,
        \item Schätzung mit der Methode der Momente.
    \end{itemize}
\end{solutionbox}

%%%%%%%%%%%%%%%%%%%%%%%%%%%%%%%%%%%%%%%%%%%%%%%%%%%%%%%%%%%%%%%%%%%
\subsection{Heuristic Estimation Using Minimum and Maximum}
%%%%%%%%%%%%%%%%%%%%%%%%%%%%%%%%%%%%%%%%%%%%%%%%%%%%%%%%%%%%%%%%%%%

\begin{notebox}
    Eine naheliegende Idee ist:

    \[
        \hat a = \min_{1 \leq n \leq N} x^{(n)}
    \]
    und
    \[
        \hat b = \max_{1 \leq n \leq N} x^{(n)}.
    \]

    Man verwendet also den kleinsten beobachteten Datenpunkt als Schätzer für
    die untere Grenze und den größten beobachteten Datenpunkt als Schätzer für
    die obere Grenze.
\end{notebox}

\begin{calculationbox}
    Für Daten
    \[
        x^{(1)},\dots,x^{(N)}
    \]
    definieren wir:

    \[
        x_{\min}
        =
        \min_{1 \leq n \leq N} x^{(n)}
    \]

    und

    \[
        x_{\max}
        =
        \max_{1 \leq n \leq N} x^{(n)}.
    \]

    Der heuristische Schätzer lautet dann:

    \[
        \boxed{
            \hat a_{\mathrm{minmax}}
            =
            x_{\min}
        }
    \]

    und

    \[
        \boxed{
            \hat b_{\mathrm{minmax}}
            =
            x_{\max}.
        }
    \]
\end{calculationbox}

\begin{solutionbox}
    Diese Methode ist intuitiv, weil bei einer Gleichverteilung auf \([a,b]\)
    alle Datenpunkte innerhalb dieses Intervalls liegen müssen.

    Der kleinste beobachtete Wert ist daher ein Hinweis auf die untere Grenze,
    und der größte beobachtete Wert ist ein Hinweis auf die obere Grenze.

    Allerdings gilt typischerweise:

    \[
        x_{\min} > a
    \]
    und
    \[
        x_{\max} < b.
    \]

    Das bedeutet: Für endliche Stichproben unterschätzt diese Methode oft die
    Intervalllänge.
\end{solutionbox}

\begin{answerbox}
    Der Min-Max-Schätzer für eine Gleichverteilung ist:

    \[
        \boxed{
            \hat a = x_{\min}
        }
    \]

    \[
        \boxed{
            \hat b = x_{\max}
        }
    \]

    Er ist einfach und plausibel, verwendet aber nur die beiden extremen
    Datenpunkte.
\end{answerbox}

%%%%%%%%%%%%%%%%%%%%%%%%%%%%%%%%%%%%%%%%%%%%%%%%%%%%%%%%%%%%%%%%%%%
\subsection{Estimation Using the Method of Moments}
%%%%%%%%%%%%%%%%%%%%%%%%%%%%%%%%%%%%%%%%%%%%%%%%%%%%%%%%%%%%%%%%%%%

\begin{notebox}
    Die Methode der Momente schätzt Parameter, indem theoretische Momente einer
    Verteilung mit empirischen Momenten der Daten gleichgesetzt werden.

    Für eine Gleichverteilung \(X \sim \operatorname{Unif}(a,b)\) gilt:

    \[
        \mathbb{E}[X]
        =
        \frac{a+b}{2}
    \]

    und

    \[
        \operatorname{Var}(X)
        =
        \frac{(b-a)^2}{12}.
    \]
\end{notebox}

\begin{calculationbox}
    Zuerst berechnet man aus den Daten den empirischen Mittelwert:

    \[
        \bar x
        =
        \frac{1}{N}
        \sum_{n=1}^{N} x^{(n)}.
    \]

    Dann berechnet man die empirische Varianz mit Faktor \(1/N\):

    \[
        v
        =
        \frac{1}{N}
        \sum_{n=1}^{N}
        \left(x^{(n)}-\bar x\right)^2.
    \]

    Die Methode der Momente setzt nun:

    \[
        \bar x
        =
        \frac{a+b}{2}
    \]

    und

    \[
        v
        =
        \frac{(b-a)^2}{12}.
    \]
\end{calculationbox}

\begin{calculationbox}
    Aus
    \[
        \bar x
        =
        \frac{a+b}{2}
    \]
    folgt:
    \[
        a+b
        =
        2\bar x.
    \]

    Aus
    \[
        v
        =
        \frac{(b-a)^2}{12}
    \]
    folgt:
    \[
        (b-a)^2
        =
        12v.
    \]

    Also:
    \[
        b-a
        =
        \sqrt{12v}.
    \]

    Setze
    \[
        m = \bar x
    \]
    und
    \[
        d = \frac{b-a}{2}.
    \]

    Dann gilt:
    \[
        d
        =
        \frac{\sqrt{12v}}{2}
        =
        \sqrt{3v}.
    \]

    Damit erhält man:

    \[
        \hat a_{\mathrm{MoM}}
        =
        \bar x - \sqrt{3v}
    \]

    und

    \[
        \hat b_{\mathrm{MoM}}
        =
        \bar x + \sqrt{3v}.
    \]
\end{calculationbox}

\begin{answerbox}
    Für \(X \sim \operatorname{Unif}(a,b)\) liefert die Methode der Momente:

    \[
        \boxed{
            \hat a_{\mathrm{MoM}}
            =
            \bar x - \sqrt{3v}
        }
    \]

    \[
        \boxed{
            \hat b_{\mathrm{MoM}}
            =
            \bar x + \sqrt{3v}
        }
    \]

    wobei

    \[
        \boxed{
            \bar x
            =
            \frac{1}{N}
            \sum_{n=1}^{N} x^{(n)}
        }
    \]

    und

    \[
        \boxed{
            v
            =
            \frac{1}{N}
            \sum_{n=1}^{N}
            \left(x^{(n)}-\bar x\right)^2
        }.
    \]
\end{answerbox}

%%%%%%%%%%%%%%%%%%%%%%%%%%%%%%%%%%%%%%%%%%%%%%%%%%%%%%%%%%%%%%%%%%%
\subsection{Which Estimation Method Is Better?}
%%%%%%%%%%%%%%%%%%%%%%%%%%%%%%%%%%%%%%%%%%%%%%%%%%%%%%%%%%%%%%%%%%%

\begin{notebox}
    Für die Gleichverteilung haben wir zwei Schätzmethoden betrachtet:

    \begin{itemize}
        \item Min-Max-Schätzung,
        \item Methode der Momente.
    \end{itemize}

    Die Frage ist:

    \[
        \text{Welche Methode ist besser?}
    \]
\end{notebox}

\begin{solutionbox}
    Die Antwort hängt davon ab, welches Kriterium man betrachtet.

    Die Min-Max-Schätzung verwendet nur zwei Datenpunkte:
    \[
        x_{\min}
        \qquad \text{und} \qquad
        x_{\max}.
    \]

    Dadurch ist sie sehr direkt für die Intervallgrenzen, aber auch stark von
    Extremwerten abhängig.

    Die Methode der Momente verwendet dagegen alle Datenpunkte über Mittelwert
    und Varianz:
    \[
        \bar x
        \qquad \text{und} \qquad
        v.
    \]

    Dadurch nutzt sie mehr Information aus den Daten, kann aber Intervallgrenzen
    liefern, die nicht exakt zu Minimum und Maximum der beobachteten Daten
    passen.
\end{solutionbox}

\begin{answerbox}
    Es gibt nicht immer eine universell beste Methode.

    \begin{itemize}
        \item Der Min-Max-Schätzer ist natürlich für Intervallgrenzen und eng mit
        der Maximum-Likelihood-Idee für die Gleichverteilung verbunden.

        \item Die Methode der Momente ist systematisch und verwendet Mittelwert
        und Varianz, also Informationen aus allen Datenpunkten.

        \item Für endliche Stichproben kann man Methoden anhand von Bias, Varianz
        und mittlerem quadratischem Fehler vergleichen.
    \end{itemize}

    Wichtige Kriterien sind daher:

    \[
        \boxed{
            \text{Bias, Varianz und Mean Squared Error.}
        }
    \]
\end{answerbox}

%%%%%%%%%%%%%%%%%%%%%%%%%%%%%%%%%%%%%%%%%%%%%%%%%%%%%%%%%%%%%%%%%%%
\subsection{Mean and Variance of Functions of Random Numbers}
%%%%%%%%%%%%%%%%%%%%%%%%%%%%%%%%%%%%%%%%%%%%%%%%%%%%%%%%%%%%%%%%%%%

\begin{notebox}
    Oft interessiert man sich nicht nur für eine Zufallsvariable \(X\), sondern
    für eine Funktion dieser Zufallsvariable:

    \[
        Y = f(X).
    \]

    Dann sind auch Mittelwert und Varianz von \(Y\) relevant.
\end{notebox}

\begin{calculationbox}
    Für eine diskrete Zufallsvariable gilt:

    \[
        \mathbb{E}[f(X)]
        =
        \sum_x f(x)P(X=x).
    \]

    Für eine kontinuierliche Zufallsvariable mit Dichte \(p(x)\) gilt:

    \[
        \mathbb{E}[f(X)]
        =
        \int f(x)p(x)\,dx.
    \]

    Die Varianz von \(f(X)\) ist:

    \[
        \operatorname{Var}(f(X))
        =
        \mathbb{E}\left[
            \left(f(X)-\mathbb{E}[f(X)]\right)^2
        \right].
    \]
\end{calculationbox}

\begin{calculationbox}
    Äquivalent kann man die Varianz schreiben als:

    \[
        \operatorname{Var}(f(X))
        =
        \mathbb{E}[f(X)^2]
        -
        \mathbb{E}[f(X)]^2.
    \]

    Diese Darstellung ist oft rechnerisch einfacher.
\end{calculationbox}

\begin{answerbox}
    Für eine Funktion \(Y=f(X)\) gilt:

    \[
        \boxed{
            \mathbb{E}[Y]
            =
            \mathbb{E}[f(X)]
        }
    \]

    und

    \[
        \boxed{
            \operatorname{Var}(Y)
            =
            \mathbb{E}[f(X)^2]
            -
            \mathbb{E}[f(X)]^2.
        }
    \]
\end{answerbox}

%%%%%%%%%%%%%%%%%%%%%%%%%%%%%%%%%%%%%%%%%%%%%%%%%%%%%%%%%%%%%%%%%%%
\subsection{i.i.d. Random Variables}
%%%%%%%%%%%%%%%%%%%%%%%%%%%%%%%%%%%%%%%%%%%%%%%%%%%%%%%%%%%%%%%%%%%

\begin{notebox}
    In der Statistik und im Machine Learning nimmt man häufig an, dass Daten
    \emph{i.i.d.} sind.

    Die Abkürzung steht für:

    \[
        \text{independent and identically distributed}.
    \]

    Das bedeutet:

    \begin{itemize}
        \item Die Zufallsvariablen sind unabhängig.
        \item Alle Zufallsvariablen haben dieselbe Verteilung.
    \end{itemize}
\end{notebox}

\begin{solutionbox}
    Für Daten
    \[
        X_1,\dots,X_N
    \]
    bedeutet i.i.d.:

    \[
        X_1,\dots,X_N
        \sim p(x)
    \]

    und die gemeinsame Verteilung faktorisiert:

    \[
        p(x_1,\dots,x_N)
        =
        \prod_{n=1}^{N} p(x_n).
    \]

    Diese Faktorisierung ist extrem wichtig für Maximum-Likelihood-Schätzung.
\end{solutionbox}

\begin{answerbox}
    Die i.i.d.-Annahme lautet:

    \[
        \boxed{
            X_1,\dots,X_N
            \text{ sind unabhängig und identisch verteilt.}
        }
    \]

    Dann gilt:

    \[
        \boxed{
            p(x_1,\dots,x_N)
            =
            \prod_{n=1}^{N}p(x_n).
        }
    \]
\end{answerbox}

%%%%%%%%%%%%%%%%%%%%%%%%%%%%%%%%%%%%%%%%%%%%%%%%%%%%%%%%%%%%%%%%%%%
\subsection{Statistical Estimators}
%%%%%%%%%%%%%%%%%%%%%%%%%%%%%%%%%%%%%%%%%%%%%%%%%%%%%%%%%%%%%%%%%%%

\begin{notebox}
    Ein statistischer Schätzer ist eine Funktion der Daten, die einen unbekannten
    Parameter einer Verteilung approximieren soll.

    Sei \(\theta\) ein unbekannter Parameter und seien

    \[
        X_1,\dots,X_N
    \]

    Zufallsvariablen, aus denen die Daten entstehen.

    Ein Schätzer ist dann eine Funktion

    \[
        \hat\theta
        =
        \hat\theta(X_1,\dots,X_N).
    \]
\end{notebox}

\begin{solutionbox}
    Wichtig ist:

    Da der Schätzer von Zufallsvariablen abhängt, ist der Schätzer selbst wieder
    eine Zufallsvariable.

    Für konkrete beobachtete Daten

    \[
        x_1,\dots,x_N
    \]

    erhält man einen konkreten Schätzwert:

    \[
        \hat\theta(x_1,\dots,x_N).
    \]

    Vor der Beobachtung der Daten ist \(\hat\theta\) aber zufällig.
\end{solutionbox}

\begin{answerbox}
    Ein Schätzer ist:

    \[
        \boxed{
            \hat\theta
            =
            \hat\theta(X_1,\dots,X_N)
        }
    \]

    Er ist eine Funktion der Daten und damit selbst eine Zufallsvariable.

    Ein guter Schätzer sollte den unbekannten Parameter \(\theta\) möglichst gut
    approximieren.
\end{answerbox}

%%%%%%%%%%%%%%%%%%%%%%%%%%%%%%%%%%%%%%%%%%%%%%%%%%%%%%%%%%%%%%%%%%%
\subsection{Bias and Variance of Estimators}
%%%%%%%%%%%%%%%%%%%%%%%%%%%%%%%%%%%%%%%%%%%%%%%%%%%%%%%%%%%%%%%%%%%

\begin{notebox}
    Zwei zentrale Eigenschaften eines Schätzers sind:

    \begin{itemize}
        \item Bias,
        \item Varianz.
    \end{itemize}

    Der Bias misst, ob ein Schätzer systematisch danebenliegt.

    Die Varianz misst, wie stark der Schätzer zwischen verschiedenen Stichproben
    schwankt.
\end{notebox}

\begin{calculationbox}
    Der Bias eines Schätzers \(\hat\theta\) für einen Parameter \(\theta\) ist:

    \[
        \operatorname{Bias}(\hat\theta)
        =
        \mathbb{E}[\hat\theta] - \theta.
    \]

    Ist

    \[
        \mathbb{E}[\hat\theta] = \theta,
    \]

    dann ist der Schätzer unverzerrt oder unbiased.

    Also:

    \[
        \boxed{
            \operatorname{Bias}(\hat\theta)=0
            \quad \Longleftrightarrow \quad
            \hat\theta \text{ ist unverzerrt.}
        }
    \]
\end{calculationbox}

\begin{calculationbox}
    Die Varianz eines Schätzers ist:

    \[
        \operatorname{Var}(\hat\theta)
        =
        \mathbb{E}\left[
            \left(\hat\theta-\mathbb{E}[\hat\theta]\right)^2
        \right].
    \]

    Sie beschreibt, wie stark die Schätzung über verschiedene mögliche
    Datensätze schwankt.
\end{calculationbox}

\begin{calculationbox}
    Ein häufiges Fehlermaß ist der mittlere quadratische Fehler:

    \[
        \operatorname{MSE}(\hat\theta)
        =
        \mathbb{E}\left[
            \left(\hat\theta-\theta\right)^2
        \right].
    \]

    Dieser zerfällt in Bias und Varianz:

    \[
        \operatorname{MSE}(\hat\theta)
        =
        \operatorname{Var}(\hat\theta)
        +
        \operatorname{Bias}(\hat\theta)^2.
    \]
\end{calculationbox}

\begin{answerbox}
    Für einen Schätzer \(\hat\theta\) gilt:

    \[
        \boxed{
            \operatorname{Bias}(\hat\theta)
            =
            \mathbb{E}[\hat\theta]-\theta
        }
    \]

    \[
        \boxed{
            \operatorname{Var}(\hat\theta)
            =
            \mathbb{E}
            \left[
                \left(\hat\theta-\mathbb{E}[\hat\theta]\right)^2
            \right]
        }
    \]

    \[
        \boxed{
            \operatorname{MSE}(\hat\theta)
            =
            \operatorname{Var}(\hat\theta)
            +
            \operatorname{Bias}(\hat\theta)^2
        }
    \]
\end{answerbox}

%%%%%%%%%%%%%%%%%%%%%%%%%%%%%%%%%%%%%%%%%%%%%%%%%%%%%%%%%%%%%%%%%%%
\subsection{Bias and Variance of the Sample Mean Estimator}
%%%%%%%%%%%%%%%%%%%%%%%%%%%%%%%%%%%%%%%%%%%%%%%%%%%%%%%%%%%%%%%%%%%

\begin{notebox}
    Sei

    \[
        X_1,\dots,X_N
    \]

    eine i.i.d. Stichprobe mit

    \[
        \mathbb{E}[X_i]=\mu
    \]

    und

    \[
        \operatorname{Var}(X_i)=\sigma^2.
    \]

    Der Stichprobenmittelwert ist:

    \[
        \bar X
        =
        \frac{1}{N}
        \sum_{n=1}^{N} X_n.
    \]
\end{notebox}

\begin{calculationbox}
    Erwartungswert des Stichprobenmittelwerts:

    \[
        \mathbb{E}[\bar X]
        =
        \mathbb{E}
        \left[
            \frac{1}{N}
            \sum_{n=1}^{N} X_n
        \right].
    \]

    Wegen Linearität des Erwartungswerts:

    \[
        \mathbb{E}[\bar X]
        =
        \frac{1}{N}
        \sum_{n=1}^{N}
        \mathbb{E}[X_n].
    \]

    Da alle \(X_n\) denselben Erwartungswert \(\mu\) haben:

    \[
        \mathbb{E}[\bar X]
        =
        \frac{1}{N}
        \sum_{n=1}^{N}
        \mu
        =
        \frac{1}{N}N\mu
        =
        \mu.
    \]

    Also:

    \[
        \operatorname{Bias}(\bar X)
        =
        \mathbb{E}[\bar X]-\mu
        =
        0.
    \]
\end{calculationbox}

\begin{calculationbox}
    Varianz des Stichprobenmittelwerts:

    \[
        \operatorname{Var}(\bar X)
        =
        \operatorname{Var}
        \left(
            \frac{1}{N}
            \sum_{n=1}^{N}X_n
        \right).
    \]

    Konstanten können quadratisch herausgezogen werden:

    \[
        \operatorname{Var}(\bar X)
        =
        \frac{1}{N^2}
        \operatorname{Var}
        \left(
            \sum_{n=1}^{N}X_n
        \right).
    \]

    Wegen Unabhängigkeit gilt:

    \[
        \operatorname{Var}
        \left(
            \sum_{n=1}^{N}X_n
        \right)
        =
        \sum_{n=1}^{N}\operatorname{Var}(X_n).
    \]

    Da alle \(X_n\) Varianz \(\sigma^2\) haben:

    \[
        \operatorname{Var}(\bar X)
        =
        \frac{1}{N^2}
        \sum_{n=1}^{N}
        \sigma^2
        =
        \frac{1}{N^2}N\sigma^2
        =
        \frac{\sigma^2}{N}.
    \]
\end{calculationbox}

\begin{answerbox}
    Für den Stichprobenmittelwert gilt:

    \[
        \boxed{
            \mathbb{E}[\bar X]=\mu
        }
    \]

    \[
        \boxed{
            \operatorname{Bias}(\bar X)=0
        }
    \]

    \[
        \boxed{
            \operatorname{Var}(\bar X)
            =
            \frac{\sigma^2}{N}
        }
    \]

    Der Stichprobenmittelwert ist also unverzerrt, und seine Varianz sinkt mit
    wachsender Stichprobengröße \(N\).
\end{answerbox}

%%%%%%%%%%%%%%%%%%%%%%%%%%%%%%%%%%%%%%%%%%%%%%%%%%%%%%%%%%%%%%%%%%%
\subsection{Bias of the Naive Sample Variance Estimator}
%%%%%%%%%%%%%%%%%%%%%%%%%%%%%%%%%%%%%%%%%%%%%%%%%%%%%%%%%%%%%%%%%%%

\begin{notebox}
    Eine naive Schätzung der Varianz ist:

    \[
        \hat\sigma_N^2
        =
        \frac{1}{N}
        \sum_{n=1}^{N}
        (X_n-\bar X)^2.
    \]

    Diese Formel sieht natürlich aus, ist aber für die wahre Varianz
    \(\sigma^2\) verzerrt.
\end{notebox}

\begin{calculationbox}
    Man kann zeigen:

    \[
        \mathbb{E}
        \left[
            \frac{1}{N}
            \sum_{n=1}^{N}
            (X_n-\bar X)^2
        \right]
        =
        \frac{N-1}{N}
        \sigma^2.
    \]

    Daher ist der Bias:

    \[
        \operatorname{Bias}(\hat\sigma_N^2)
        =
        \mathbb{E}[\hat\sigma_N^2]-\sigma^2
    \]

    also:

    \[
        \operatorname{Bias}(\hat\sigma_N^2)
        =
        \frac{N-1}{N}\sigma^2-\sigma^2.
    \]

    Zusammenfassen ergibt:

    \[
        \operatorname{Bias}(\hat\sigma_N^2)
        =
        -\frac{1}{N}\sigma^2.
    \]
\end{calculationbox}

\begin{answerbox}
    Der naive Varianzschätzer

    \[
        \boxed{
            \hat\sigma_N^2
            =
            \frac{1}{N}
            \sum_{n=1}^{N}
            (X_n-\bar X)^2
        }
    \]

    ist verzerrt:

    \[
        \boxed{
            \mathbb{E}[\hat\sigma_N^2]
            =
            \frac{N-1}{N}\sigma^2
        }
    \]

    und damit:

    \[
        \boxed{
            \operatorname{Bias}(\hat\sigma_N^2)
            =
            -\frac{\sigma^2}{N}.
        }
    \]
\end{answerbox}

%%%%%%%%%%%%%%%%%%%%%%%%%%%%%%%%%%%%%%%%%%%%%%%%%%%%%%%%%%%%%%%%%%%
\subsection{An Unbiased Estimator for the Sample Variance}
%%%%%%%%%%%%%%%%%%%%%%%%%%%%%%%%%%%%%%%%%%%%%%%%%%%%%%%%%%%%%%%%%%%

\begin{notebox}
    Um einen unverzerrten Varianzschätzer zu erhalten, ersetzt man den Faktor
    \(1/N\) durch \(1/(N-1)\).

    Der korrigierte Schätzer lautet:

    \[
        S^2
        =
        \frac{1}{N-1}
        \sum_{n=1}^{N}
        (X_n-\bar X)^2.
    \]
\end{notebox}

\begin{calculationbox}
    Aus der vorherigen Beziehung wissen wir:

    \[
        \mathbb{E}
        \left[
            \frac{1}{N}
            \sum_{n=1}^{N}
            (X_n-\bar X)^2
        \right]
        =
        \frac{N-1}{N}
        \sigma^2.
    \]

    Multiplizieren wir den naiven Schätzer mit

    \[
        \frac{N}{N-1},
    \]

    erhalten wir:

    \[
        \frac{N}{N-1}
        \cdot
        \frac{1}{N}
        \sum_{n=1}^{N}
        (X_n-\bar X)^2
        =
        \frac{1}{N-1}
        \sum_{n=1}^{N}
        (X_n-\bar X)^2.
    \]

    Daher gilt:

    \[
        \mathbb{E}[S^2]
        =
        \sigma^2.
    \]
\end{calculationbox}

\begin{answerbox}
    Der unverzerrte Schätzer für die Varianz ist:

    \[
        \boxed{
            S^2
            =
            \frac{1}{N-1}
            \sum_{n=1}^{N}
            (X_n-\bar X)^2
        }
    \]

    Für ihn gilt:

    \[
        \boxed{
            \mathbb{E}[S^2]
            =
            \sigma^2.
        }
    \]

    Die Korrektur von \(N\) zu \(N-1\) heißt Bessel-Korrektur.
\end{answerbox}

%%%%%%%%%%%%%%%%%%%%%%%%%%%%%%%%%%%%%%%%%%%%%%%%%%%%%%%%%%%%%%%%%%%
\subsection{Maximum Likelihood Estimators}
%%%%%%%%%%%%%%%%%%%%%%%%%%%%%%%%%%%%%%%%%%%%%%%%%%%%%%%%%%%%%%%%%%%

\begin{notebox}
    Maximum Likelihood Estimation, kurz MLE, ist eine der wichtigsten Methoden
    zur Parameterschätzung.

    Idee:

    \[
        \text{Wähle den Parameter, unter dem die beobachteten Daten am wahrscheinlichsten sind.}
    \]

    Sei \(\theta\) ein Parameter und seien die Daten

    \[
        x_1,\dots,x_N.
    \]
\end{notebox}

\begin{calculationbox}
    Die Likelihood ist die Wahrscheinlichkeit der beobachteten Daten als
    Funktion des Parameters:

    \[
        L(\theta)
        =
        p(x_1,\dots,x_N \mid \theta).
    \]

    Unter der i.i.d.-Annahme gilt:

    \[
        L(\theta)
        =
        \prod_{n=1}^{N}
        p(x_n \mid \theta).
    \]

    Der Maximum-Likelihood-Schätzer ist:

    \[
        \hat\theta_{\mathrm{ML}}
        =
        \operatorname*{arg\,max}_{\theta}
        L(\theta).
    \]
\end{calculationbox}

\begin{calculationbox}
    Da Produkte oft schwer zu optimieren sind, verwendet man meistens die
    Log-Likelihood:

    \[
        \ell(\theta)
        =
        \log L(\theta).
    \]

    Wegen der Monotonie des Logarithmus gilt:

    \[
        \operatorname*{arg\,max}_{\theta} L(\theta)
        =
        \operatorname*{arg\,max}_{\theta} \ell(\theta).
    \]

    Unter der i.i.d.-Annahme:

    \[
        \ell(\theta)
        =
        \log
        \prod_{n=1}^{N}
        p(x_n \mid \theta)
        =
        \sum_{n=1}^{N}
        \log p(x_n \mid \theta).
    \]
\end{calculationbox}

\begin{answerbox}
    Der Maximum-Likelihood-Schätzer ist:

    \[
        \boxed{
            \hat\theta_{\mathrm{ML}}
            =
            \operatorname*{arg\,max}_{\theta}
            p(x_1,\dots,x_N \mid \theta)
        }
    \]

    Für i.i.d. Daten:

    \[
        \boxed{
            \hat\theta_{\mathrm{ML}}
            =
            \operatorname*{arg\,max}_{\theta}
            \prod_{n=1}^{N}
            p(x_n \mid \theta)
        }
    \]

    Äquivalent maximiert man die Log-Likelihood:

    \[
        \boxed{
            \hat\theta_{\mathrm{ML}}
            =
            \operatorname*{arg\,max}_{\theta}
            \sum_{n=1}^{N}
            \log p(x_n \mid \theta).
        }
    \]
\end{answerbox}

%%%%%%%%%%%%%%%%%%%%%%%%%%%%%%%%%%%%%%%%%%%%%%%%%%%%%%%%%%%%%%%%%%%
\subsection{Maximum A-Posteriori Estimators}
%%%%%%%%%%%%%%%%%%%%%%%%%%%%%%%%%%%%%%%%%%%%%%%%%%%%%%%%%%%%%%%%%%%

\begin{notebox}
    Maximum a-posteriori Estimation, kurz MAP, erweitert Maximum Likelihood um
    Vorwissen über den Parameter.

    Dieses Vorwissen wird als Prior modelliert:

    \[
        p(\theta).
    \]

    Gesucht ist nun der wahrscheinlichste Parameter nach Beobachtung der Daten:

    \[
        p(\theta \mid x_1,\dots,x_N).
    \]
\end{notebox}

\begin{calculationbox}
    Nach Bayes' Regel gilt:

    \[
        p(\theta \mid x_1,\dots,x_N)
        =
        \frac{
            p(x_1,\dots,x_N \mid \theta)p(\theta)
        }{
            p(x_1,\dots,x_N)
        }.
    \]

    Der Nenner hängt nicht von \(\theta\) ab.

    Daher ist:

    \[
        \hat\theta_{\mathrm{MAP}}
        =
        \operatorname*{arg\,max}_{\theta}
        p(\theta \mid x_1,\dots,x_N)
    \]

    äquivalent zu:

    \[
        \hat\theta_{\mathrm{MAP}}
        =
        \operatorname*{arg\,max}_{\theta}
        p(x_1,\dots,x_N \mid \theta)p(\theta).
    \]
\end{calculationbox}

\begin{calculationbox}
    Für i.i.d. Daten gilt:

    \[
        p(x_1,\dots,x_N \mid \theta)
        =
        \prod_{n=1}^{N} p(x_n \mid \theta).
    \]

    Also:

    \[
        \hat\theta_{\mathrm{MAP}}
        =
        \operatorname*{arg\,max}_{\theta}
        \left[
            \prod_{n=1}^{N}
            p(x_n \mid \theta)
        \right]
        p(\theta).
    \]

    In Log-Form:

    \[
        \hat\theta_{\mathrm{MAP}}
        =
        \operatorname*{arg\,max}_{\theta}
        \left[
            \sum_{n=1}^{N}
            \log p(x_n \mid \theta)
            +
            \log p(\theta)
        \right].
    \]
\end{calculationbox}

\begin{answerbox}
    Der MAP-Schätzer ist:

    \[
        \boxed{
            \hat\theta_{\mathrm{MAP}}
            =
            \operatorname*{arg\,max}_{\theta}
            p(\theta \mid x_1,\dots,x_N)
        }
    \]

    Äquivalent:

    \[
        \boxed{
            \hat\theta_{\mathrm{MAP}}
            =
            \operatorname*{arg\,max}_{\theta}
            p(x_1,\dots,x_N \mid \theta)p(\theta)
        }
    \]

    In Log-Form:

    \[
        \boxed{
            \hat\theta_{\mathrm{MAP}}
            =
            \operatorname*{arg\,max}_{\theta}
            \left[
                \sum_{n=1}^{N}
                \log p(x_n \mid \theta)
                +
                \log p(\theta)
            \right].
        }
    \]
\end{answerbox}

%%%%%%%%%%%%%%%%%%%%%%%%%%%%%%%%%%%%%%%%%%%%%%%%%%%%%%%%%%%%%%%%%%%
\subsection{Bayesian Estimators}
%%%%%%%%%%%%%%%%%%%%%%%%%%%%%%%%%%%%%%%%%%%%%%%%%%%%%%%%%%%%%%%%%%%

\begin{notebox}
    Bayessche Schätzung geht noch einen Schritt weiter als MAP.

    Bei ML und MAP wird am Ende ein einzelner Parameterwert geschätzt:

    \[
        \hat\theta_{\mathrm{ML}}
        \qquad \text{oder} \qquad
        \hat\theta_{\mathrm{MAP}}.
    \]

    Bei Bayesscher Schätzung betrachtet man stattdessen die gesamte
    Posterior-Verteilung:

    \[
        p(\theta \mid x_1,\dots,x_N).
    \]
\end{notebox}

\begin{solutionbox}
    Die Posterior-Verteilung beschreibt die Unsicherheit über den Parameter nach
    Beobachtung der Daten.

    Bayes' Regel liefert:

    \[
        p(\theta \mid x_1,\dots,x_N)
        =
        \frac{
            p(x_1,\dots,x_N \mid \theta)p(\theta)
        }{
            p(x_1,\dots,x_N)
        }.
    \]

    Dabei ist:

    \[
        p(x_1,\dots,x_N)
        =
        \int
        p(x_1,\dots,x_N \mid \theta)p(\theta)
        \,d\theta.
    \]

    Diese Größe normalisiert den Posterior.
\end{solutionbox}

\begin{answerbox}
    Bayessche Schätzung liefert nicht nur einen Punktwert, sondern eine ganze
    Verteilung:

    \[
        \boxed{
            p(\theta \mid x_1,\dots,x_N)
        }
    \]

    MAP wählt nur den wahrscheinlichsten Wert dieser Verteilung:

    \[
        \boxed{
            \hat\theta_{\mathrm{MAP}}
            =
            \operatorname*{arg\,max}_{\theta}
            p(\theta \mid x_1,\dots,x_N).
        }
    \]

    Der Bayessche Ansatz behält dagegen die Unsicherheit über \(\theta\) bei.
\end{answerbox}

%%%%%%%%%%%%%%%%%%%%%%%%%%%%%%%%%%%%%%%%%%%%%%%%%%%%%%%%%%%%%%%%%%%
\subsection{Standard Derivation Technique for Maximum Likelihood Estimators}
%%%%%%%%%%%%%%%%%%%%%%%%%%%%%%%%%%%%%%%%%%%%%%%%%%%%%%%%%%%%%%%%%%%

\begin{notebox}
    Für Maximum-Likelihood-Schätzer gibt es ein Standardvorgehen.
\end{notebox}

\begin{calculationbox}
    Vorgehen zur Herleitung eines ML-Schätzers:

    \begin{enumerate}
        \item Modell bzw. Dichte aufschreiben:
        \[
            p(x \mid \theta).
        \]

        \item Likelihood für die Daten bilden:
        \[
            L(\theta)
            =
            \prod_{n=1}^{N}
            p(x_n \mid \theta).
        \]

        \item Log-Likelihood bilden:
        \[
            \ell(\theta)
            =
            \sum_{n=1}^{N}
            \log p(x_n \mid \theta).
        \]

        \item Nach dem Parameter ableiten:
        \[
            \frac{\partial}{\partial \theta}
            \ell(\theta).
        \]

        \item Ableitung gleich null setzen:
        \[
            \frac{\partial}{\partial \theta}
            \ell(\theta)
            =
            0.
        \]

        \item Nach \(\theta\) auflösen.
    \end{enumerate}
\end{calculationbox}

\begin{answerbox}
    Standardform für ML:

    \[
        \boxed{
            \hat\theta_{\mathrm{ML}}
            =
            \operatorname*{arg\,max}_{\theta}
            \sum_{n=1}^{N}
            \log p(x_n \mid \theta)
        }
    \]

    Praktisch rechnet man meistens:

    \[
        \boxed{
            \frac{\partial}{\partial \theta}
            \ell(\theta)
            =
            0
        }
    \]

    und löst nach \(\theta\) auf.
\end{answerbox}

%%%%%%%%%%%%%%%%%%%%%%%%%%%%%%%%%%%%%%%%%%%%%%%%%%%%%%%%%%%%%%%%%%%
\subsection{Standard Derivation Technique for MAP Estimators}
%%%%%%%%%%%%%%%%%%%%%%%%%%%%%%%%%%%%%%%%%%%%%%%%%%%%%%%%%%%%%%%%%%%

\begin{notebox}
    Für MAP-Schätzer ist die Herleitung fast gleich wie bei ML, aber zusätzlich
    kommt der Prior \(p(\theta)\) hinzu.
\end{notebox}

\begin{calculationbox}
    Vorgehen zur Herleitung eines MAP-Schätzers:

    \begin{enumerate}
        \item Modell bzw. Likelihood aufschreiben:
        \[
            p(x_1,\dots,x_N \mid \theta).
        \]

        \item Prior auf den Parameter aufschreiben:
        \[
            p(\theta).
        \]

        \item Posterior bis auf Normierung bilden:
        \[
            p(\theta \mid x_1,\dots,x_N)
            \propto
            p(x_1,\dots,x_N \mid \theta)p(\theta).
        \]

        \item Log-Posterior bilden:
        \[
            \log p(\theta \mid x_1,\dots,x_N)
            =
            \sum_{n=1}^{N}
            \log p(x_n \mid \theta)
            +
            \log p(\theta)
            +
            \text{const}.
        \]

        \item Nach \(\theta\) ableiten und gleich null setzen.
    \end{enumerate}
\end{calculationbox}

\begin{answerbox}
    Standardform für MAP:

    \[
        \boxed{
            \hat\theta_{\mathrm{MAP}}
            =
            \operatorname*{arg\,max}_{\theta}
            \left[
                \sum_{n=1}^{N}
                \log p(x_n \mid \theta)
                +
                \log p(\theta)
            \right]
        }
    \]

    Der Unterschied zu ML ist der zusätzliche Term:

    \[
        \boxed{
            \log p(\theta).
        }
    \]
\end{answerbox}

%%%%%%%%%%%%%%%%%%%%%%%%%%%%%%%%%%%%%%%%%%%%%%%%%%%%%%%%%%%%%%%%%%%
\subsection{Prüfungsrelevante Kurzfassung}
%%%%%%%%%%%%%%%%%%%%%%%%%%%%%%%%%%%%%%%%%%%%%%%%%%%%%%%%%%%%%%%%%%%

\begin{answerbox}
    Lecture 3 behandelt Schätzer für unbekannte Parameter von Verteilungen.

    Die wichtigsten Punkte sind:

    \begin{itemize}
        \item Für eine Gleichverteilung \(X \sim \operatorname{Unif}(a,b)\) gilt:
        \[
            \mathbb{E}[X]=\frac{a+b}{2},
            \qquad
            \operatorname{Var}(X)=\frac{(b-a)^2}{12}.
        \]

        \item Min-Max-Schätzer:
        \[
            \hat a = x_{\min},
            \qquad
            \hat b = x_{\max}.
        \]

        \item Methode der Momente:
        \[
            \hat a_{\mathrm{MoM}}
            =
            \bar x-\sqrt{3v},
            \qquad
            \hat b_{\mathrm{MoM}}
            =
            \bar x+\sqrt{3v}.
        \]

        \item i.i.d. bedeutet unabhängig und identisch verteilt:
        \[
            p(x_1,\dots,x_N)
            =
            \prod_{n=1}^{N}p(x_n).
        \]

        \item Ein Schätzer ist eine Funktion der Daten:
        \[
            \hat\theta
            =
            \hat\theta(X_1,\dots,X_N).
        \]

        \item Bias:
        \[
            \operatorname{Bias}(\hat\theta)
            =
            \mathbb{E}[\hat\theta]-\theta.
        \]

        \item Varianz des Schätzers:
        \[
            \operatorname{Var}(\hat\theta)
            =
            \mathbb{E}
            \left[
                \left(\hat\theta-\mathbb{E}[\hat\theta]\right)^2
            \right].
        \]

        \item Stichprobenmittelwert:
        \[
            \bar X
            =
            \frac{1}{N}
            \sum_{n=1}^{N}X_n.
        \]

        \item Für den Stichprobenmittelwert gilt:
        \[
            \mathbb{E}[\bar X]=\mu,
            \qquad
            \operatorname{Var}(\bar X)=\frac{\sigma^2}{N}.
        \]

        \item Unverzerrte Stichprobenvarianz:
        \[
            S^2
            =
            \frac{1}{N-1}
            \sum_{n=1}^{N}(X_n-\bar X)^2.
        \]

        \item Maximum Likelihood:
        \[
            \hat\theta_{\mathrm{ML}}
            =
            \operatorname*{arg\,max}_{\theta}
            \sum_{n=1}^{N}
            \log p(x_n \mid \theta).
        \]

        \item Maximum a-posteriori:
        \[
            \hat\theta_{\mathrm{MAP}}
            =
            \operatorname*{arg\,max}_{\theta}
            \left[
                \sum_{n=1}^{N}
                \log p(x_n \mid \theta)
                +
                \log p(\theta)
            \right].
        \]

        \item Bayessche Schätzung betrachtet die ganze Posterior-Verteilung:
        \[
            p(\theta \mid x_1,\dots,x_N).
        \]
    \end{itemize}

    Die Kernaussage lautet:

    \[
        \boxed{
            \text{Lernen bedeutet oft, unbekannte Parameter aus Daten zu schätzen.}
        }
    \]
\end{answerbox}

%%%%%%%%%%%%%%%%%%%%%%%%%%%%%%%%%%%%%%%%%%%%%%%%%%%%%%%%%%%%%%%%%%%
\section{Lecture 4: Maximum Likelihood and Regression}
%%%%%%%%%%%%%%%%%%%%%%%%%%%%%%%%%%%%%%%%%%%%%%%%%%%%%%%%%%%%%%%%%%%

\subsection{Overview}

\begin{notebox}
    Lecture 4 verbindet die Maximum-Likelihood-Schätzung aus Lecture 3 mit
    Regressions- und Klassifikationsmodellen.

    Die zentralen Themen sind:

    \begin{itemize}
        \item Wiederholung von Schätzern und Maximum-Likelihood-Schätzern,
        \item Maximum-Likelihood-Schätzung für die Gleichverteilung,
        \item Maximum-Likelihood-Schätzung für Gauß-Verteilungen,
        \item Diskussion der Maximum-Likelihood-Methode,
        \item Regression mit gelabelten Daten,
        \item lineare Regression,
        \item Overfitting,
        \item Maximum-a-posteriori-Schätzung und Ridge Regression,
        \item Klassifikation als Regression,
        \item logistische Regression,
        \item Gradient Descent.
    \end{itemize}
\end{notebox}

\begin{answerbox}
    Die Kernaussage von Lecture 4 ist:

    \[
        \boxed{
            \text{Viele Lernprobleme lassen sich als Parameterschätzung
            mit Maximum Likelihood oder MAP formulieren.}
        }
    \]

    Regression und Klassifikation können dadurch probabilistisch interpretiert
    werden.
\end{answerbox}

%%%%%%%%%%%%%%%%%%%%%%%%%%%%%%%%%%%%%%%%%%%%%%%%%%%%%%%%%%%%%%%%%%%
\subsection{Recap: Estimators}
%%%%%%%%%%%%%%%%%%%%%%%%%%%%%%%%%%%%%%%%%%%%%%%%%%%%%%%%%%%%%%%%%%%

\begin{notebox}
    Ein Schätzer ist eine Funktion der Daten, mit der ein unbekannter Parameter
    einer Verteilung geschätzt wird.

    Sei \(\theta\) ein unbekannter Parameter und seien
    \[
        X_1,\dots,X_N
    \]
    Zufallsvariablen.

    Ein Schätzer ist dann:
    \[
        \hat\theta
        =
        \hat\theta(X_1,\dots,X_N).
    \]
\end{notebox}

\begin{solutionbox}
    Wichtig ist, dass ein Schätzer selbst eine Zufallsvariable ist, weil er von
    zufälligen Daten abhängt.

    Für konkrete beobachtete Daten
    \[
        x_1,\dots,x_N
    \]
    erhält man einen konkreten Wert:
    \[
        \hat\theta(x_1,\dots,x_N).
    \]

    Vor der Beobachtung der Daten ist \(\hat\theta\) aber zufällig.
\end{solutionbox}

\begin{answerbox}
    Ein Schätzer ist:

    \[
        \boxed{
            \hat\theta
            =
            \hat\theta(X_1,\dots,X_N)
        }
    \]

    Ein guter Schätzer sollte den wahren Parameter \(\theta\) möglichst genau
    approximieren.
\end{answerbox}

%%%%%%%%%%%%%%%%%%%%%%%%%%%%%%%%%%%%%%%%%%%%%%%%%%%%%%%%%%%%%%%%%%%
\subsection{Recap: Maximum Likelihood Estimators}
%%%%%%%%%%%%%%%%%%%%%%%%%%%%%%%%%%%%%%%%%%%%%%%%%%%%%%%%%%%%%%%%%%%

\begin{notebox}
    Die Maximum-Likelihood-Methode wählt den Parameter, unter dem die
    beobachteten Daten am wahrscheinlichsten sind.

    Für Daten
    \[
        x_1,\dots,x_N
    \]
    und Parameter \(\theta\) ist die Likelihood:
    \[
        L(\theta)
        =
        p(x_1,\dots,x_N \mid \theta).
    \]
\end{notebox}

\begin{calculationbox}
    Unter der i.i.d.-Annahme gilt:
    \[
        p(x_1,\dots,x_N \mid \theta)
        =
        \prod_{n=1}^{N}
        p(x_n \mid \theta).
    \]

    Der Maximum-Likelihood-Schätzer ist:
    \[
        \hat\theta_{\mathrm{ML}}
        =
        \operatorname*{arg\,max}_{\theta}
        p(x_1,\dots,x_N \mid \theta).
    \]

    Für i.i.d. Daten:
    \[
        \hat\theta_{\mathrm{ML}}
        =
        \operatorname*{arg\,max}_{\theta}
        \prod_{n=1}^{N}
        p(x_n \mid \theta).
    \]
\end{calculationbox}

\begin{calculationbox}
    Meistens verwendet man die Log-Likelihood:
    \[
        \ell(\theta)
        =
        \log L(\theta).
    \]

    Da der Logarithmus streng monoton wachsend ist, gilt:
    \[
        \operatorname*{arg\,max}_{\theta} L(\theta)
        =
        \operatorname*{arg\,max}_{\theta} \ell(\theta).
    \]

    Für i.i.d. Daten folgt:
    \[
        \ell(\theta)
        =
        \sum_{n=1}^{N}
        \log p(x_n \mid \theta).
    \]
\end{calculationbox}

\begin{answerbox}
    Maximum Likelihood:

    \[
        \boxed{
            \hat\theta_{\mathrm{ML}}
            =
            \operatorname*{arg\,max}_{\theta}
            \prod_{n=1}^{N}
            p(x_n \mid \theta)
        }
    \]

    Äquivalent in Log-Form:

    \[
        \boxed{
            \hat\theta_{\mathrm{ML}}
            =
            \operatorname*{arg\,max}_{\theta}
            \sum_{n=1}^{N}
            \log p(x_n \mid \theta)
        }
    \]
\end{answerbox}

%%%%%%%%%%%%%%%%%%%%%%%%%%%%%%%%%%%%%%%%%%%%%%%%%%%%%%%%%%%%%%%%%%%
\subsection{Uniform Distribution as Maximum Likelihood Estimation}
%%%%%%%%%%%%%%%%%%%%%%%%%%%%%%%%%%%%%%%%%%%%%%%%%%%%%%%%%%%%%%%%%%%

\begin{notebox}
    Wir betrachten Daten
    \[
        x_1,\dots,x_N
    \]
    aus einer Gleichverteilung:
    \[
        X \sim \operatorname{Unif}(a,b).
    \]

    Die Dichte lautet:
    \[
        p(x \mid a,b)
        =
        \begin{cases}
            \dfrac{1}{b-a}, & a \leq x \leq b,\\[0.8em]
            0, & \text{sonst}.
        \end{cases}
    \]
\end{notebox}

\begin{calculationbox}
    Für i.i.d. Daten ist die Likelihood:
    \[
        L(a,b)
        =
        \prod_{n=1}^{N}
        p(x_n \mid a,b).
    \]

    Wenn alle Datenpunkte im Intervall \([a,b]\) liegen, gilt:
    \[
        L(a,b)
        =
        \prod_{n=1}^{N}
        \frac{1}{b-a}
        =
        \frac{1}{(b-a)^N}.
    \]

    Wenn mindestens ein Datenpunkt außerhalb von \([a,b]\) liegt, ist:
    \[
        L(a,b)=0.
    \]
\end{calculationbox}

\begin{calculationbox}
    Um die Likelihood zu maximieren, muss das Intervall \([a,b]\) alle Daten
    enthalten.

    Gleichzeitig soll \(b-a\) möglichst klein sein, da
    \[
        L(a,b)
        =
        \frac{1}{(b-a)^N}
    \]
    größer wird, wenn \(b-a\) kleiner wird.

    Das kleinste Intervall, das alle Daten enthält, ist:
    \[
        a = x_{\min}
    \]
    und
    \[
        b = x_{\max}.
    \]
\end{calculationbox}

\begin{answerbox}
    Für die Gleichverteilung liefert Maximum Likelihood:

    \[
        \boxed{
            \hat a_{\mathrm{ML}}
            =
            x_{\min}
        }
    \]

    \[
        \boxed{
            \hat b_{\mathrm{ML}}
            =
            x_{\max}
        }
    \]

    Das entspricht der Min-Max-Methode aus Lecture 3.
\end{answerbox}

%%%%%%%%%%%%%%%%%%%%%%%%%%%%%%%%%%%%%%%%%%%%%%%%%%%%%%%%%%%%%%%%%%%
\subsection{Testing Maximum Likelihood Estimators for Gaussian Distributions}
%%%%%%%%%%%%%%%%%%%%%%%%%%%%%%%%%%%%%%%%%%%%%%%%%%%%%%%%%%%%%%%%%%%

\begin{notebox}
    Ein wichtiges Testbeispiel für Maximum Likelihood ist die Normalverteilung.

    Sei:
    \[
        X \sim \mathcal{N}(\mu,\sigma^2).
    \]

    Die Dichte lautet:
    \[
        p(x \mid \mu,\sigma^2)
        =
        \frac{1}{\sqrt{2\pi\sigma^2}}
        \exp\left(
            -\frac{(x-\mu)^2}{2\sigma^2}
        \right).
    \]
\end{notebox}

\begin{calculationbox}
    Für i.i.d. Daten \(x_1,\dots,x_N\) ist die Likelihood:
    \[
        L(\mu,\sigma^2)
        =
        \prod_{n=1}^{N}
        \frac{1}{\sqrt{2\pi\sigma^2}}
        \exp\left(
            -\frac{(x_n-\mu)^2}{2\sigma^2}
        \right).
    \]

    Die Log-Likelihood ist:
    \[
        \ell(\mu,\sigma^2)
        =
        -\frac{N}{2}\log(2\pi)
        -\frac{N}{2}\log(\sigma^2)
        -\frac{1}{2\sigma^2}
        \sum_{n=1}^{N}(x_n-\mu)^2.
    \]
\end{calculationbox}

\begin{calculationbox}
    Ableitung nach \(\mu\):

    \[
        \frac{\partial \ell}{\partial \mu}
        =
        \frac{1}{\sigma^2}
        \sum_{n=1}^{N}
        (x_n-\mu).
    \]

    Setze gleich null:
    \[
        \frac{1}{\sigma^2}
        \sum_{n=1}^{N}
        (x_n-\mu)
        =
        0.
    \]

    Also:
    \[
        \sum_{n=1}^{N}x_n
        -
        N\mu
        =
        0.
    \]

    Damit:
    \[
        \hat\mu_{\mathrm{ML}}
        =
        \frac{1}{N}
        \sum_{n=1}^{N}
        x_n
        =
        \bar x.
    \]
\end{calculationbox}

\begin{calculationbox}
    Ableitung nach \(\sigma^2\) liefert:

    \[
        \hat\sigma^2_{\mathrm{ML}}
        =
        \frac{1}{N}
        \sum_{n=1}^{N}
        (x_n-\hat\mu_{\mathrm{ML}})^2.
    \]

    Dieser Schätzer ist der ML-Schätzer für die Varianz, aber er ist im
    Allgemeinen verzerrt.

    Der unverzerrte Schätzer verwendet stattdessen:
    \[
        \frac{1}{N-1}
        \sum_{n=1}^{N}
        (x_n-\bar x)^2.
    \]
\end{calculationbox}

\begin{answerbox}
    Für eine Gauß-Verteilung gilt:

    \[
        \boxed{
            \hat\mu_{\mathrm{ML}}
            =
            \bar x
            =
            \frac{1}{N}
            \sum_{n=1}^{N}
            x_n
        }
    \]

    \[
        \boxed{
            \hat\sigma^2_{\mathrm{ML}}
            =
            \frac{1}{N}
            \sum_{n=1}^{N}
            (x_n-\bar x)^2
        }
    \]

    Der ML-Varianzschätzer verwendet \(1/N\), nicht \(1/(N-1)\).
\end{answerbox}

%%%%%%%%%%%%%%%%%%%%%%%%%%%%%%%%%%%%%%%%%%%%%%%%%%%%%%%%%%%%%%%%%%%
\subsection{Discussion of the Maximum Likelihood Method}
%%%%%%%%%%%%%%%%%%%%%%%%%%%%%%%%%%%%%%%%%%%%%%%%%%%%%%%%%%%%%%%%%%%

\begin{notebox}
    Maximum Likelihood ist eine sehr allgemeine Methode zur Parameterschätzung.

    Man braucht dafür:

    \begin{itemize}
        \item ein probabilistisches Modell \(p(x \mid \theta)\),
        \item beobachtete Daten \(x_1,\dots,x_N\),
        \item eine Optimierung der Likelihood oder Log-Likelihood.
    \end{itemize}
\end{notebox}

\begin{solutionbox}
    Vorteile von Maximum Likelihood:

    \begin{itemize}
        \item systematische Methode,
        \item auf viele Modelle anwendbar,
        \item oft rechnerisch gut handhabbar über die Log-Likelihood,
        \item direkte probabilistische Interpretation.
    \end{itemize}

    Mögliche Nachteile:

    \begin{itemize}
        \item kann bei kleinen Datenmengen stark schwanken,
        \item nutzt kein Vorwissen über Parameter,
        \item kann zu Overfitting führen,
        \item Optimierungsproblem kann schwierig sein.
    \end{itemize}
\end{solutionbox}

\begin{answerbox}
    Maximum Likelihood bedeutet:

    \[
        \boxed{
            \text{Wähle den Parameter, unter dem die beobachteten Daten maximal plausibel sind.}
        }
    \]

    MAP erweitert diese Idee um Vorwissen durch einen Prior:
    \[
        \boxed{
            p(\theta).
        }
    \]
\end{answerbox}

%%%%%%%%%%%%%%%%%%%%%%%%%%%%%%%%%%%%%%%%%%%%%%%%%%%%%%%%%%%%%%%%%%%
\subsection{Regression: Problem Definition and Maximum Likelihood for Labeled Data}
%%%%%%%%%%%%%%%%%%%%%%%%%%%%%%%%%%%%%%%%%%%%%%%%%%%%%%%%%%%%%%%%%%%

\begin{notebox}
    Bei Regression möchte man aus Eingaben \(x\) kontinuierliche Zielwerte \(t\)
    vorhersagen.

    Die Trainingsdaten sind gelabelte Paare:
    \[
        \mathcal{D}
        =
        \set{(x_1,t_1),\dots,(x_N,t_N)}.
    \]

    Dabei ist:

    \[
        x_n = \text{Input}
    \]
    und
    \[
        t_n = \text{Target bzw. Label}.
    \]
\end{notebox}

\begin{solutionbox}
    In einem probabilistischen Regressionsmodell beschreibt man nicht nur eine
    deterministische Funktion
    \[
        y(x,w),
    \]
    sondern eine Verteilung über Zielwerte:
    \[
        p(t \mid x,w).
    \]

    Der Parametervektor \(w\) bestimmt das Modell.

    Für gelabelte Daten lautet die Likelihood:
    \[
        L(w)
        =
        p(t_1,\dots,t_N \mid x_1,\dots,x_N,w).
    \]

    Unter der Annahme unabhängiger Daten:
    \[
        L(w)
        =
        \prod_{n=1}^{N}
        p(t_n \mid x_n,w).
    \]
\end{solutionbox}

\begin{answerbox}
    Regression mit Maximum Likelihood:

    \[
        \boxed{
            \hat w_{\mathrm{ML}}
            =
            \operatorname*{arg\,max}_{w}
            \prod_{n=1}^{N}
            p(t_n \mid x_n,w)
        }
    \]

    In Log-Form:

    \[
        \boxed{
            \hat w_{\mathrm{ML}}
            =
            \operatorname*{arg\,max}_{w}
            \sum_{n=1}^{N}
            \log p(t_n \mid x_n,w)
        }
    \]
\end{answerbox}

%%%%%%%%%%%%%%%%%%%%%%%%%%%%%%%%%%%%%%%%%%%%%%%%%%%%%%%%%%%%%%%%%%%
\subsection{Examples of Regression Problems}
%%%%%%%%%%%%%%%%%%%%%%%%%%%%%%%%%%%%%%%%%%%%%%%%%%%%%%%%%%%%%%%%%%%

\begin{notebox}
    Regression wird verwendet, wenn die Zielvariable kontinuierlich ist.

    Beispiele:

    \begin{itemize}
        \item Vorhersage von Hauspreisen,
        \item Vorhersage von Temperatur,
        \item Vorhersage von Blutdruck oder medizinischen Messwerten,
        \item Vorhersage von Verkaufszahlen,
        \item Schätzung einer physikalischen Größe aus Messdaten.
    \end{itemize}
\end{notebox}

\begin{answerbox}
    Regression bedeutet:

    \[
        \boxed{
            x \longmapsto t
        }
    \]

    wobei \(t\) ein kontinuierlicher Zielwert ist.

    Im Unterschied zur Klassifikation ist das Ziel nicht eine Klasse, sondern
    ein numerischer Wert.
\end{answerbox}

%%%%%%%%%%%%%%%%%%%%%%%%%%%%%%%%%%%%%%%%%%%%%%%%%%%%%%%%%%%%%%%%%%%
\subsection{Linear Regression: Definition}
%%%%%%%%%%%%%%%%%%%%%%%%%%%%%%%%%%%%%%%%%%%%%%%%%%%%%%%%%%%%%%%%%%%

\begin{notebox}
    Bei linearer Regression wird der Zielwert durch eine lineare Funktion der
    Eingabe beschrieben.

    Für eine eindimensionale Eingabe:
    \[
        y(x,w)
        =
        w_0 + w_1x.
    \]

    Allgemeiner mit Basisfunktionen \(\phi_j(x)\):
    \[
        y(x,w)
        =
        \sum_{j=0}^{M}
        w_j \phi_j(x).
    \]

    In Vektorform:
    \[
        y(x,w)
        =
        w^T\phi(x).
    \]
\end{notebox}

\begin{solutionbox}
    Die Funktion ist linear in den Parametern \(w_j\), auch wenn die
    Basisfunktionen \(\phi_j(x)\) nichtlinear in \(x\) sein können.

    Beispiele für Basisfunktionen:

    \[
        \phi_0(x)=1,
        \qquad
        \phi_1(x)=x,
        \qquad
        \phi_2(x)=x^2.
    \]

    Dann ist:
    \[
        y(x,w)
        =
        w_0 + w_1x + w_2x^2.
    \]

    Dieses Modell ist nichtlinear in \(x\), aber linear in den Parametern.
\end{solutionbox}

\begin{answerbox}
    Lineare Regression verwendet Modelle der Form:

    \[
        \boxed{
            y(x,w)
            =
            w^T\phi(x)
        }
    \]

    Wichtig:

    \[
        \boxed{
            \text{linear in den Parametern } w.
        }
    \]
\end{answerbox}

%%%%%%%%%%%%%%%%%%%%%%%%%%%%%%%%%%%%%%%%%%%%%%%%%%%%%%%%%%%%%%%%%%%
\subsection{Linear Regression via Maximum Likelihood}
%%%%%%%%%%%%%%%%%%%%%%%%%%%%%%%%%%%%%%%%%%%%%%%%%%%%%%%%%%%%%%%%%%%

\begin{notebox}
    Für die ML-Herleitung der linearen Regression nimmt man meist ein
    Gaußsches Rauschmodell an.

    Das Modell lautet:
    \[
        t_n
        =
        y(x_n,w) + \varepsilon_n,
    \]
    wobei
    \[
        \varepsilon_n
        \sim \mathcal{N}(0,\sigma^2).
    \]

    Daher:
    \[
        p(t_n \mid x_n,w,\sigma^2)
        =
        \mathcal{N}(t_n \mid y(x_n,w),\sigma^2).
    \]
\end{notebox}

\begin{calculationbox}
    Für unabhängige Daten ist die Likelihood:
    \[
        L(w)
        =
        \prod_{n=1}^{N}
        \mathcal{N}(t_n \mid y(x_n,w),\sigma^2).
    \]

    Die Log-Likelihood ist bis auf Konstanten:
    \[
        \ell(w)
        =
        -\frac{1}{2\sigma^2}
        \sum_{n=1}^{N}
        \left(t_n-y(x_n,w)\right)^2
        + \text{const}.
    \]

    Maximierung der Log-Likelihood ist daher äquivalent zur Minimierung der
    Summe quadratischer Fehler:
    \[
        E(w)
        =
        \frac{1}{2}
        \sum_{n=1}^{N}
        \left(t_n-y(x_n,w)\right)^2.
    \]
\end{calculationbox}

\begin{answerbox}
    Lineare Regression mit Gaußschem Rauschen führt zu Least Squares:

    \[
        \boxed{
            \hat w_{\mathrm{ML}}
            =
            \operatorname*{arg\,min}_{w}
            \frac{1}{2}
            \sum_{n=1}^{N}
            \left(t_n-y(x_n,w)\right)^2
        }
    \]

    Das bedeutet:

    \[
        \boxed{
            \text{Least Squares ist Maximum Likelihood unter Gaußschem Rauschen.}
        }
    \]
\end{answerbox}

%%%%%%%%%%%%%%%%%%%%%%%%%%%%%%%%%%%%%%%%%%%%%%%%%%%%%%%%%%%%%%%%%%%
\subsection{Overfitting}
%%%%%%%%%%%%%%%%%%%%%%%%%%%%%%%%%%%%%%%%%%%%%%%%%%%%%%%%%%%%%%%%%%%

\begin{notebox}
    Overfitting tritt auf, wenn ein Modell die Trainingsdaten sehr gut erklärt,
    aber schlecht auf neue Daten generalisiert.

    Typisch ist:

    \[
        \text{Training Error klein, Test Error groß.}
    \]
\end{notebox}

\begin{solutionbox}
    Overfitting entsteht häufig, wenn das Modell zu flexibel ist, also zu viele
    Parameter oder zu komplexe Basisfunktionen verwendet.

    Dann kann das Modell nicht nur die zugrundeliegende Struktur lernen, sondern
    auch Rauschen in den Trainingsdaten.

    In Regression kann dies zum Beispiel passieren, wenn ein Polynom sehr hohen
    Grades verwendet wird.
\end{solutionbox}

\begin{answerbox}
    Overfitting bedeutet:

    \[
        \boxed{
            \text{Das Modell passt die Trainingsdaten zu stark an und generalisiert schlecht.}
        }
    \]

    Gegenmaßnahmen sind zum Beispiel:

    \begin{itemize}
        \item Regularisierung,
        \item einfachere Modelle,
        \item mehr Daten,
        \item Validierungsdaten zur Modellwahl.
    \end{itemize}
\end{answerbox}

%%%%%%%%%%%%%%%%%%%%%%%%%%%%%%%%%%%%%%%%%%%%%%%%%%%%%%%%%%%%%%%%%%%
\subsection{Linear Regression via MAP: Ridge Regression}
%%%%%%%%%%%%%%%%%%%%%%%%%%%%%%%%%%%%%%%%%%%%%%%%%%%%%%%%%%%%%%%%%%%

\begin{notebox}
    MAP-Schätzung erweitert Maximum Likelihood um einen Prior auf die Parameter.

    Für lineare Regression verwendet man häufig einen Gauß-Prior:
    \[
        p(w)
        =
        \mathcal{N}(w \mid 0,\alpha^{-1}I).
    \]

    Dieser Prior bevorzugt kleine Gewichte.
\end{notebox}

\begin{calculationbox}
    Die MAP-Schätzung maximiert:
    \[
        p(w \mid \mathcal{D})
        \propto
        p(\mathcal{D} \mid w)p(w).
    \]

    In Log-Form:
    \[
        \log p(w \mid \mathcal{D})
        =
        \log p(\mathcal{D} \mid w)
        +
        \log p(w)
        +
        \text{const}.
    \]

    Der Likelihood-Term mit Gaußschem Rauschen führt zu:
    \[
        -\frac{1}{2\sigma^2}
        \sum_{n=1}^{N}
        \left(t_n-y(x_n,w)\right)^2.
    \]

    Der Gauß-Prior führt zu:
    \[
        -\frac{\alpha}{2}
        w^Tw.
    \]
\end{calculationbox}

\begin{calculationbox}
    Maximierung des Log-Posteriors ist äquivalent zur Minimierung:

    \[
        E_{\mathrm{ridge}}(w)
        =
        \frac{1}{2}
        \sum_{n=1}^{N}
        \left(t_n-y(x_n,w)\right)^2
        +
        \frac{\lambda}{2}
        w^Tw.
    \]

    Dabei ist \(\lambda\) ein Regularisierungsparameter.

    Der zusätzliche Term
    \[
        \frac{\lambda}{2}w^Tw
    \]
    bestraft große Gewichte.
\end{calculationbox}

\begin{answerbox}
    Ridge Regression ergibt sich als MAP-Schätzung mit Gauß-Prior auf \(w\):

    \[
        \boxed{
            \hat w_{\mathrm{MAP}}
            =
            \operatorname*{arg\,min}_{w}
            \left[
                \frac{1}{2}
                \sum_{n=1}^{N}
                \left(t_n-y(x_n,w)\right)^2
                +
                \frac{\lambda}{2}
                w^Tw
            \right]
        }
    \]

    Der Regularisierungsterm reduziert Overfitting.
\end{answerbox}

%%%%%%%%%%%%%%%%%%%%%%%%%%%%%%%%%%%%%%%%%%%%%%%%%%%%%%%%%%%%%%%%%%%
\subsection{Classification as Regression}
%%%%%%%%%%%%%%%%%%%%%%%%%%%%%%%%%%%%%%%%%%%%%%%%%%%%%%%%%%%%%%%%%%%

\begin{notebox}
    Klassifikation kann teilweise als Regressionsproblem formuliert werden.

    Statt einen kontinuierlichen Zielwert vorherzusagen, soll das Modell eine
    Klassenwahrscheinlichkeit ausgeben.

    Für binäre Klassifikation:
    \[
        t \in \{0,1\}.
    \]

    Gesucht ist:
    \[
        P(t=1 \mid x).
    \]
\end{notebox}

\begin{solutionbox}
    Ein lineares Modell
    \[
        y(x,w)=w^Tx
    \]
    kann beliebige reelle Werte annehmen.

    Für Wahrscheinlichkeiten brauchen wir aber Werte im Intervall:
    \[
        [0,1].
    \]

    Daher verwendet man eine nichtlineare Funktion, die reelle Zahlen auf
    Wahrscheinlichkeiten abbildet.

    Bei logistischer Regression ist das die Sigmoidfunktion.
\end{solutionbox}

\begin{answerbox}
    Klassifikation als Regression bedeutet:

    \[
        \boxed{
            \text{Wir modellieren eine Klassenwahrscheinlichkeit als Funktion von } x.
        }
    \]

    Für binäre Klassifikation:

    \[
        \boxed{
            P(t=1 \mid x,w)
            \in [0,1].
        }
    \]
\end{answerbox}

%%%%%%%%%%%%%%%%%%%%%%%%%%%%%%%%%%%%%%%%%%%%%%%%%%%%%%%%%%%%%%%%%%%
\subsection{Logistic Regression: Definition}
%%%%%%%%%%%%%%%%%%%%%%%%%%%%%%%%%%%%%%%%%%%%%%%%%%%%%%%%%%%%%%%%%%%

\begin{notebox}
    Logistische Regression ist ein Modell für binäre Klassifikation.

    Die Zielvariable ist:
    \[
        t \in \{0,1\}.
    \]

    Das Modell beschreibt:
    \[
        P(t=1 \mid x,w).
    \]
\end{notebox}

\begin{calculationbox}
    Zuerst bildet man eine lineare Kombination:
    \[
        a
        =
        w^T x.
    \]

    Dann wird diese durch die Sigmoidfunktion transformiert:
    \[
        \sigma(a)
        =
        \frac{1}{1+\exp(-a)}.
    \]

    Damit:
    \[
        P(t=1 \mid x,w)
        =
        \sigma(w^Tx).
    \]

    Und:
    \[
        P(t=0 \mid x,w)
        =
        1-\sigma(w^Tx).
    \]
\end{calculationbox}

\begin{answerbox}
    Logistische Regression:

    \[
        \boxed{
            P(t=1 \mid x,w)
            =
            \sigma(w^Tx)
        }
    \]

    mit

    \[
        \boxed{
            \sigma(a)
            =
            \frac{1}{1+\exp(-a)}
        }
    \]

    Die Sigmoidfunktion macht aus einem reellen Wert eine Wahrscheinlichkeit.
\end{answerbox}

%%%%%%%%%%%%%%%%%%%%%%%%%%%%%%%%%%%%%%%%%%%%%%%%%%%%%%%%%%%%%%%%%%%
\subsection{Excursion: Gradient Descent}
%%%%%%%%%%%%%%%%%%%%%%%%%%%%%%%%%%%%%%%%%%%%%%%%%%%%%%%%%%%%%%%%%%%

\begin{notebox}
    Viele Modelle haben keine geschlossene Lösung für die optimalen Parameter.

    Dann verwendet man iterative Optimierungsverfahren wie Gradient Descent.

    Ziel ist es, eine Fehlerfunktion \(E(w)\) zu minimieren.
\end{notebox}

\begin{calculationbox}
    Gradient Descent aktualisiert die Parameter in Richtung des negativen
    Gradienten:

    \[
        w^{(\tau+1)}
        =
        w^{(\tau)}
        -
        \eta
        \nabla E(w^{(\tau)}).
    \]

    Dabei ist:

    \begin{itemize}
        \item \(\tau\) der Iterationsindex,
        \item \(\eta>0\) die Lernrate,
        \item \(\nabla E(w)\) der Gradient der Fehlerfunktion.
    \end{itemize}
\end{calculationbox}

\begin{solutionbox}
    Der Gradient zeigt in Richtung des stärksten Anstiegs der Funktion.

    Da man minimieren möchte, geht man in die entgegengesetzte Richtung:
    \[
        -\nabla E(w).
    \]

    Die Lernrate \(\eta\) bestimmt, wie groß jeder Schritt ist.

    Ist \(\eta\) zu groß, kann das Verfahren instabil werden.

    Ist \(\eta\) zu klein, konvergiert das Verfahren langsam.
\end{solutionbox}

\begin{answerbox}
    Gradient Descent:

    \[
        \boxed{
            w^{(\tau+1)}
            =
            w^{(\tau)}
            -
            \eta
            \nabla E(w^{(\tau)})
        }
    \]

    Es ist ein iteratives Verfahren zur Minimierung einer Fehlerfunktion.
\end{answerbox}

%%%%%%%%%%%%%%%%%%%%%%%%%%%%%%%%%%%%%%%%%%%%%%%%%%%%%%%%%%%%%%%%%%%
\subsection{Logistic Regression via Maximum Likelihood}
%%%%%%%%%%%%%%%%%%%%%%%%%%%%%%%%%%%%%%%%%%%%%%%%%%%%%%%%%%%%%%%%%%%

\begin{notebox}
    Logistische Regression kann mit Maximum Likelihood hergeleitet werden.

    Für binäre Labels
    \[
        t_n \in \{0,1\}
    \]
    modelliert man:
    \[
        P(t_n=1 \mid x_n,w)
        =
        y_n
        =
        \sigma(w^Tx_n).
    \]
\end{notebox}

\begin{calculationbox}
    Die Bernoulli-Wahrscheinlichkeit für ein Label \(t_n\) ist:
    \[
        p(t_n \mid x_n,w)
        =
        y_n^{t_n}
        (1-y_n)^{1-t_n}.
    \]

    Für i.i.d. Daten ist die Likelihood:
    \[
        L(w)
        =
        \prod_{n=1}^{N}
        y_n^{t_n}
        (1-y_n)^{1-t_n}.
    \]

    Die Log-Likelihood ist:
    \[
        \ell(w)
        =
        \sum_{n=1}^{N}
        \left[
            t_n\log y_n
            +
            (1-t_n)\log(1-y_n)
        \right].
    \]
\end{calculationbox}

\begin{calculationbox}
    Maximum Likelihood maximiert \(\ell(w)\).

    Äquivalent minimiert man die negative Log-Likelihood:
    \[
        E(w)
        =
        -
        \sum_{n=1}^{N}
        \left[
            t_n\log y_n
            +
            (1-t_n)\log(1-y_n)
        \right].
    \]

    Diese Fehlerfunktion heißt Cross-Entropy Loss.
\end{calculationbox}

\begin{answerbox}
    Logistische Regression mit Maximum Likelihood führt zur Cross-Entropy:

    \[
        \boxed{
            E(w)
            =
            -
            \sum_{n=1}^{N}
            \left[
                t_n\log y_n
                +
                (1-t_n)\log(1-y_n)
            \right]
        }
    \]

    mit

    \[
        \boxed{
            y_n
            =
            \sigma(w^Tx_n).
        }
    \]

    Dieser Fehler wird typischerweise mit Gradient Descent minimiert.
\end{answerbox}

%%%%%%%%%%%%%%%%%%%%%%%%%%%%%%%%%%%%%%%%%%%%%%%%%%%%%%%%%%%%%%%%%%%
\subsection{Summary}
%%%%%%%%%%%%%%%%%%%%%%%%%%%%%%%%%%%%%%%%%%%%%%%%%%%%%%%%%%%%%%%%%%%

\begin{answerbox}
    Lecture 4 zeigt, wie Maximum Likelihood und MAP-Schätzung zu konkreten
    Machine-Learning-Modellen führen.

    Die wichtigsten Punkte sind:

    \begin{itemize}
        \item Maximum Likelihood wählt Parameter, die die Daten am plausibelsten
        machen:
        \[
            \hat\theta_{\mathrm{ML}}
            =
            \operatorname*{arg\,max}_{\theta}
            \sum_{n=1}^{N}
            \log p(x_n \mid \theta).
        \]

        \item Für eine Gleichverteilung liefert ML:
        \[
            \hat a_{\mathrm{ML}}=x_{\min},
            \qquad
            \hat b_{\mathrm{ML}}=x_{\max}.
        \]

        \item Für eine Gauß-Verteilung liefert ML:
        \[
            \hat\mu_{\mathrm{ML}}=\bar x,
            \qquad
            \hat\sigma^2_{\mathrm{ML}}
            =
            \frac{1}{N}
            \sum_{n=1}^{N}(x_n-\bar x)^2.
        \]

        \item Regression verwendet gelabelte Daten:
        \[
            (x_n,t_n).
        \]

        \item Lineare Regression modelliert:
        \[
            y(x,w)=w^T\phi(x).
        \]

        \item Lineare Regression mit Gaußschem Rauschen führt zu Least Squares:
        \[
            \hat w_{\mathrm{ML}}
            =
            \operatorname*{arg\,min}_{w}
            \frac{1}{2}
            \sum_{n=1}^{N}
            (t_n-y(x_n,w))^2.
        \]

        \item MAP mit Gauß-Prior auf \(w\) führt zu Ridge Regression:
        \[
            \hat w_{\mathrm{MAP}}
            =
            \operatorname*{arg\,min}_{w}
            \left[
                \frac{1}{2}
                \sum_{n=1}^{N}
                (t_n-y(x_n,w))^2
                +
                \frac{\lambda}{2}w^Tw
            \right].
        \]

        \item Logistische Regression modelliert:
        \[
            P(t=1 \mid x,w)
            =
            \sigma(w^Tx).
        \]

        \item Die Sigmoidfunktion ist:
        \[
            \sigma(a)
            =
            \frac{1}{1+\exp(-a)}.
        \]

        \item Logistische Regression mit ML führt zur Cross-Entropy:
        \[
            E(w)
            =
            -
            \sum_{n=1}^{N}
            \left[
                t_n\log y_n
                +
                (1-t_n)\log(1-y_n)
            \right].
        \]

        \item Gradient Descent aktualisiert Parameter durch:
        \[
            w^{(\tau+1)}
            =
            w^{(\tau)}
            -
            \eta
            \nabla E(w^{(\tau)}).
        \]
    \end{itemize}

    Die Kernaussage lautet:

    \[
        \boxed{
            \text{Regression und Klassifikation können als probabilistische
            Parameterschätzung formuliert werden.}
        }
    \]
\end{answerbox}

%%%%%%%%%%%%%%%%%%%%%%%%%%%%%%%%%%%%%%%%%%%%%%%%%%%%%%%%%%%%%%%%%%%
\section{Lecture 5: Neural Networks -- Perceptrons and Multi-Layer Perceptrons}
%%%%%%%%%%%%%%%%%%%%%%%%%%%%%%%%%%%%%%%%%%%%%%%%%%%%%%%%%%%%%%%%%%%

\subsection{Overview}

\begin{notebox}
    Lecture 5 führt neuronale Netze ein. Der erste Teil behandelt einfache
    Perzeptrons, der zweite Teil erweitert diese Idee zu Multi-Layer
    Perceptrons.

    Die zentralen Themen sind:

    \begin{itemize}
        \item biologische und künstliche neuronale Netze,
        \item Definition des Perzeptrons,
        \item Lernen in Perzeptrons,
        \item Perzeptrons als Klassifikatoren und ihre Grenzen,
        \item das XOR-Problem,
        \item Erweiterung zu Multi-Layer Perceptrons,
        \item Herleitung von Lernregeln für MLPs,
        \item mechanistische Sicht auf Lernen durch Backpropagation.
    \end{itemize}
\end{notebox}

\begin{answerbox}
    Die Kernaussage von Lecture 5 ist:

    \[
        \boxed{
            \text{Neuronale Netze lernen Funktionen durch gewichtete
            Kombinationen, Nichtlinearitäten und schrittweise Anpassung ihrer Parameter.}
        }
    \]

    Ein einzelnes Perzeptron kann nur linear separierbare Probleme lösen.
    Multi-Layer Perceptrons erweitern diese Idee und können deutlich komplexere
    Funktionen darstellen.
\end{answerbox}

%%%%%%%%%%%%%%%%%%%%%%%%%%%%%%%%%%%%%%%%%%%%%%%%%%%%%%%%%%%%%%%%%%%
\subsection{Part A: Biological and Artificial Neural Networks}
%%%%%%%%%%%%%%%%%%%%%%%%%%%%%%%%%%%%%%%%%%%%%%%%%%%%%%%%%%%%%%%%%%%

\begin{notebox}
    Die Motivation für künstliche neuronale Netze stammt historisch aus
    biologischen Nervensystemen.

    Ein biologisches Neuron erhält Eingangssignale über Dendriten, verarbeitet
    diese im Zellkörper und sendet bei ausreichender Aktivierung ein Signal über
    das Axon weiter.

    Die grobe Analogie in künstlichen neuronalen Netzen ist:

    \begin{itemize}
        \item Eingänge entsprechen Eingangssignalen,
        \item Gewichte entsprechen der Stärke von Verbindungen,
        \item eine gewichtete Summe entspricht der Integration von Signalen,
        \item eine Aktivierungsfunktion entscheidet über die Ausgabe.
    \end{itemize}
\end{notebox}

\begin{solutionbox}
    Ein künstliches Neuron erhält Eingaben
    \[
        x_1,\dots,x_D.
    \]

    Jede Eingabe hat ein Gewicht
    \[
        w_1,\dots,w_D.
    \]

    Zusätzlich gibt es oft einen Bias-Term \(b\). Zuerst wird eine gewichtete
    Summe gebildet:
    \[
        a
        =
        \sum_{j=1}^{D} w_j x_j + b.
    \]

    Danach wird diese Aktivierung durch eine Funktion \(g\) transformiert:
    \[
        y
        =
        g(a).
    \]

    Das ist die Grundstruktur eines künstlichen Neurons.
\end{solutionbox}

\begin{answerbox}
    Ein künstliches Neuron berechnet:

    \[
        \boxed{
            a
            =
            \sum_{j=1}^{D} w_j x_j + b
        }
    \]

    und anschließend:

    \[
        \boxed{
            y
            =
            g(a)
        }
    \]

    Dabei sind \(w_j\) die Gewichte, \(b\) der Bias und \(g\) die
    Aktivierungsfunktion.
\end{answerbox}

%%%%%%%%%%%%%%%%%%%%%%%%%%%%%%%%%%%%%%%%%%%%%%%%%%%%%%%%%%%%%%%%%%%
\subsection{Definition of the Perceptron}
%%%%%%%%%%%%%%%%%%%%%%%%%%%%%%%%%%%%%%%%%%%%%%%%%%%%%%%%%%%%%%%%%%%

\begin{notebox}
    Das Perzeptron ist eines der einfachsten künstlichen neuronalen Modelle.

    Es ist ein binärer Klassifikator. Das bedeutet: Es entscheidet zwischen zwei
    Klassen.

    Die Eingabe ist ein Vektor:
    \[
        x
        =
        (x_1,\dots,x_D)^T.
    \]

    Das Perzeptron besitzt Gewichte:
    \[
        w
        =
        (w_1,\dots,w_D)^T
    \]
    und einen Bias \(b\).
\end{notebox}

\begin{calculationbox}
    Zuerst berechnet das Perzeptron die Aktivierung:
    \[
        a
        =
        w^T x + b.
    \]

    Dann wird eine Schwellenfunktion angewendet.

    Eine typische Definition ist:
    \[
        y
        =
        \begin{cases}
            1, & w^T x + b \geq 0,\\
            0, & w^T x + b < 0.
        \end{cases}
    \]

    Alternativ verwendet man Labels \(\{-1,+1\}\):
    \[
        y
        =
        \operatorname{sign}(w^T x + b).
    \]
\end{calculationbox}

\begin{solutionbox}
    Die Entscheidungsgrenze des Perzeptrons ist durch
    \[
        w^T x + b = 0
    \]
    gegeben.

    In zwei Dimensionen ist das eine Gerade.

    In höheren Dimensionen ist es eine Hyperebene.

    Das Perzeptron klassifiziert danach, auf welcher Seite dieser Hyperebene ein
    Datenpunkt liegt.
\end{solutionbox}

\begin{answerbox}
    Das Perzeptron ist definiert durch:

    \[
        \boxed{
            a = w^T x + b
        }
    \]

    und

    \[
        \boxed{
            y =
            \begin{cases}
                1, & a \geq 0,\\
                0, & a < 0.
            \end{cases}
        }
    \]

    Die Entscheidungsgrenze ist:

    \[
        \boxed{
            w^T x + b = 0.
        }
    \]
\end{answerbox}

%%%%%%%%%%%%%%%%%%%%%%%%%%%%%%%%%%%%%%%%%%%%%%%%%%%%%%%%%%%%%%%%%%%
\subsection{Learning in Perceptrons}
%%%%%%%%%%%%%%%%%%%%%%%%%%%%%%%%%%%%%%%%%%%%%%%%%%%%%%%%%%%%%%%%%%%

\begin{notebox}
    Lernen bedeutet beim Perzeptron, die Gewichte \(w\) und den Bias \(b\)
    so anzupassen, dass die Trainingsdaten möglichst korrekt klassifiziert
    werden.

    Die Trainingsdaten sind gelabelte Beispiele:
    \[
        \mathcal{D}
        =
        \set{(x^{(1)},t^{(1)}),\dots,(x^{(N)},t^{(N)})}.
    \]

    Dabei ist \(t^{(n)}\) das gewünschte Label.
\end{notebox}

\begin{solutionbox}
    Das Perzeptron macht für einen Datenpunkt \(x^{(n)}\) eine Vorhersage:
    \[
        y^{(n)}
        =
        f(w^T x^{(n)} + b).
    \]

    Wenn die Vorhersage korrekt ist, müssen die Gewichte nicht verändert werden.

    Wenn die Vorhersage falsch ist, werden die Gewichte so angepasst, dass die
    Entscheidung beim nächsten Mal eher in Richtung des richtigen Labels geht.
\end{solutionbox}

\begin{calculationbox}
    Für Labels \(t \in \{-1,+1\}\) und Perzeptron-Ausgabe
    \[
        y = \operatorname{sign}(w^T x + b)
    \]
    kann die Perzeptron-Lernregel geschrieben werden als:

    \[
        w_{\text{new}}
        =
        w_{\text{old}}
        +
        \eta t x
    \]
    falls der Punkt falsch klassifiziert wurde.

    Für den Bias:
    \[
        b_{\text{new}}
        =
        b_{\text{old}}
        +
        \eta t.
    \]

    Dabei ist \(\eta>0\) die Lernrate.
\end{calculationbox}

\begin{sidecalcbox}
    Warum macht diese Regel Sinn?

    Wenn \(t=+1\), der Punkt aber falsch klassifiziert wurde, dann ist
    \[
        w^T x + b < 0.
    \]

    Durch
    \[
        w \leftarrow w + \eta x
    \]
    wird die Aktivierung
    \[
        w^T x + b
    \]
    für diesen Punkt größer.

    Wenn \(t=-1\), der Punkt aber falsch klassifiziert wurde, dann ist
    \[
        w^T x + b > 0.
    \]

    Durch
    \[
        w \leftarrow w - \eta x
    \]
    wird die Aktivierung kleiner.
\end{sidecalcbox}

\begin{answerbox}
    Die Perzeptron-Lernregel für falsch klassifizierte Punkte lautet:

    \[
        \boxed{
            w \leftarrow w + \eta t x
        }
    \]

    \[
        \boxed{
            b \leftarrow b + \eta t
        }
    \]

    mit Labels \(t \in \{-1,+1\}\).

    Die Regel verändert die Entscheidungsgrenze so, dass der falsch
    klassifizierte Punkt eher auf die richtige Seite verschoben wird.
\end{answerbox}

%%%%%%%%%%%%%%%%%%%%%%%%%%%%%%%%%%%%%%%%%%%%%%%%%%%%%%%%%%%%%%%%%%%
\subsection{Perceptrons as Classifiers}
%%%%%%%%%%%%%%%%%%%%%%%%%%%%%%%%%%%%%%%%%%%%%%%%%%%%%%%%%%%%%%%%%%%

\begin{notebox}
    Ein Perzeptron ist ein linearer Klassifikator.

    Das bedeutet: Es trennt Klassen durch eine lineare Entscheidungsgrenze.

    In zwei Dimensionen:
    \[
        w_1x_1 + w_2x_2 + b = 0
    \]
    ist eine Gerade.

    In \(D\) Dimensionen:
    \[
        w^T x + b = 0
    \]
    ist eine Hyperebene.
\end{notebox}

\begin{solutionbox}
    Ein Datensatz heißt linear separierbar, wenn es eine Gerade beziehungsweise
    Hyperebene gibt, die die Klassen vollständig trennt.

    Für linear separierbare Daten kann ein Perzeptron prinzipiell eine passende
    Entscheidungsgrenze lernen.

    Wenn die Daten nicht linear separierbar sind, reicht ein einzelnes Perzeptron
    nicht aus.
\end{solutionbox}

\begin{answerbox}
    Das Perzeptron ist ein linearer Klassifikator:

    \[
        \boxed{
            y = \operatorname{sign}(w^Tx+b)
        }
    \]

    Es kann nur Entscheidungsgrenzen der Form

    \[
        \boxed{
            w^Tx+b=0
        }
    \]

    darstellen.

    Daher kann es nur linear separierbare Klassifikationsprobleme lösen.
\end{answerbox}

%%%%%%%%%%%%%%%%%%%%%%%%%%%%%%%%%%%%%%%%%%%%%%%%%%%%%%%%%%%%%%%%%%%
\subsection{Limitations of Perceptrons: The XOR Problem}
%%%%%%%%%%%%%%%%%%%%%%%%%%%%%%%%%%%%%%%%%%%%%%%%%%%%%%%%%%%%%%%%%%%

\begin{notebox}
    Das klassische Beispiel für die Grenzen eines einzelnen Perzeptrons ist das
    XOR-Problem.

    Die XOR-Funktion ist definiert durch:

    \[
        0 \oplus 0 = 0,
        \qquad
        0 \oplus 1 = 1,
    \]
    \[
        1 \oplus 0 = 1,
        \qquad
        1 \oplus 1 = 0.
    \]

    Die positiven Beispiele liegen also diagonal gegenüber, ebenso die negativen.
\end{notebox}

\begin{solutionbox}
    Die vier Punkte sind:

    \[
        (0,0) \mapsto 0,
        \qquad
        (1,0) \mapsto 1,
    \]
    \[
        (0,1) \mapsto 1,
        \qquad
        (1,1) \mapsto 0.
    \]

    Diese Punkte können nicht durch eine einzige Gerade getrennt werden.

    Daher kann ein einzelnes Perzeptron XOR nicht darstellen.
\end{solutionbox}

\begin{answerbox}
    Das XOR-Problem zeigt:

    \[
        \boxed{
            \text{Ein einzelnes Perzeptron kann nur linear separierbare Probleme lösen.}
        }
    \]

    Da XOR nicht linear separierbar ist, gilt:

    \[
        \boxed{
            \text{Ein einzelnes Perzeptron kann XOR nicht lösen.}
        }
    \]

    Dafür braucht man mehrere Schichten beziehungsweise nichtlineare
    Zwischenrepräsentationen.
\end{answerbox}

%%%%%%%%%%%%%%%%%%%%%%%%%%%%%%%%%%%%%%%%%%%%%%%%%%%%%%%%%%%%%%%%%%%
\subsection{Part B: Extension to Multi-Layer Perceptrons}
%%%%%%%%%%%%%%%%%%%%%%%%%%%%%%%%%%%%%%%%%%%%%%%%%%%%%%%%%%%%%%%%%%%

\begin{notebox}
    Ein Multi-Layer Perceptron, kurz MLP, erweitert das einfache Perzeptron durch
    mehrere Schichten.

    Typische Schichten sind:

    \begin{itemize}
        \item Eingabeschicht,
        \item eine oder mehrere versteckte Schichten,
        \item Ausgabeschicht.
    \end{itemize}

    Die versteckten Schichten erlauben es, nichtlineare Zwischenrepräsentationen
    zu lernen.
\end{notebox}

\begin{solutionbox}
    Eine Schicht eines neuronalen Netzes berechnet typischerweise:

    \[
        a^{(\ell)}
        =
        W^{(\ell)} h^{(\ell-1)} + b^{(\ell)}
    \]

    und danach:

    \[
        h^{(\ell)}
        =
        g(a^{(\ell)}).
    \]

    Dabei ist:

    \begin{itemize}
        \item \(\ell\) der Schichtindex,
        \item \(W^{(\ell)}\) die Gewichtsmatrix der Schicht,
        \item \(b^{(\ell)}\) der Bias-Vektor,
        \item \(h^{(\ell-1)}\) die Eingabe der Schicht,
        \item \(h^{(\ell)}\) die Ausgabe der Schicht,
        \item \(g\) eine nichtlineare Aktivierungsfunktion.
    \end{itemize}
\end{solutionbox}

\begin{answerbox}
    Eine MLP-Schicht hat die Form:

    \[
        \boxed{
            a^{(\ell)}
            =
            W^{(\ell)}h^{(\ell-1)} + b^{(\ell)}
        }
    \]

    \[
        \boxed{
            h^{(\ell)}
            =
            g(a^{(\ell)})
        }
    \]

    Durch mehrere solche Schichten entstehen komplexe nichtlineare Funktionen.
\end{answerbox}

%%%%%%%%%%%%%%%%%%%%%%%%%%%%%%%%%%%%%%%%%%%%%%%%%%%%%%%%%%%%%%%%%%%
\subsection{Why Nonlinear Activation Functions Are Necessary}
%%%%%%%%%%%%%%%%%%%%%%%%%%%%%%%%%%%%%%%%%%%%%%%%%%%%%%%%%%%%%%%%%%%

\begin{notebox}
    Die Nichtlinearität in den Aktivierungsfunktionen ist entscheidend.

    Wenn alle Aktivierungsfunktionen linear wären, dann wäre die Verkettung
    mehrerer Schichten wieder nur eine lineare Funktion.
\end{notebox}

\begin{calculationbox}
    Betrachte zwei lineare Schichten ohne Nichtlinearität:

    \[
        h^{(1)}
        =
        W^{(1)}x + b^{(1)}
    \]

    und

    \[
        y
        =
        W^{(2)}h^{(1)} + b^{(2)}.
    \]

    Einsetzen liefert:

    \[
        y
        =
        W^{(2)}(W^{(1)}x + b^{(1)}) + b^{(2)}.
    \]

    Also:

    \[
        y
        =
        W^{(2)}W^{(1)}x
        +
        W^{(2)}b^{(1)}
        +
        b^{(2)}.
    \]

    Das ist wieder von der Form:

    \[
        y
        =
        \widetilde W x + \widetilde b.
    \]

    Also wieder linear.
\end{calculationbox}

\begin{answerbox}
    Ohne Nichtlinearitäten gilt:

    \[
        \boxed{
            \text{Mehrere lineare Schichten ergeben zusammen wieder nur eine lineare Schicht.}
        }
    \]

    Deshalb braucht ein MLP nichtlineare Aktivierungsfunktionen, zum Beispiel
    Sigmoid, \(\tanh\) oder ReLU.
\end{answerbox}

%%%%%%%%%%%%%%%%%%%%%%%%%%%%%%%%%%%%%%%%%%%%%%%%%%%%%%%%%%%%%%%%%%%
\subsection{Learning Rules for Multi-Layer Perceptrons}
%%%%%%%%%%%%%%%%%%%%%%%%%%%%%%%%%%%%%%%%%%%%%%%%%%%%%%%%%%%%%%%%%%%

\begin{notebox}
    Beim Lernen eines MLPs sollen alle Gewichte und Biases so angepasst werden,
    dass eine Fehlerfunktion minimiert wird.

    Sei die Fehlerfunktion:
    \[
        E(\theta),
    \]
    wobei \(\theta\) alle Parameter des Netzes zusammenfasst.

    Die Parameter sind:
    \[
        \theta
        =
        \set{W^{(1)},b^{(1)},W^{(2)},b^{(2)},\dots}.
    \]
\end{notebox}

\begin{solutionbox}
    Das Standardprinzip ist Gradient Descent.

    Für einen Parameter \(\theta_i\) lautet die Update-Regel:
    \[
        \theta_i
        \leftarrow
        \theta_i
        -
        \eta
        \frac{\partial E}{\partial \theta_i}.
    \]

    Für Gewichtsmatrizen schreibt man entsprechend:
    \[
        W^{(\ell)}
        \leftarrow
        W^{(\ell)}
        -
        \eta
        \frac{\partial E}{\partial W^{(\ell)}}.
    \]

    Für Biases:
    \[
        b^{(\ell)}
        \leftarrow
        b^{(\ell)}
        -
        \eta
        \frac{\partial E}{\partial b^{(\ell)}}.
    \]
\end{solutionbox}

\begin{answerbox}
    Lernen in MLPs bedeutet Gradientenabstieg auf einer Fehlerfunktion:

    \[
        \boxed{
            \theta
            \leftarrow
            \theta
            -
            \eta
            \nabla_\theta E(\theta)
        }
    \]

    Für jede Schicht:

    \[
        \boxed{
            W^{(\ell)}
            \leftarrow
            W^{(\ell)}
            -
            \eta
            \frac{\partial E}{\partial W^{(\ell)}}
        }
    \]

    \[
        \boxed{
            b^{(\ell)}
            \leftarrow
            b^{(\ell)}
            -
            \eta
            \frac{\partial E}{\partial b^{(\ell)}}
        }
    \]
\end{answerbox}

%%%%%%%%%%%%%%%%%%%%%%%%%%%%%%%%%%%%%%%%%%%%%%%%%%%%%%%%%%%%%%%%%%%
\subsection{Mechanistic View of Learning: Backpropagation}
%%%%%%%%%%%%%%%%%%%%%%%%%%%%%%%%%%%%%%%%%%%%%%%%%%%%%%%%%%%%%%%%%%%

\begin{notebox}
    Backpropagation ist der Algorithmus, mit dem Gradienten in mehrschichtigen
    neuronalen Netzen effizient berechnet werden.

    Die zentrale Idee ist die Kettenregel der Differentialrechnung.

    Da ein MLP eine Verkettung vieler Funktionen ist, kann man die Ableitungen
    von hinten nach vorne durch das Netzwerk propagieren.
\end{notebox}

\begin{solutionbox}
    Der Lernprozess besteht aus zwei Hauptphasen:

    \begin{enumerate}
        \item \textbf{Forward pass:}
        Die Eingabe wird Schicht für Schicht durch das Netzwerk gerechnet, bis
        eine Ausgabe entsteht.

        \item \textbf{Backward pass:}
        Der Fehler wird von der Ausgabeschicht zurück durch das Netzwerk
        propagiert, um die Gradienten aller Gewichte zu berechnen.
    \end{enumerate}
\end{solutionbox}

\begin{calculationbox}
    Für eine Schicht gilt:

    \[
        a^{(\ell)}
        =
        W^{(\ell)}h^{(\ell-1)} + b^{(\ell)}
    \]

    und

    \[
        h^{(\ell)}
        =
        g(a^{(\ell)}).
    \]

    Um die Gewichte zu aktualisieren, benötigt man:

    \[
        \frac{\partial E}{\partial W^{(\ell)}}.
    \]

    Mit der Kettenregel zerlegt man diese Ableitung in lokale Ableitungen:

    \[
        \frac{\partial E}{\partial W^{(\ell)}}
        =
        \frac{\partial E}{\partial a^{(\ell)}}
        \frac{\partial a^{(\ell)}}{\partial W^{(\ell)}}.
    \]

    Definiert man den Fehlerterm

    \[
        \delta^{(\ell)}
        =
        \frac{\partial E}{\partial a^{(\ell)}},
    \]

    dann gilt für die Gewichtgradienten:

    \[
        \frac{\partial E}{\partial W^{(\ell)}}
        =
        \delta^{(\ell)}
        \left(h^{(\ell-1)}\right)^T.
    \]
\end{calculationbox}

\begin{calculationbox}
    Für die Biases gilt entsprechend:

    \[
        \frac{\partial E}{\partial b^{(\ell)}}
        =
        \delta^{(\ell)}.
    \]

    Die Fehlerterme \(\delta^{(\ell)}\) werden rückwärts berechnet.

    Für eine versteckte Schicht gilt schematisch:

    \[
        \delta^{(\ell)}
        =
        \left(
            (W^{(\ell+1)})^T
            \delta^{(\ell+1)}
        \right)
        \odot
        g'(a^{(\ell)}).
    \]

    Dabei bezeichnet \(\odot\) die komponentenweise Multiplikation.
\end{calculationbox}

\begin{answerbox}
    Backpropagation besteht aus:

    \begin{itemize}
        \item Forward pass: Ausgaben berechnen.
        \item Loss berechnen.
        \item Backward pass: Gradienten mit der Kettenregel berechnen.
        \item Parameter mit Gradient Descent aktualisieren.
    \end{itemize}

    Zentraler Fehlerterm:

    \[
        \boxed{
            \delta^{(\ell)}
            =
            \frac{\partial E}{\partial a^{(\ell)}}
        }
    \]

    Gewichtgradient:

    \[
        \boxed{
            \frac{\partial E}{\partial W^{(\ell)}}
            =
            \delta^{(\ell)}
            \left(h^{(\ell-1)}\right)^T
        }
    \]

    Hidden-layer-Rekursion:

    \[
        \boxed{
            \delta^{(\ell)}
            =
            \left(
                (W^{(\ell+1)})^T
                \delta^{(\ell+1)}
            \right)
            \odot
            g'(a^{(\ell)})
        }
    \]
\end{answerbox}

%%%%%%%%%%%%%%%%%%%%%%%%%%%%%%%%%%%%%%%%%%%%%%%%%%%%%%%%%%%%%%%%%%%
\subsection{Perceptrons, MLPs and the XOR Problem}
%%%%%%%%%%%%%%%%%%%%%%%%%%%%%%%%%%%%%%%%%%%%%%%%%%%%%%%%%%%%%%%%%%%

\begin{notebox}
    Das XOR-Problem motiviert den Übergang vom einzelnen Perzeptron zum MLP.

    Ein einzelnes Perzeptron kann nur eine lineare Entscheidungsgrenze bilden.

    XOR benötigt aber eine nichtlineare Trennung.
\end{notebox}

\begin{solutionbox}
    Ein MLP mit mindestens einer versteckten Schicht kann zuerst eine neue
    Repräsentation der Eingabe lernen.

    In dieser versteckten Repräsentation kann ein Problem linear separierbar
    werden, obwohl es im ursprünglichen Eingaberaum nicht linear separierbar war.

    Genau dadurch können MLPs Probleme wie XOR lösen.
\end{solutionbox}

\begin{answerbox}
    Der Übergang ist:

    \[
        \boxed{
            \text{Perzeptron}
            =
            \text{linearer Klassifikator}
        }
    \]

    \[
        \boxed{
            \text{MLP}
            =
            \text{mehrschichtiges nichtlineares Modell}
        }
    \]

    Ein MLP kann durch versteckte Schichten komplexere Entscheidungsgrenzen
    darstellen.
\end{answerbox}

%%%%%%%%%%%%%%%%%%%%%%%%%%%%%%%%%%%%%%%%%%%%%%%%%%%%%%%%%%%%%%%%%%%
\subsection{Prüfungsrelevante Kurzfassung}
%%%%%%%%%%%%%%%%%%%%%%%%%%%%%%%%%%%%%%%%%%%%%%%%%%%%%%%%%%%%%%%%%%%

\begin{answerbox}
    Lecture 5 behandelt Perzeptrons und Multi-Layer Perceptrons.

    Die wichtigsten Punkte sind:

    \begin{itemize}
        \item Ein künstliches Neuron berechnet:
        \[
            a = w^Tx+b,
            \qquad
            y=g(a).
        \]

        \item Ein Perzeptron ist ein binärer linearer Klassifikator:
        \[
            y
            =
            \operatorname{sign}(w^Tx+b).
        \]

        \item Die Entscheidungsgrenze ist:
        \[
            w^Tx+b=0.
        \]

        \item Für falsch klassifizierte Punkte lautet die Perzeptron-Lernregel:
        \[
            w \leftarrow w+\eta tx,
            \qquad
            b \leftarrow b+\eta t.
        \]

        \item Ein einzelnes Perzeptron kann nur linear separierbare Probleme
        lösen.

        \item XOR ist nicht linear separierbar und kann daher nicht durch ein
        einzelnes Perzeptron gelöst werden.

        \item Ein MLP besteht aus mehreren Schichten:
        \[
            a^{(\ell)}
            =
            W^{(\ell)}h^{(\ell-1)}+b^{(\ell)},
            \qquad
            h^{(\ell)}
            =
            g(a^{(\ell)}).
        \]

        \item Nichtlineare Aktivierungsfunktionen sind notwendig, weil mehrere
        lineare Schichten zusammen wieder nur eine lineare Funktion ergeben.

        \item Lernen in MLPs erfolgt durch Gradient Descent:
        \[
            \theta
            \leftarrow
            \theta
            -
            \eta\nabla_\theta E(\theta).
        \]

        \item Backpropagation berechnet die Gradienten effizient mit der
        Kettenregel.

        \item Der Fehlerterm einer Schicht ist:
        \[
            \delta^{(\ell)}
            =
            \frac{\partial E}{\partial a^{(\ell)}}.
        \]

        \item Für Gewichtsmatrizen gilt:
        \[
            \frac{\partial E}{\partial W^{(\ell)}}
            =
            \delta^{(\ell)}
            \left(h^{(\ell-1)}\right)^T.
        \]
    \end{itemize}

    Die Kernaussage lautet:

    \[
        \boxed{
            \text{MLPs erweitern Perzeptrons durch versteckte Schichten und
            Nichtlinearitäten, sodass komplexe Funktionen gelernt werden können.}
        }
    \]
\end{answerbox}

%%%%%%%%%%%%%%%%%%%%%%%%%%%%%%%%%%%%%%%%%%%%%%%%%%%%%%%%%%%%%%%%%%%
\section{Lecture 6: Neural Networks -- DNNs and Their Capabilities}
%%%%%%%%%%%%%%%%%%%%%%%%%%%%%%%%%%%%%%%%%%%%%%%%%%%%%%%%%%%%%%%%%%%

\subsection{Overview}

\begin{notebox}
    Lecture 6 erweitert die Diskussion zu Multi-Layer Perceptrons und betrachtet,
    warum tiefe neuronale Netze leistungsfähig sind.

    Die zentralen Themen sind:

    \begin{itemize}
        \item Verallgemeinerung der Backpropagation-Herleitung,
        \item Eigenschaften von Multi-Layer Perceptrons,
        \item MLPs und das XOR-Problem,
        \item Overfitting und Generalization Error bei MLPs,
        \item Validation Sets für MLPs,
        \item MLPs und Maximum Likelihood Estimation,
        \item Generalisierung und Validierung,
        \item andere neuronale Netzwerkarchitekturen.
    \end{itemize}
\end{notebox}

\begin{answerbox}
    Die Kernaussage von Lecture 6 ist:

    \[
        \boxed{
            \text{Tiefe neuronale Netze lernen komplexe Funktionen durch
            Komposition vieler einfacher nichtlinearer Transformationen.}
        }
    \]

    Ihre Leistungsfähigkeit hängt aber nicht nur von der Modellklasse ab,
    sondern auch von Training, Regularisierung, Datenmenge und Validierung.
\end{answerbox}

%%%%%%%%%%%%%%%%%%%%%%%%%%%%%%%%%%%%%%%%%%%%%%%%%%%%%%%%%%%%%%%%%%%
\subsection{Part C.1: Generalization of Backpropagation}
%%%%%%%%%%%%%%%%%%%%%%%%%%%%%%%%%%%%%%%%%%%%%%%%%%%%%%%%%%%%%%%%%%%

\begin{notebox}
    Backpropagation ist eine systematische Anwendung der Kettenregel auf
    mehrschichtige neuronale Netze.

    Ein MLP besteht aus mehreren Schichten der Form
    \[
        a^{(\ell)}
        =
        W^{(\ell)}h^{(\ell-1)} + b^{(\ell)}
    \]
    und
    \[
        h^{(\ell)}
        =
        g^{(\ell)}(a^{(\ell)}).
    \]

    Dabei ist:
    \[
        h^{(0)} = x
    \]
    die Eingabe des Netzes.
\end{notebox}

\begin{solutionbox}
    Der Forward Pass berechnet Schicht für Schicht:

    \[
        x = h^{(0)}
        \longrightarrow
        h^{(1)}
        \longrightarrow
        h^{(2)}
        \longrightarrow
        \dots
        \longrightarrow
        h^{(L)}.
    \]

    Die letzte Ausgabe \(h^{(L)}\) wird mit dem Zielwert \(t\) verglichen.
    Daraus entsteht eine Fehlerfunktion:
    \[
        E = E(h^{(L)},t).
    \]

    Beim Backward Pass werden dann die Ableitungen dieser Fehlerfunktion nach
    allen Parametern berechnet.
\end{solutionbox}

\begin{answerbox}
    Ein MLP ist eine Komposition von Funktionen:

    \[
        \boxed{
            h^{(L)}
            =
            f^{(L)} \circ f^{(L-1)} \circ \dots \circ f^{(1)}(x)
        }
    \]

    Backpropagation berechnet effizient:

    \[
        \boxed{
            \nabla_\theta E
        }
    \]

    für alle Parameter \(\theta\) des Netzwerks.
\end{answerbox}

%%%%%%%%%%%%%%%%%%%%%%%%%%%%%%%%%%%%%%%%%%%%%%%%%%%%%%%%%%%%%%%%%%%
\subsection{Generalization of the Backpropagation Derivation}
%%%%%%%%%%%%%%%%%%%%%%%%%%%%%%%%%%%%%%%%%%%%%%%%%%%%%%%%%%%%%%%%%%%

\begin{notebox}
    Die Backpropagation-Herleitung aus Lecture 5 lässt sich auf beliebig viele
    Schichten verallgemeinern.

    Die zentrale Größe ist der Fehlerterm einer Schicht:
    \[
        \delta^{(\ell)}
        =
        \frac{\partial E}{\partial a^{(\ell)}}.
    \]

    Dieser Fehlerterm beschreibt, wie stark eine Änderung der Voraktivierung
    \(a^{(\ell)}\) den Gesamtfehler beeinflusst.
\end{notebox}

\begin{calculationbox}
    Für die letzte Schicht \(L\) gilt allgemein:
    \[
        \delta^{(L)}
        =
        \frac{\partial E}{\partial a^{(L)}}.
    \]

    Falls
    \[
        h^{(L)} = g^{(L)}(a^{(L)}),
    \]
    dann erhält man mit der Kettenregel:
    \[
        \delta^{(L)}
        =
        \frac{\partial E}{\partial h^{(L)}}
        \odot
        g'^{(L)}(a^{(L)}).
    \]

    Dabei bezeichnet \(\odot\) komponentenweise Multiplikation.
\end{calculationbox}

\begin{calculationbox}
    Für eine versteckte Schicht \(\ell\) hängt der Fehler indirekt über die
    nächste Schicht von \(a^{(\ell)}\) ab.

    Die Rekursion lautet:
    \[
        \delta^{(\ell)}
        =
        \left(
            (W^{(\ell+1)})^T
            \delta^{(\ell+1)}
        \right)
        \odot
        g'^{(\ell)}(a^{(\ell)}).
    \]

    Damit können die Fehlerterme rückwärts berechnet werden:
    \[
        \delta^{(L)}
        \longrightarrow
        \delta^{(L-1)}
        \longrightarrow
        \dots
        \longrightarrow
        \delta^{(1)}.
    \]
\end{calculationbox}

\begin{calculationbox}
    Sind die Fehlerterme \(\delta^{(\ell)}\) bekannt, erhält man die Gradienten:

    \[
        \frac{\partial E}{\partial W^{(\ell)}}
        =
        \delta^{(\ell)}
        \left(h^{(\ell-1)}\right)^T.
    \]

    Für den Bias gilt:
    \[
        \frac{\partial E}{\partial b^{(\ell)}}
        =
        \delta^{(\ell)}.
    \]

    Die Parameter werden dann mit Gradient Descent aktualisiert:

    \[
        W^{(\ell)}
        \leftarrow
        W^{(\ell)}
        -
        \eta
        \frac{\partial E}{\partial W^{(\ell)}}
    \]

    und

    \[
        b^{(\ell)}
        \leftarrow
        b^{(\ell)}
        -
        \eta
        \frac{\partial E}{\partial b^{(\ell)}}.
    \]
\end{calculationbox}

\begin{answerbox}
    Die allgemeine Backpropagation-Rekursion lautet:

    \[
        \boxed{
            \delta^{(\ell)}
            =
            \left(
                (W^{(\ell+1)})^T
                \delta^{(\ell+1)}
            \right)
            \odot
            g'^{(\ell)}(a^{(\ell)})
        }
    \]

    und die Gradienten sind:

    \[
        \boxed{
            \frac{\partial E}{\partial W^{(\ell)}}
            =
            \delta^{(\ell)}
            \left(h^{(\ell-1)}\right)^T
        }
    \]

    \[
        \boxed{
            \frac{\partial E}{\partial b^{(\ell)}}
            =
            \delta^{(\ell)}
        }
    \]
\end{answerbox}

%%%%%%%%%%%%%%%%%%%%%%%%%%%%%%%%%%%%%%%%%%%%%%%%%%%%%%%%%%%%%%%%%%%
\subsection{Properties of Multi-Layer Perceptrons}
%%%%%%%%%%%%%%%%%%%%%%%%%%%%%%%%%%%%%%%%%%%%%%%%%%%%%%%%%%%%%%%%%%%

\begin{notebox}
    Multi-Layer Perceptrons haben gegenüber einzelnen Perzeptrons deutlich
    stärkere Ausdrucksfähigkeit.

    Die wichtigsten Eigenschaften sind:

    \begin{itemize}
        \item mehrere Schichten,
        \item nichtlineare Aktivierungsfunktionen,
        \item lernbare Gewichtsmatrizen und Biases,
        \item hierarchische Repräsentationen,
        \item flexible Entscheidungsgrenzen.
    \end{itemize}
\end{notebox}

\begin{solutionbox}
    Ein einzelnes Perzeptron kann nur lineare Entscheidungsgrenzen darstellen:
    \[
        w^Tx+b=0.
    \]

    Ein MLP kann dagegen durch versteckte Schichten eine Eingabe zuerst in neue
    Merkmalsräume transformieren.

    Dadurch kann ein ursprünglich nicht linear separierbares Problem in einer
    versteckten Repräsentation linear separierbar werden.
\end{solutionbox}

\begin{answerbox}
    Die zentrale Eigenschaft von MLPs ist:

    \[
        \boxed{
            \text{MLPs lernen nicht nur eine Entscheidung, sondern auch eine
            geeignete Repräsentation der Eingabe.}
        }
    \]

    Dadurch können sie komplexere Funktionen modellieren als einfache lineare
    Klassifikatoren.
\end{answerbox}

%%%%%%%%%%%%%%%%%%%%%%%%%%%%%%%%%%%%%%%%%%%%%%%%%%%%%%%%%%%%%%%%%%%
\subsection{MLPs and the XOR Problem}
%%%%%%%%%%%%%%%%%%%%%%%%%%%%%%%%%%%%%%%%%%%%%%%%%%%%%%%%%%%%%%%%%%%

\begin{notebox}
    Das XOR-Problem ist das klassische Beispiel dafür, warum mehrere Schichten
    und Nichtlinearitäten nützlich sind.

    Die XOR-Funktion ist:
    \[
        0 \oplus 0 = 0,
        \qquad
        0 \oplus 1 = 1,
    \]
    \[
        1 \oplus 0 = 1,
        \qquad
        1 \oplus 1 = 0.
    \]

    Die beiden positiven Punkte liegen diagonal gegenüber. Daher ist XOR nicht
    linear separierbar.
\end{notebox}

\begin{solutionbox}
    Ein einzelnes Perzeptron kann XOR nicht lösen, weil es nur eine Gerade als
    Entscheidungsgrenze erzeugen kann.

    Ein MLP mit versteckter Schicht kann dagegen mehrere lineare Schnitte
    kombinieren und danach nichtlinear transformieren.

    Dadurch kann es eine nichtlineare Entscheidungsregion erzeugen, die XOR
    korrekt klassifiziert.
\end{solutionbox}

\begin{answerbox}
    XOR zeigt den Unterschied zwischen Perzeptron und MLP:

    \[
        \boxed{
            \text{Perzeptron: nur linear separierbare Probleme}
        }
    \]

    \[
        \boxed{
            \text{MLP: nichtlineare Entscheidungsgrenzen durch versteckte Schichten}
        }
    \]
\end{answerbox}

%%%%%%%%%%%%%%%%%%%%%%%%%%%%%%%%%%%%%%%%%%%%%%%%%%%%%%%%%%%%%%%%%%%
\subsection{Overfitting and Generalization Error for MLPs}
%%%%%%%%%%%%%%%%%%%%%%%%%%%%%%%%%%%%%%%%%%%%%%%%%%%%%%%%%%%%%%%%%%%

\begin{notebox}
    MLPs sind sehr flexible Modelle. Diese Flexibilität ist einerseits nützlich,
    kann aber auch zu Overfitting führen.

    Overfitting bedeutet:

    \[
        \text{Das Modell passt die Trainingsdaten sehr gut an, generalisiert aber schlecht.}
    \]

    Typisches Muster:
    \[
        E_{\mathrm{train}} \text{ klein},
        \qquad
        E_{\mathrm{test}} \text{ groß}.
    \]
\end{notebox}

\begin{solutionbox}
    Ein MLP kann viele Parameter haben. Wenn das Modell im Vergleich zur
    Datenmenge zu komplex ist, kann es nicht nur die zugrundeliegende Struktur,
    sondern auch das Rauschen der Trainingsdaten lernen.

    Dann sinkt der Trainingsfehler weiter, während der Fehler auf neuen Daten
    steigt.

    Das Ziel ist daher nicht nur, den Trainingsfehler zu minimieren, sondern
    eine gute Generalisierung zu erreichen.
\end{solutionbox}

\begin{calculationbox}
    Der Generalization Error beschreibt den erwarteten Fehler auf neuen,
    unabhängigen Daten:

    \[
        E_{\mathrm{gen}}
        =
        \mathbb{E}_{(x,t)\sim p_{\mathrm{data}}}
        \left[
            L(f_\theta(x),t)
        \right].
    \]

    Dieser Fehler ist in der Praxis unbekannt, da die wahre Datenverteilung
    \(p_{\mathrm{data}}\) unbekannt ist.

    Deshalb approximiert man ihn durch einen Test- oder Validierungsfehler:
    \[
        E_{\mathrm{val}}
        =
        \frac{1}{N_{\mathrm{val}}}
        \sum_{n=1}^{N_{\mathrm{val}}}
        L(f_\theta(x_n),t_n).
    \]
\end{calculationbox}

\begin{answerbox}
    Overfitting erkennt man typischerweise daran:

    \[
        \boxed{
            E_{\mathrm{train}} \downarrow
            \quad \text{aber} \quad
            E_{\mathrm{val}} \uparrow
        }
    \]

    Gute Modelle sollen nicht nur auf Trainingsdaten gut sein, sondern auf neuen
    Daten:

    \[
        \boxed{
            \text{Ziel: kleiner Generalization Error.}
        }
    \]
\end{answerbox}

%%%%%%%%%%%%%%%%%%%%%%%%%%%%%%%%%%%%%%%%%%%%%%%%%%%%%%%%%%%%%%%%%%%
\subsection{Validation Sets for MLPs}
%%%%%%%%%%%%%%%%%%%%%%%%%%%%%%%%%%%%%%%%%%%%%%%%%%%%%%%%%%%%%%%%%%%

\begin{notebox}
    Ein Validation Set ist ein Teil der Daten, der nicht zum direkten Training
    der Gewichte verwendet wird.

    Es dient dazu, Modellentscheidungen zu treffen, zum Beispiel:

    \begin{itemize}
        \item Wahl der Netzwerkgröße,
        \item Wahl der Lernrate,
        \item Wahl der Regularisierung,
        \item Entscheidung über Early Stopping,
        \item Vergleich verschiedener Modelle.
    \end{itemize}
\end{notebox}

\begin{solutionbox}
    Typische Aufteilung der Daten:

    \[
        \text{Training Set}
        \qquad
        \text{Validation Set}
        \qquad
        \text{Test Set}.
    \]

    Das Training Set wird verwendet, um Parameter wie Gewichte und Biases zu
    lernen.

    Das Validation Set wird verwendet, um Hyperparameter und Modellvarianten zu
    wählen.

    Das Test Set wird erst am Ende verwendet, um die finale Generalisierung zu
    bewerten.
\end{solutionbox}

\begin{answerbox}
    Die Rollen der Datensätze sind:

    \[
        \boxed{
            \text{Training Set: Parameter lernen}
        }
    \]

    \[
        \boxed{
            \text{Validation Set: Hyperparameter und Modellwahl}
        }
    \]

    \[
        \boxed{
            \text{Test Set: finale Bewertung}
        }
    \]

    Das Test Set sollte nicht zur Modellwahl verwendet werden.
\end{answerbox}

%%%%%%%%%%%%%%%%%%%%%%%%%%%%%%%%%%%%%%%%%%%%%%%%%%%%%%%%%%%%%%%%%%%
\subsection{Part C.2: MLPs and Maximum Likelihood Estimation}
%%%%%%%%%%%%%%%%%%%%%%%%%%%%%%%%%%%%%%%%%%%%%%%%%%%%%%%%%%%%%%%%%%%

\begin{notebox}
    MLPs können probabilistisch interpretiert werden.

    Dazu modelliert das Netzwerk Parameter einer Wahrscheinlichkeitsverteilung.

    Für Klassifikation gibt ein MLP zum Beispiel Klassenwahrscheinlichkeiten aus:
    \[
        p(t=k \mid x,\theta).
    \]

    Für Regression kann ein MLP den Mittelwert einer Gauß-Verteilung modellieren:
    \[
        p(t \mid x,\theta)
        =
        \mathcal{N}(t \mid f_\theta(x),\sigma^2).
    \]
\end{notebox}

\begin{calculationbox}
    Für gelabelte Daten
    \[
        \mathcal{D}
        =
        \set{(x_1,t_1),\dots,(x_N,t_N)}
    \]
    lautet die Likelihood:
    \[
        L(\theta)
        =
        \prod_{n=1}^{N}
        p(t_n \mid x_n,\theta).
    \]

    Die Log-Likelihood ist:
    \[
        \ell(\theta)
        =
        \sum_{n=1}^{N}
        \log p(t_n \mid x_n,\theta).
    \]

    Maximum Likelihood bedeutet:
    \[
        \hat\theta_{\mathrm{ML}}
        =
        \operatorname*{arg\,max}_{\theta}
        \ell(\theta).
    \]
\end{calculationbox}

\begin{calculationbox}
    In der Praxis minimiert man oft die negative Log-Likelihood:

    \[
        E(\theta)
        =
        -
        \sum_{n=1}^{N}
        \log p(t_n \mid x_n,\theta).
    \]

    Für binäre Klassifikation mit
    \[
        y_n = p(t_n=1 \mid x_n,\theta)
    \]
    erhält man die Cross-Entropy:
    \[
        E(\theta)
        =
        -
        \sum_{n=1}^{N}
        \left[
            t_n\log y_n
            +
            (1-t_n)\log(1-y_n)
        \right].
    \]
\end{calculationbox}

\begin{answerbox}
    MLP-Training kann als Maximum Likelihood formuliert werden:

    \[
        \boxed{
            \hat\theta_{\mathrm{ML}}
            =
            \operatorname*{arg\,max}_{\theta}
            \sum_{n=1}^{N}
            \log p(t_n \mid x_n,\theta)
        }
    \]

    Äquivalent:

    \[
        \boxed{
            \hat\theta_{\mathrm{ML}}
            =
            \operatorname*{arg\,min}_{\theta}
            \left[
                -
                \sum_{n=1}^{N}
                \log p(t_n \mid x_n,\theta)
            \right]
        }
    \]

    Viele Loss-Funktionen entstehen also aus probabilistischen Annahmen.
\end{answerbox}

%%%%%%%%%%%%%%%%%%%%%%%%%%%%%%%%%%%%%%%%%%%%%%%%%%%%%%%%%%%%%%%%%%%
\subsection{Part C.3: Generalization and Validation}
%%%%%%%%%%%%%%%%%%%%%%%%%%%%%%%%%%%%%%%%%%%%%%%%%%%%%%%%%%%%%%%%%%%

\begin{notebox}
    Generalisierung beschreibt die Fähigkeit eines Modells, auf neuen,
    ungesehenen Daten gut zu funktionieren.

    Ein Modell soll also nicht nur die Trainingsdaten auswendig lernen, sondern
    die zugrundeliegende Struktur der Daten erfassen.
\end{notebox}

\begin{solutionbox}
    Generalisierung hängt von mehreren Faktoren ab:

    \begin{itemize}
        \item Größe und Qualität des Trainingsdatensatzes,
        \item Modellkomplexität,
        \item Regularisierung,
        \item Optimierungsverfahren,
        \item Wahl der Hyperparameter,
        \item Verteilung von Trainings- und Testdaten.
    \end{itemize}

    Ein Validation Set hilft dabei, diese Faktoren praktisch zu kontrollieren.
\end{solutionbox}

\begin{calculationbox}
    Häufig betrachtet man drei Fehlergrößen:

    \[
        E_{\mathrm{train}}
        =
        \frac{1}{N_{\mathrm{train}}}
        \sum_{n=1}^{N_{\mathrm{train}}}
        L(f_\theta(x_n),t_n),
    \]

    \[
        E_{\mathrm{val}}
        =
        \frac{1}{N_{\mathrm{val}}}
        \sum_{n=1}^{N_{\mathrm{val}}}
        L(f_\theta(x_n),t_n),
    \]

    \[
        E_{\mathrm{test}}
        =
        \frac{1}{N_{\mathrm{test}}}
        \sum_{n=1}^{N_{\mathrm{test}}}
        L(f_\theta(x_n),t_n).
    \]

    Der Trainingsfehler misst Anpassung an die Trainingsdaten.

    Der Validierungsfehler hilft bei Modellwahl.

    Der Testfehler schätzt die finale Generalisierung.
\end{calculationbox}

\begin{answerbox}
    Gute Modellwahl bedeutet:

    \[
        \boxed{
            \text{Nicht das Modell mit minimalem Trainingsfehler wählen,
            sondern das mit gutem Validierungsfehler.}
        }
    \]

    Ein Validation Set dient zur Modell- und Hyperparameterwahl.

    Ein Test Set dient zur finalen unabhängigen Bewertung.
\end{answerbox}

%%%%%%%%%%%%%%%%%%%%%%%%%%%%%%%%%%%%%%%%%%%%%%%%%%%%%%%%%%%%%%%%%%%
\subsection{Regularization and Early Stopping}
%%%%%%%%%%%%%%%%%%%%%%%%%%%%%%%%%%%%%%%%%%%%%%%%%%%%%%%%%%%%%%%%%%%

\begin{notebox}
    Regularisierung bezeichnet Methoden, die Overfitting reduzieren sollen.

    Sie schränken das Modell direkt oder indirekt ein, damit es besser
    generalisiert.
\end{notebox}

\begin{solutionbox}
    Typische Regularisierungsmethoden für MLPs sind:

    \begin{itemize}
        \item Gewichtsnorm-Strafen wie \(L^2\)-Regularisierung,
        \item kleinere Netzwerke,
        \item Dropout,
        \item Data Augmentation,
        \item Early Stopping.
    \end{itemize}

    Bei Early Stopping beobachtet man den Validierungsfehler während des
    Trainings.

    Wenn der Trainingsfehler weiter sinkt, aber der Validierungsfehler wieder
    steigt, stoppt man das Training.
\end{solutionbox}

\begin{calculationbox}
    Ein typischer regularisierter Loss hat die Form:

    \[
        E_{\mathrm{reg}}(\theta)
        =
        E_{\mathrm{data}}(\theta)
        +
        \lambda \Omega(\theta).
    \]

    Für \(L^2\)-Regularisierung:

    \[
        \Omega(\theta)
        =
        \frac{1}{2}
        \|\theta\|^2.
    \]

    Also:
    \[
        E_{\mathrm{reg}}(\theta)
        =
        E_{\mathrm{data}}(\theta)
        +
        \frac{\lambda}{2}
        \|\theta\|^2.
    \]
\end{calculationbox}

\begin{answerbox}
    Regularisierung soll Overfitting reduzieren:

    \[
        \boxed{
            E_{\mathrm{reg}}(\theta)
            =
            E_{\mathrm{data}}(\theta)
            +
            \lambda \Omega(\theta)
        }
    \]

    Early Stopping nutzt den Validierungsfehler, um das Training rechtzeitig zu
    beenden.
\end{answerbox}

%%%%%%%%%%%%%%%%%%%%%%%%%%%%%%%%%%%%%%%%%%%%%%%%%%%%%%%%%%%%%%%%%%%
\subsection{Part C.4: Other Neural Networks}
%%%%%%%%%%%%%%%%%%%%%%%%%%%%%%%%%%%%%%%%%%%%%%%%%%%%%%%%%%%%%%%%%%%

\begin{notebox}
    MLPs sind nicht die einzige Art neuronaler Netze.

    Für verschiedene Datenarten und Aufgaben verwendet man spezialisierte
    Architekturen.
\end{notebox}

\begin{solutionbox}
    Beispiele für andere neuronale Netzwerke sind:

    \begin{itemize}
        \item \textbf{Convolutional Neural Networks:}
        besonders wichtig für Bilder und räumliche Daten.

        \item \textbf{Recurrent Neural Networks:}
        historisch wichtig für Sequenzen und Zeitreihen.

        \item \textbf{Long Short-Term Memory Networks:}
        spezielle rekurrente Netze zur Verarbeitung längerer Abhängigkeiten.

        \item \textbf{Autoencoder:}
        Netze, die Daten komprimieren und rekonstruieren lernen.

        \item \textbf{Transformer:}
        Architekturen mit Attention-Mechanismus, besonders wichtig für Sprache
        und viele moderne Deep-Learning-Systeme.
    \end{itemize}
\end{solutionbox}

\begin{answerbox}
    Unterschiedliche neuronale Netzwerkarchitekturen sind für unterschiedliche
    Datenstrukturen geeignet:

    \[
        \boxed{
            \text{MLPs: allgemeine Vektordaten}
        }
    \]

    \[
        \boxed{
            \text{CNNs: Bilder und räumliche Muster}
        }
    \]

    \[
        \boxed{
            \text{RNNs/LSTMs/Transformer: Sequenzen}
        }
    \]

    \[
        \boxed{
            \text{Autoencoder: Repräsentationslernen und Rekonstruktion}
        }
    \]
\end{answerbox}

%%%%%%%%%%%%%%%%%%%%%%%%%%%%%%%%%%%%%%%%%%%%%%%%%%%%%%%%%%%%%%%%%%%
\subsection{Prüfungsrelevante Kurzfassung}
%%%%%%%%%%%%%%%%%%%%%%%%%%%%%%%%%%%%%%%%%%%%%%%%%%%%%%%%%%%%%%%%%%%

\begin{answerbox}
    Lecture 6 behandelt DNNs und ihre Fähigkeiten.

    Die wichtigsten Punkte sind:

    \begin{itemize}
        \item Ein MLP besteht aus Schichten:
        \[
            a^{(\ell)}
            =
            W^{(\ell)}h^{(\ell-1)} + b^{(\ell)},
            \qquad
            h^{(\ell)}
            =
            g^{(\ell)}(a^{(\ell)}).
        \]

        \item Backpropagation berechnet Gradienten effizient mit der Kettenregel.

        \item Der Fehlerterm einer Schicht ist:
        \[
            \delta^{(\ell)}
            =
            \frac{\partial E}{\partial a^{(\ell)}}.
        \]

        \item Für versteckte Schichten gilt:
        \[
            \delta^{(\ell)}
            =
            \left(
                (W^{(\ell+1)})^T
                \delta^{(\ell+1)}
            \right)
            \odot
            g'^{(\ell)}(a^{(\ell)}).
        \]

        \item Gewichtgradient:
        \[
            \frac{\partial E}{\partial W^{(\ell)}}
            =
            \delta^{(\ell)}
            \left(h^{(\ell-1)}\right)^T.
        \]

        \item MLPs können XOR lösen, weil sie nichtlineare versteckte
        Repräsentationen lernen.

        \item Overfitting bedeutet:
        \[
            E_{\mathrm{train}} \text{ klein},
            \qquad
            E_{\mathrm{val/test}} \text{ groß}.
        \]

        \item Validation Sets dienen zur Modellwahl und Hyperparameterwahl.

        \item MLP-Training kann als Maximum Likelihood formuliert werden:
        \[
            \hat\theta_{\mathrm{ML}}
            =
            \operatorname*{arg\,max}_{\theta}
            \sum_{n=1}^{N}
            \log p(t_n \mid x_n,\theta).
        \]

        \item Negative Log-Likelihood wird als Loss minimiert:
        \[
            E(\theta)
            =
            -
            \sum_{n=1}^{N}
            \log p(t_n \mid x_n,\theta).
        \]

        \item Regularisierung reduziert Overfitting:
        \[
            E_{\mathrm{reg}}(\theta)
            =
            E_{\mathrm{data}}(\theta)
            +
            \lambda \Omega(\theta).
        \]

        \item Andere Netztypen sind zum Beispiel CNNs, RNNs, LSTMs,
        Autoencoder und Transformer.
    \end{itemize}

    Die Kernaussage lautet:

    \[
        \boxed{
            \text{DNNs sind leistungsfähig, weil sie durch viele Schichten
            hierarchische nichtlineare Repräsentationen lernen.}
        }
    \]
\end{answerbox}

%%%%%%%%%%%%%%%%%%%%%%%%%%%%%%%%%%%%%%%%%%%%%%%%%%%%%%%%%%%%%%%%%%%
\section{Lecture 7: Dimensionality Reduction with PCA}
%%%%%%%%%%%%%%%%%%%%%%%%%%%%%%%%%%%%%%%%%%%%%%%%%%%%%%%%%%%%%%%%%%%

\subsection{Overview}

\begin{notebox}
    Lecture 7 behandelt Dimensionsreduktion mit Principal Component Analysis,
    kurz PCA.

    Die zentrale Idee ist:

    \[
        \text{Hochdimensionale Daten sollen durch wenige informative Richtungen beschrieben werden.}
    \]

    Die wichtigsten Themen sind:

    \begin{itemize}
        \item Projektionsoperatoren,
        \item das eindimensionale PCA-Problem,
        \item Maximum-Variance-Ansatz,
        \item Herleitung der Varianzformel,
        \item optimale Lösung über Eigenvektoren,
        \item das allgemeine niedrigdimensionale Hyperplane-Problem,
        \item Anwendungen von PCA,
        \item Diskussion der Stärken und Grenzen von PCA.
    \end{itemize}
\end{notebox}

\begin{answerbox}
    Die Kernaussage von Lecture 7 ist:

    \[
        \boxed{
            \text{PCA findet Richtungen maximaler Varianz und projiziert Daten auf diese Richtungen.}
        }
    \]

    Dadurch kann man hochdimensionale Daten visualisieren, interpretieren oder
    komprimieren.
\end{answerbox}

%%%%%%%%%%%%%%%%%%%%%%%%%%%%%%%%%%%%%%%%%%%%%%%%%%%%%%%%%%%%%%%%%%%
\subsection{Preparation: Projection Operators}
%%%%%%%%%%%%%%%%%%%%%%%%%%%%%%%%%%%%%%%%%%%%%%%%%%%%%%%%%%%%%%%%%%%

\begin{notebox}
    Eine Projektion bildet einen Vektor auf einen Unterraum ab.

    Für PCA interessiert uns besonders die Projektion eines Datenpunkts \(x\)
    auf eine Richtung \(u\).

    Dabei nehmen wir an, dass \(u\) normiert ist:

    \[
        \|u\| = 1.
    \]
\end{notebox}

\begin{calculationbox}
    Die Projektion eines Vektors \(x \in \mathbb{R}^D\) auf die Richtung \(u\)
    ist:

    \[
        \operatorname{proj}_u(x)
        =
        (u^T x)u.
    \]

    Der Skalar

    \[
        z = u^T x
    \]

    ist die Koordinate von \(x\) entlang der Richtung \(u\).

    Die Matrixform des Projektionsoperators ist:

    \[
        P_u
        =
        uu^T.
    \]

    Damit:

    \[
        \operatorname{proj}_u(x)
        =
        P_u x
        =
        uu^T x.
    \]
\end{calculationbox}

\begin{solutionbox}
    Wenn die Daten nicht um den Ursprung zentriert sind, projiziert man
    normalerweise die zentrierten Daten.

    Für Datenpunkte

    \[
        x_1,\dots,x_N
    \]

    berechnet man zuerst den Mittelwert:

    \[
        \bar x
        =
        \frac{1}{N}
        \sum_{n=1}^{N} x_n.
    \]

    Danach betrachtet man:

    \[
        \tilde x_n
        =
        x_n - \bar x.
    \]

    PCA wird typischerweise auf diesen zentrierten Daten durchgeführt.
\end{solutionbox}

\begin{answerbox}
    Projektion auf eine normierte Richtung \(u\):

    \[
        \boxed{
            \operatorname{proj}_u(x)
            =
            (u^Tx)u
        }
    \]

    Projektionsmatrix:

    \[
        \boxed{
            P_u = uu^T
        }
    \]

    Zentrierte Daten:

    \[
        \boxed{
            \tilde x_n = x_n-\bar x
        }
    \]
\end{answerbox}

%%%%%%%%%%%%%%%%%%%%%%%%%%%%%%%%%%%%%%%%%%%%%%%%%%%%%%%%%%%%%%%%%%%
\subsection{The One-Dimensional Problem: Finding a Line Through Data}
%%%%%%%%%%%%%%%%%%%%%%%%%%%%%%%%%%%%%%%%%%%%%%%%%%%%%%%%%%%%%%%%%%%

\begin{notebox}
    Im eindimensionalen PCA-Problem sucht man eine Gerade, auf die die Daten
    projiziert werden sollen.

    Die Gerade soll möglichst viel Information der Daten erhalten.

    PCA formalisiert diese Idee durch den Maximum-Variance-Ansatz:

    \[
        \text{Wähle die Richtung, in der die projizierten Daten maximale Varianz haben.}
    \]
\end{notebox}

\begin{solutionbox}
    Gegeben seien zentrierte Datenpunkte:

    \[
        \tilde x_1,\dots,\tilde x_N \in \mathbb{R}^D.
    \]

    Wir suchen eine normierte Richtung

    \[
        u \in \mathbb{R}^D,
        \qquad
        \|u\|=1.
    \]

    Jeder Datenpunkt wird auf diese Richtung projiziert. Die eindimensionale
    Koordinate des Datenpunkts entlang \(u\) ist:

    \[
        z_n
        =
        u^T \tilde x_n.
    \]

    PCA sucht die Richtung \(u\), sodass die Varianz der \(z_n\) maximal wird.
\end{solutionbox}

\begin{answerbox}
    Das eindimensionale PCA-Problem lautet:

    \[
        \boxed{
            \text{Finde eine Richtung } u
            \text{ mit } \|u\|=1,
            \text{ sodass die Projektionen } u^T\tilde x_n
            \text{ maximale Varianz haben.}
        }
    \]
\end{answerbox}

%%%%%%%%%%%%%%%%%%%%%%%%%%%%%%%%%%%%%%%%%%%%%%%%%%%%%%%%%%%%%%%%%%%
\subsection{Maximum Variance Approach in One Dimension}
%%%%%%%%%%%%%%%%%%%%%%%%%%%%%%%%%%%%%%%%%%%%%%%%%%%%%%%%%%%%%%%%%%%

\begin{notebox}
    Die Varianz der projizierten Daten misst, wie stark sich die Daten entlang
    der Richtung \(u\) ausbreiten.

    Eine Richtung mit großer Varianz enthält viel Struktur der Daten.

    Eine Richtung mit kleiner Varianz enthält wenig Variation und ist für die
    Beschreibung der Daten weniger informativ.
\end{notebox}

\begin{calculationbox}
    Die Projektionskoordinate des \(n\)-ten Datenpunkts ist:

    \[
        z_n
        =
        u^T\tilde x_n.
    \]

    Da die Daten zentriert sind, gilt:

    \[
        \frac{1}{N}
        \sum_{n=1}^{N}
        \tilde x_n
        =
        0.
    \]

    Daher ist auch der Mittelwert der Projektionen null:

    \[
        \bar z
        =
        \frac{1}{N}
        \sum_{n=1}^{N}
        z_n
        =
        \frac{1}{N}
        \sum_{n=1}^{N}
        u^T\tilde x_n
        =
        u^T
        \left(
            \frac{1}{N}
            \sum_{n=1}^{N}
            \tilde x_n
        \right)
        =
        0.
    \]
\end{calculationbox}

\begin{calculationbox}
    Die Varianz der Projektionen ist daher:

    \[
        \operatorname{Var}(z)
        =
        \frac{1}{N}
        \sum_{n=1}^{N}
        z_n^2.
    \]

    Einsetzen von

    \[
        z_n = u^T\tilde x_n
    \]

    liefert:

    \[
        \operatorname{Var}(z)
        =
        \frac{1}{N}
        \sum_{n=1}^{N}
        (u^T\tilde x_n)^2.
    \]

    Da

    \[
        (u^T\tilde x_n)^2
        =
        u^T\tilde x_n \tilde x_n^T u,
    \]

    folgt:

    \[
        \operatorname{Var}(z)
        =
        u^T
        \left(
            \frac{1}{N}
            \sum_{n=1}^{N}
            \tilde x_n \tilde x_n^T
        \right)
        u.
    \]
\end{calculationbox}

\begin{calculationbox}
    Die Matrix

    \[
        S
        =
        \frac{1}{N}
        \sum_{n=1}^{N}
        \tilde x_n \tilde x_n^T
    \]

    ist die empirische Kovarianzmatrix der Daten.

    Damit lautet die Varianz entlang der Richtung \(u\):

    \[
        \operatorname{Var}(z)
        =
        u^T S u.
    \]

    Das Optimierungsproblem ist:

    \[
        \max_u
        \quad
        u^T S u
        \qquad
        \text{unter der Nebenbedingung}
        \qquad
        u^Tu=1.
    \]
\end{calculationbox}

\begin{answerbox}
    Die Varianz der Projektionen auf \(u\) ist:

    \[
        \boxed{
            \operatorname{Var}(z)
            =
            u^T S u
        }
    \]

    mit Kovarianzmatrix:

    \[
        \boxed{
            S
            =
            \frac{1}{N}
            \sum_{n=1}^{N}
            \tilde x_n \tilde x_n^T
        }
    \]

    Das PCA-Optimierungsproblem ist:

    \[
        \boxed{
            \max_{u}
            u^T S u
            \quad
            \text{s.t.}
            \quad
            u^T u = 1.
        }
    \]
\end{answerbox}

%%%%%%%%%%%%%%%%%%%%%%%%%%%%%%%%%%%%%%%%%%%%%%%%%%%%%%%%%%%%%%%%%%%
\subsection{Optimal Solution in Terms of Eigenvectors}
%%%%%%%%%%%%%%%%%%%%%%%%%%%%%%%%%%%%%%%%%%%%%%%%%%%%%%%%%%%%%%%%%%%

\begin{notebox}
    Das PCA-Optimierungsproblem führt auf ein Eigenwertproblem.

    Die optimale Richtung ist der Eigenvektor der Kovarianzmatrix zum größten
    Eigenwert.
\end{notebox}

\begin{calculationbox}
    Wir maximieren:

    \[
        u^T S u
    \]

    unter der Nebenbedingung:

    \[
        u^T u = 1.
    \]

    Dafür verwenden wir einen Lagrange-Multiplikator \(\lambda\):

    \[
        \mathcal{L}(u,\lambda)
        =
        u^T S u
        -
        \lambda(u^Tu-1).
    \]

    Ableitung nach \(u\):

    \[
        \frac{\partial \mathcal{L}}{\partial u}
        =
        2Su - 2\lambda u.
    \]

    Setze gleich null:

    \[
        2Su - 2\lambda u = 0.
    \]

    Also:

    \[
        Su = \lambda u.
    \]
\end{calculationbox}

\begin{solutionbox}
    Die Gleichung

    \[
        Su = \lambda u
    \]

    ist die Eigenwertgleichung der Kovarianzmatrix \(S\).

    Für einen normierten Eigenvektor \(u\) gilt:

    \[
        u^T S u
        =
        u^T \lambda u
        =
        \lambda u^T u
        =
        \lambda.
    \]

    Die Varianz entlang einer Eigenrichtung ist also genau der zugehörige
    Eigenwert.

    Daher maximiert der Eigenvektor mit dem größten Eigenwert die Varianz.
\end{solutionbox}

\begin{answerbox}
    Die erste Hauptkomponente ist:

    \[
        \boxed{
            u_1
            =
            \text{Eigenvektor von } S \text{ zum größten Eigenwert } \lambda_1.
        }
    \]

    Die erklärte Varianz entlang dieser Richtung ist:

    \[
        \boxed{
            \operatorname{Var}(u_1^T\tilde x)
            =
            \lambda_1.
        }
    \]
\end{answerbox}

%%%%%%%%%%%%%%%%%%%%%%%%%%%%%%%%%%%%%%%%%%%%%%%%%%%%%%%%%%%%%%%%%%%
\subsection{The General Problem: Finding a Low-Dimensional Hyperplane}
%%%%%%%%%%%%%%%%%%%%%%%%%%%%%%%%%%%%%%%%%%%%%%%%%%%%%%%%%%%%%%%%%%%

\begin{notebox}
    Im allgemeinen PCA-Problem sucht man nicht nur eine Richtung, sondern einen
    niedrigdimensionalen Unterraum.

    Ziel:

    \[
        \text{Finde einen } M\text{-dimensionalen Unterraum in } \mathbb{R}^D,
        \quad M < D,
    \]

    der möglichst viel Varianz der Daten erhält.
\end{notebox}

\begin{solutionbox}
    Statt einer Richtung \(u\) betrachtet man mehrere orthonormale Richtungen:

    \[
        u_1,\dots,u_M.
    \]

    Diese bilden die Spalten einer Matrix:

    \[
        U_M
        =
        \begin{pmatrix}
        | & & | \\
        u_1 & \cdots & u_M \\
        | & & |
        \end{pmatrix}
        \in \mathbb{R}^{D \times M}.
    \]

    Die Orthonormalität bedeutet:

    \[
        U_M^T U_M = I_M.
    \]

    Die Projektion eines zentrierten Datenpunkts \(\tilde x\) auf den
    \(M\)-dimensionalen Unterraum ist:

    \[
        z
        =
        U_M^T \tilde x.
    \]

    Dabei gilt:

    \[
        z \in \mathbb{R}^M.
    \]
\end{solutionbox}

\begin{answerbox}
    Allgemeine PCA sucht eine Matrix

    \[
        \boxed{
            U_M = (u_1,\dots,u_M)
        }
    \]

    mit

    \[
        \boxed{
            U_M^T U_M = I_M
        }
    \]

    sodass die Projektionen

    \[
        \boxed{
            z_n = U_M^T \tilde x_n
        }
    \]

    möglichst viel Varianz erhalten.
\end{answerbox}

%%%%%%%%%%%%%%%%%%%%%%%%%%%%%%%%%%%%%%%%%%%%%%%%%%%%%%%%%%%%%%%%%%%
\subsection{Maximum Variance Approach in the General Case}
%%%%%%%%%%%%%%%%%%%%%%%%%%%%%%%%%%%%%%%%%%%%%%%%%%%%%%%%%%%%%%%%%%%

\begin{calculationbox}
    Die Gesamtvarianz der Projektionen auf den Unterraum ist die Summe der
    Varianzen entlang der gewählten Richtungen:

    \[
        J(U_M)
        =
        \sum_{j=1}^{M}
        u_j^T S u_j.
    \]

    In Matrixform kann man das schreiben als:

    \[
        J(U_M)
        =
        \operatorname{tr}(U_M^T S U_M).
    \]

    Das Optimierungsproblem lautet:

    \[
        \max_{U_M}
        \operatorname{tr}(U_M^T S U_M)
    \]

    unter der Nebenbedingung:

    \[
        U_M^T U_M = I_M.
    \]
\end{calculationbox}

\begin{solutionbox}
    Die optimale Lösung wird durch die Eigenvektoren der Kovarianzmatrix gegeben.

    Sei:

    \[
        Su_j = \lambda_j u_j
    \]

    mit geordneten Eigenwerten:

    \[
        \lambda_1 \geq \lambda_2 \geq \dots \geq \lambda_D.
    \]

    Dann wählt PCA die ersten \(M\) Eigenvektoren:

    \[
        u_1,\dots,u_M.
    \]

    Diese Richtungen enthalten maximal viel Varianz unter allen
    \(M\)-dimensionalen linearen Projektionen.
\end{solutionbox}

\begin{answerbox}
    Die allgemeine PCA-Lösung ist:

    \[
        \boxed{
            U_M
            =
            (u_1,\dots,u_M)
        }
    \]

    wobei \(u_1,\dots,u_M\) die Eigenvektoren von \(S\) zu den größten
    Eigenwerten sind:

    \[
        \boxed{
            Su_j=\lambda_j u_j,
            \qquad
            \lambda_1 \geq \lambda_2 \geq \dots
        }
    \]

    Die erhaltene Varianz ist:

    \[
        \boxed{
            \sum_{j=1}^{M}\lambda_j.
        }
    \]
\end{answerbox}

%%%%%%%%%%%%%%%%%%%%%%%%%%%%%%%%%%%%%%%%%%%%%%%%%%%%%%%%%%%%%%%%%%%
\subsection{Projection and Reconstruction}
%%%%%%%%%%%%%%%%%%%%%%%%%%%%%%%%%%%%%%%%%%%%%%%%%%%%%%%%%%%%%%%%%%%

\begin{notebox}
    PCA kann auch als Projektion und anschließende Rekonstruktion verstanden
    werden.

    Man reduziert zuerst die Dimension und rekonstruiert danach näherungsweise
    die ursprünglichen Daten.
\end{notebox}

\begin{calculationbox}
    Für einen zentrierten Datenpunkt \(\tilde x\) ist die niedrigdimensionale
    Darstellung:

    \[
        z
        =
        U_M^T \tilde x.
    \]

    Die Rekonstruktion im ursprünglichen Raum ist:

    \[
        \hat x
        =
        U_M z + \bar x.
    \]

    Einsetzen von \(z\):

    \[
        \hat x
        =
        U_M U_M^T \tilde x + \bar x.
    \]

    Für den ursprünglichen Datenpunkt \(x\) gilt:

    \[
        \tilde x = x-\bar x.
    \]

    Also:

    \[
        \hat x
        =
        U_M U_M^T (x-\bar x) + \bar x.
    \]
\end{calculationbox}

\begin{answerbox}
    PCA-Projektion:

    \[
        \boxed{
            z = U_M^T(x-\bar x)
        }
    \]

    PCA-Rekonstruktion:

    \[
        \boxed{
            \hat x
            =
            U_M z + \bar x
            =
            U_MU_M^T(x-\bar x)+\bar x
        }
    \]
\end{answerbox}

%%%%%%%%%%%%%%%%%%%%%%%%%%%%%%%%%%%%%%%%%%%%%%%%%%%%%%%%%%%%%%%%%%%
\subsection{Explained Variance}
%%%%%%%%%%%%%%%%%%%%%%%%%%%%%%%%%%%%%%%%%%%%%%%%%%%%%%%%%%%%%%%%%%%

\begin{notebox}
    Die Eigenwerte der Kovarianzmatrix geben an, wie viel Varianz entlang der
    jeweiligen Hauptkomponenten liegt.

    Große Eigenwerte bedeuten wichtige Richtungen.

    Kleine Eigenwerte bedeuten Richtungen mit wenig Variation.
\end{notebox}

\begin{calculationbox}
    Die gesamte Varianz der Daten ist:

    \[
        \sum_{j=1}^{D}
        \lambda_j.
    \]

    Wenn man nur die ersten \(M\) Hauptkomponenten verwendet, ist die erklärte
    Varianz:

    \[
        \sum_{j=1}^{M}
        \lambda_j.
    \]

    Der Anteil erklärter Varianz ist daher:

    \[
        r_M
        =
        \frac{
            \sum_{j=1}^{M}
            \lambda_j
        }{
            \sum_{j=1}^{D}
            \lambda_j
        }.
    \]
\end{calculationbox}

\begin{answerbox}
    Anteil erklärter Varianz:

    \[
        \boxed{
            r_M
            =
            \frac{
                \sum_{j=1}^{M}
                \lambda_j
            }{
                \sum_{j=1}^{D}
                \lambda_j
            }
        }
    \]

    Diese Größe hilft bei der Wahl der reduzierten Dimension \(M\).
\end{answerbox}

%%%%%%%%%%%%%%%%%%%%%%%%%%%%%%%%%%%%%%%%%%%%%%%%%%%%%%%%%%%%%%%%%%%
\subsection{Example Applications}
%%%%%%%%%%%%%%%%%%%%%%%%%%%%%%%%%%%%%%%%%%%%%%%%%%%%%%%%%%%%%%%%%%%

\begin{notebox}
    PCA wird in vielen Bereichen verwendet, wenn hochdimensionale Daten
    analysiert oder vereinfacht werden sollen.

    Typische Anwendungen sind:

    \begin{itemize}
        \item niedrigdimensionale Visualisierung hochdimensionaler Daten,
        \item Dateninterpretation,
        \item Datenkompression.
    \end{itemize}
\end{notebox}

\begin{solutionbox}
    \textbf{Visualisierung:}

    Hochdimensionale Daten können auf zwei oder drei Hauptkomponenten projiziert
    werden.

    Dadurch kann man sie in 2D oder 3D darstellen:

    \[
        x_n \in \mathbb{R}^D
        \quad \longrightarrow \quad
        z_n \in \mathbb{R}^2
        \text{ oder }
        \mathbb{R}^3.
    \]

    So kann man Cluster, Ausreißer oder grobe Strukturen erkennen.
\end{solutionbox}

\begin{solutionbox}
    \textbf{Dateninterpretation:}

    Die Hauptkomponenten zeigen Richtungen, in denen die Daten besonders stark
    variieren.

    Wenn die ursprünglichen Features interpretierbar sind, kann man untersuchen,
    welche Features stark zu einer Hauptkomponente beitragen.

    Dadurch können wichtige Variationsmuster sichtbar werden.
\end{solutionbox}

\begin{solutionbox}
    \textbf{Datenkompression:}

    Wenn die ersten wenigen Hauptkomponenten bereits einen großen Teil der
    Varianz erklären, kann man die Daten in reduzierter Dimension speichern.

    Statt eines \(D\)-dimensionalen Vektors speichert man nur:

    \[
        z \in \mathbb{R}^M,
        \qquad
        M \ll D.
    \]

    Die Rekonstruktion ist dann näherungsweise möglich durch:

    \[
        \hat x
        =
        U_M z + \bar x.
    \]
\end{solutionbox}

\begin{answerbox}
    PCA-Anwendungen:

    \begin{itemize}
        \item \textbf{Visualisierung:}
        Projektion auf zwei oder drei Dimensionen.

        \item \textbf{Interpretation:}
        Analyse wichtiger Varianzrichtungen.

        \item \textbf{Kompression:}
        Speicherung der Daten mit weniger Koordinaten.
    \end{itemize}

    PCA ersetzt hochdimensionale Daten durch:

    \[
        \boxed{
            z = U_M^T(x-\bar x)
            \in \mathbb{R}^M,
            \qquad
            M \ll D.
        }
    \]
\end{answerbox}

%%%%%%%%%%%%%%%%%%%%%%%%%%%%%%%%%%%%%%%%%%%%%%%%%%%%%%%%%%%%%%%%%%%
\subsection{Discussion of PCA}
%%%%%%%%%%%%%%%%%%%%%%%%%%%%%%%%%%%%%%%%%%%%%%%%%%%%%%%%%%%%%%%%%%%

\begin{notebox}
    PCA ist ein lineares Verfahren zur Dimensionsreduktion.

    Es ist einfach, effizient und gut interpretierbar, hat aber auch Grenzen.
\end{notebox}

\begin{solutionbox}
    Vorteile von PCA:

    \begin{itemize}
        \item rechnerisch gut handhabbar,
        \item klare geometrische Interpretation,
        \item keine Labels nötig,
        \item nützlich für Visualisierung und Kompression,
        \item optimale lineare Projektion im Sinne maximaler Varianz.
    \end{itemize}
\end{solutionbox}

\begin{solutionbox}
    Grenzen von PCA:

    \begin{itemize}
        \item PCA ist linear und findet keine nichtlinearen Strukturen.
        \item Große Varianz muss nicht immer für die Aufgabe relevant sein.
        \item PCA ist empfindlich gegenüber Skalierung der Features.
        \item PCA ist empfindlich gegenüber Ausreißern.
        \item Die Hauptkomponenten können schwer interpretierbar sein.
    \end{itemize}
\end{solutionbox}

\begin{answerbox}
    PCA ist besonders sinnvoll, wenn:

    \[
        \boxed{
            \text{wichtige Struktur der Daten in linearen Varianzrichtungen liegt.}
        }
    \]

    PCA ist weniger geeignet, wenn:

    \[
        \boxed{
            \text{die relevante Struktur stark nichtlinear ist oder nicht mit maximaler Varianz zusammenhängt.}
        }
    \]
\end{answerbox}

%%%%%%%%%%%%%%%%%%%%%%%%%%%%%%%%%%%%%%%%%%%%%%%%%%%%%%%%%%%%%%%%%%%
\subsection{Prüfungsrelevante Kurzfassung}
%%%%%%%%%%%%%%%%%%%%%%%%%%%%%%%%%%%%%%%%%%%%%%%%%%%%%%%%%%%%%%%%%%%

\begin{answerbox}
    Lecture 7 behandelt PCA als lineare Dimensionsreduktion.

    Die wichtigsten Punkte sind:

    \begin{itemize}
        \item Daten werden zuerst zentriert:
        \[
            \tilde x_n = x_n-\bar x.
        \]

        \item Die Kovarianzmatrix ist:
        \[
            S
            =
            \frac{1}{N}
            \sum_{n=1}^{N}
            \tilde x_n\tilde x_n^T.
        \]

        \item Projektion auf eine normierte Richtung \(u\):
        \[
            z_n = u^T\tilde x_n.
        \]

        \item Die Varianz der Projektionen ist:
        \[
            \operatorname{Var}(z)
            =
            u^TSu.
        \]

        \item Das eindimensionale PCA-Problem ist:
        \[
            \max_u u^TSu
            \quad
            \text{s.t.}
            \quad
            u^Tu=1.
        \]

        \item Die Lösung erfüllt:
        \[
            Su=\lambda u.
        \]

        \item Die erste Hauptkomponente ist der Eigenvektor zum größten
        Eigenwert.

        \item Für \(M\)-dimensionale PCA wählt man die ersten \(M\)
        Eigenvektoren:
        \[
            U_M=(u_1,\dots,u_M).
        \]

        \item Projektion:
        \[
            z = U_M^T(x-\bar x).
        \]

        \item Rekonstruktion:
        \[
            \hat x
            =
            U_MU_M^T(x-\bar x)+\bar x.
        \]

        \item Anteil erklärter Varianz:
        \[
            r_M
            =
            \frac{\sum_{j=1}^{M}\lambda_j}
            {\sum_{j=1}^{D}\lambda_j}.
        \]

        \item Anwendungen:
        Visualisierung, Interpretation und Kompression.
    \end{itemize}

    Die Kernaussage lautet:

    \[
        \boxed{
            \text{PCA projiziert Daten auf die Eigenrichtungen der Kovarianzmatrix
            mit größter Varianz.}
        }
    \]
\end{answerbox}

%%%%%%%%%%%%%%%%%%%%%%%%%%%%%%%%%%%%%%%%%%%%%%%%%%%%%%%%%%%%%%%%%%%
\section{Lecture 8: Clustering -- k-means Algorithms, Hierarchical Clustering and GMMs}
%%%%%%%%%%%%%%%%%%%%%%%%%%%%%%%%%%%%%%%%%%%%%%%%%%%%%%%%%%%%%%%%%%%

\subsection{Overview}

\begin{notebox}
    Lecture 8 behandelt Clustering-Verfahren, also Methoden zum Finden von
    Gruppen in ungelabelten Daten.

    Die zentralen Themen sind:

    \begin{itemize}
        \item das Problem, Cluster in Daten zu finden,
        \item k-means,
        \item ein ad-hoc Objective für k-means,
        \item Lloyd's Algorithmus,
        \item Fuzzy-k-means,
        \item hierarchisches Clustering,
        \item agglomeratives Clustering,
        \item Gaussian Mixture Models,
        \item Maximum-Likelihood-Schätzung für GMMs,
        \item Free Energy beziehungsweise ELBO als Hilfsobjective,
        \item Update des Mittelwerts in GMMs,
        \item Diskussion der Methoden.
    \end{itemize}
\end{notebox}

\begin{answerbox}
    Die Kernaussage von Lecture 8 ist:

    \[
        \boxed{
            \text{Clustering versucht, ohne Labels sinnvolle Gruppenstrukturen in Daten zu finden.}
        }
    \]

    k-means beschreibt Cluster durch harte Zuordnungen zu Zentren, während GMMs
    Cluster probabilistisch als Mischung von Verteilungen modellieren.
\end{answerbox}

%%%%%%%%%%%%%%%%%%%%%%%%%%%%%%%%%%%%%%%%%%%%%%%%%%%%%%%%%%%%%%%%%%%
\subsection{The Problem of Finding Clusters in Data}
%%%%%%%%%%%%%%%%%%%%%%%%%%%%%%%%%%%%%%%%%%%%%%%%%%%%%%%%%%%%%%%%%%%

\begin{notebox}
    Beim Clustering sind keine Labels gegeben.

    Man beobachtet nur Datenpunkte:
    \[
        x_1,\dots,x_N \in \mathbb{R}^D.
    \]

    Ziel ist es, diese Datenpunkte in Gruppen einzuteilen, sodass Punkte
    innerhalb derselben Gruppe einander ähnlich sind und Punkte aus
    verschiedenen Gruppen möglichst unterschiedlich sind.
\end{notebox}

\begin{solutionbox}
    Clustering ist ein Beispiel für \emph{unsupervised learning}.

    Im Gegensatz zur Klassifikation gibt es keine bekannten Zielwerte:
    \[
        (x_n,t_n)
    \]
    sondern nur:
    \[
        x_n.
    \]

    Das Modell muss also selbst aus der Struktur der Daten ableiten, welche
    Gruppen sinnvoll sind.

    Typische Anwendungen sind:

    \begin{itemize}
        \item Gruppierung von Kundendaten,
        \item Segmentierung von Bildern,
        \item Finden biologischer Zelltypen,
        \item Strukturierung großer Datensätze,
        \item Vorverarbeitung für weitere Machine-Learning-Methoden.
    \end{itemize}
\end{solutionbox}

\begin{answerbox}
    Clustering bedeutet:

    \[
        \boxed{
            \text{Finde Gruppen in ungelabelten Daten } x_1,\dots,x_N.
        }
    \]

    Die Schwierigkeit ist, dass nicht vorgegeben ist, was ein ``guter'' Cluster
    genau ist. Daher braucht man ein Objective oder ein probabilistisches Modell.
\end{answerbox}

%%%%%%%%%%%%%%%%%%%%%%%%%%%%%%%%%%%%%%%%%%%%%%%%%%%%%%%%%%%%%%%%%%%
\subsection{k-means: Basic Idea}
%%%%%%%%%%%%%%%%%%%%%%%%%%%%%%%%%%%%%%%%%%%%%%%%%%%%%%%%%%%%%%%%%%%

\begin{notebox}
    k-means ist eines der einfachsten und wichtigsten Clustering-Verfahren.

    Man nimmt an, dass es \(K\) Cluster gibt.

    Jeder Cluster wird durch ein Zentrum beschrieben:
    \[
        \mu_1,\dots,\mu_K \in \mathbb{R}^D.
    \]

    Jeder Datenpunkt wird genau einem Clusterzentrum zugeordnet.
\end{notebox}

\begin{solutionbox}
    Die Grundidee von k-means ist:

    \begin{itemize}
        \item Wähle \(K\) Clusterzentren.
        \item Ordne jeden Datenpunkt dem nächstgelegenen Zentrum zu.
        \item Aktualisiere jedes Zentrum als Mittelwert der ihm zugeordneten Datenpunkte.
        \item Wiederhole diese Schritte bis zur Konvergenz.
    \end{itemize}

    k-means erzeugt also eine harte Clusterzuordnung:

    \[
        x_n \longmapsto k_n,
        \qquad
        k_n \in \{1,\dots,K\}.
    \]
\end{solutionbox}

\begin{answerbox}
    Bei k-means wird jeder Datenpunkt genau einem Cluster zugeordnet:

    \[
        \boxed{
            k_n
            =
            \operatorname*{arg\,min}_{k}
            \|x_n-\mu_k\|^2
        }
    \]

    Danach wird jedes Zentrum als Mittelwert seiner zugeordneten Punkte
    aktualisiert.
\end{answerbox}

%%%%%%%%%%%%%%%%%%%%%%%%%%%%%%%%%%%%%%%%%%%%%%%%%%%%%%%%%%%%%%%%%%%
\subsection{k-means: Definition of an Ad-hoc Objective}
%%%%%%%%%%%%%%%%%%%%%%%%%%%%%%%%%%%%%%%%%%%%%%%%%%%%%%%%%%%%%%%%%%%

\begin{notebox}
    k-means kann über ein ad-hoc Objective formuliert werden.

    Die Idee ist:

    \[
        \text{Punkte sollen möglichst nahe an ihren jeweiligen Clusterzentren liegen.}
    \]

    Dafür minimiert man die Summe der quadratischen Abstände zu den
    Clusterzentren.
\end{notebox}

\begin{calculationbox}
    Sei \(r_{nk}\) eine binäre Zuordnungsvariable:

    \[
        r_{nk}
        =
        \begin{cases}
            1, & \text{wenn } x_n \text{ Cluster } k \text{ zugeordnet ist},\\
            0, & \text{sonst}.
        \end{cases}
    \]

    Für jeden Datenpunkt gilt:

    \[
        \sum_{k=1}^{K} r_{nk} = 1.
    \]

    Das k-means Objective lautet:

    \[
        J
        =
        \sum_{n=1}^{N}
        \sum_{k=1}^{K}
        r_{nk}
        \|x_n-\mu_k\|^2.
    \]

    Ziel ist:

    \[
        \min_{\{r_{nk}\},\{\mu_k\}}
        J.
    \]
\end{calculationbox}

\begin{answerbox}
    Das k-means Objective ist:

    \[
        \boxed{
            J
            =
            \sum_{n=1}^{N}
            \sum_{k=1}^{K}
            r_{nk}
            \|x_n-\mu_k\|^2
        }
    \]

    Dabei ist \(r_{nk}\) die harte Clusterzuordnung und \(\mu_k\) das Zentrum
    von Cluster \(k\).
\end{answerbox}

%%%%%%%%%%%%%%%%%%%%%%%%%%%%%%%%%%%%%%%%%%%%%%%%%%%%%%%%%%%%%%%%%%%
\subsection{Derivation of Lloyd's Algorithm}
%%%%%%%%%%%%%%%%%%%%%%%%%%%%%%%%%%%%%%%%%%%%%%%%%%%%%%%%%%%%%%%%%%%

\begin{notebox}
    Lloyd's Algorithmus ist der Standardalgorithmus für k-means.

    Er minimiert das k-means Objective iterativ durch abwechselnde Optimierung:

    \begin{itemize}
        \item Fixiere die Zentren und optimiere die Zuordnungen.
        \item Fixiere die Zuordnungen und optimiere die Zentren.
    \end{itemize}
\end{notebox}

\begin{calculationbox}
    \textbf{Schritt 1: Update der Zuordnungen.}

    Sind die Zentren \(\mu_1,\dots,\mu_K\) fixiert, minimiert man für jeden
    Datenpunkt:

    \[
        \sum_{k=1}^{K}
        r_{nk}
        \|x_n-\mu_k\|^2.
    \]

    Da jeder Punkt genau einem Cluster zugeordnet wird, wählt man den Cluster
    mit dem kleinsten Abstand:

    \[
        r_{nk}
        =
        \begin{cases}
            1, & k = \operatorname*{arg\,min}_{j}\|x_n-\mu_j\|^2,\\
            0, & \text{sonst}.
        \end{cases}
    \]
\end{calculationbox}

\begin{calculationbox}
    \textbf{Schritt 2: Update der Zentren.}

    Sind die Zuordnungen \(r_{nk}\) fixiert, minimiert man

    \[
        J
        =
        \sum_{n=1}^{N}
        \sum_{k=1}^{K}
        r_{nk}
        \|x_n-\mu_k\|^2
    \]

    bezüglich \(\mu_k\).

    Für ein fixes \(k\) betrachtet man:

    \[
        J_k
        =
        \sum_{n=1}^{N}
        r_{nk}
        \|x_n-\mu_k\|^2.
    \]

    Ableitung nach \(\mu_k\):

    \[
        \frac{\partial J_k}{\partial \mu_k}
        =
        -2
        \sum_{n=1}^{N}
        r_{nk}
        (x_n-\mu_k).
    \]

    Gleich null setzen:

    \[
        \sum_{n=1}^{N}
        r_{nk}
        (x_n-\mu_k)
        =
        0.
    \]

    Also:

    \[
        \sum_{n=1}^{N} r_{nk}x_n
        =
        \mu_k
        \sum_{n=1}^{N} r_{nk}.
    \]

    Damit:

    \[
        \mu_k
        =
        \frac{
            \sum_{n=1}^{N} r_{nk}x_n
        }{
            \sum_{n=1}^{N} r_{nk}
        }.
    \]
\end{calculationbox}

\begin{answerbox}
    Lloyd's Algorithmus besteht aus zwei Updates:

    \[
        \boxed{
            r_{nk}
            =
            \begin{cases}
                1, & k = \operatorname*{arg\,min}_{j}\|x_n-\mu_j\|^2,\\
                0, & \text{sonst}
            \end{cases}
        }
    \]

    und

    \[
        \boxed{
            \mu_k
            =
            \frac{
                \sum_{n=1}^{N} r_{nk}x_n
            }{
                \sum_{n=1}^{N} r_{nk}
            }
        }
    \]

    Der Algorithmus verringert das k-means Objective in jedem Schritt oder lässt
    es unverändert.
\end{answerbox}

%%%%%%%%%%%%%%%%%%%%%%%%%%%%%%%%%%%%%%%%%%%%%%%%%%%%%%%%%%%%%%%%%%%
\subsection{Lloyd's Algorithm: Practical Procedure}
%%%%%%%%%%%%%%%%%%%%%%%%%%%%%%%%%%%%%%%%%%%%%%%%%%%%%%%%%%%%%%%%%%%

\begin{solutionbox}
    Praktisches Vorgehen:

    \begin{enumerate}
        \item Wähle \(K\), die Anzahl der Cluster.
        \item Initialisiere die Zentren \(\mu_1,\dots,\mu_K\), zum Beispiel zufällig.
        \item Ordne jeden Datenpunkt dem nächsten Zentrum zu.
        \item Aktualisiere jedes Zentrum als Mittelwert seiner zugeordneten Punkte.
        \item Wiederhole Zuordnung und Update bis sich nichts mehr ändert oder
        bis ein Abbruchkriterium erreicht ist.
    \end{enumerate}
\end{solutionbox}

\begin{notebox}
    k-means findet nicht garantiert das globale Minimum des Objectives.

    Das Ergebnis hängt von der Initialisierung ab.

    Deshalb startet man k-means in der Praxis oft mehrfach mit verschiedenen
    Initialisierungen und wählt die Lösung mit dem kleinsten Objective.
\end{notebox}

\begin{answerbox}
    k-means ist einfach und effizient, aber:

    \[
        \boxed{
            \text{Das Ergebnis hängt von } K \text{ und der Initialisierung ab.}
        }
    \]

    Außerdem bevorzugt k-means ungefähr kugelförmige Cluster ähnlicher Größe.
\end{answerbox}

%%%%%%%%%%%%%%%%%%%%%%%%%%%%%%%%%%%%%%%%%%%%%%%%%%%%%%%%%%%%%%%%%%%
\subsection{Fuzzy-k-means}
%%%%%%%%%%%%%%%%%%%%%%%%%%%%%%%%%%%%%%%%%%%%%%%%%%%%%%%%%%%%%%%%%%%

\begin{notebox}
    Beim normalen k-means gehört jeder Datenpunkt genau zu einem Cluster.

    Das nennt man harte Zuordnung.

    Fuzzy-k-means erlaubt dagegen weiche Zuordnungen:

    \[
        r_{nk} \in [0,1].
    \]

    Dabei beschreibt \(r_{nk}\), wie stark der Datenpunkt \(x_n\) zu Cluster
    \(k\) gehört.
\end{notebox}

\begin{solutionbox}
    Für jeden Datenpunkt gilt weiterhin:

    \[
        \sum_{k=1}^{K} r_{nk}=1.
    \]

    Aber jetzt darf ein Punkt teilweise mehreren Clustern angehören.

    Das ist nützlich, wenn Cluster überlappen oder wenn Zuordnungen unsicher
    sind.

    Fuzzy-k-means ist damit ein Zwischenschritt zwischen hartem k-means und
    probabilistischen Mischmodellen wie Gaussian Mixture Models.
\end{solutionbox}

\begin{calculationbox}
    Ein typisches Fuzzy-k-means Objective ist:

    \[
        J_m
        =
        \sum_{n=1}^{N}
        \sum_{k=1}^{K}
        r_{nk}^{m}
        \|x_n-\mu_k\|^2.
    \]

    Dabei ist

    \[
        m>1
    \]

    ein Fuzziness-Parameter.

    Für größere \(m\) werden die Zuordnungen weicher.
\end{calculationbox}

\begin{answerbox}
    Fuzzy-k-means ersetzt harte Zuordnungen durch weiche Memberships:

    \[
        \boxed{
            r_{nk}\in[0,1],
            \qquad
            \sum_{k=1}^{K}r_{nk}=1.
        }
    \]

    Dadurch kann ein Punkt teilweise zu mehreren Clustern gehören.
\end{answerbox}

%%%%%%%%%%%%%%%%%%%%%%%%%%%%%%%%%%%%%%%%%%%%%%%%%%%%%%%%%%%%%%%%%%%
\subsection{Hierarchical Clustering}
%%%%%%%%%%%%%%%%%%%%%%%%%%%%%%%%%%%%%%%%%%%%%%%%%%%%%%%%%%%%%%%%%%%

\begin{notebox}
    Hierarchisches Clustering erzeugt nicht nur eine einzelne Clusterlösung,
    sondern eine Hierarchie von Clustern.

    Das Ergebnis wird oft als Dendrogramm dargestellt.

    Es gibt zwei Haupttypen:

    \begin{itemize}
        \item agglomeratives Clustering,
        \item divisives Clustering.
    \end{itemize}
\end{notebox}

\begin{solutionbox}
    Beim \textbf{agglomerativen Clustering} startet man mit vielen kleinen
    Clustern.

    Meist startet jeder Datenpunkt als eigener Cluster.

    Dann werden schrittweise die ähnlichsten Cluster zusammengeführt.

    Beim \textbf{divisiven Clustering} startet man mit einem großen Cluster und
    teilt diesen schrittweise auf.

    In der Praxis ist agglomeratives Clustering besonders verbreitet.
\end{solutionbox}

\begin{answerbox}
    Hierarchisches Clustering liefert eine Clusterhierarchie:

    \[
        \boxed{
            \text{kleine Cluster}
            \longleftrightarrow
            \text{große Cluster}
        }
    \]

    Ein Dendrogramm zeigt, bei welchen Distanzen Cluster zusammengeführt werden.
\end{answerbox}

%%%%%%%%%%%%%%%%%%%%%%%%%%%%%%%%%%%%%%%%%%%%%%%%%%%%%%%%%%%%%%%%%%%
\subsection{Agglomerative Clustering}
%%%%%%%%%%%%%%%%%%%%%%%%%%%%%%%%%%%%%%%%%%%%%%%%%%%%%%%%%%%%%%%%%%%

\begin{notebox}
    Agglomeratives Clustering ist ein bottom-up Verfahren.

    Es startet mit \(N\) Clustern:

    \[
        C_1=\{x_1\},\dots,C_N=\{x_N\}.
    \]

    Danach werden wiederholt die zwei ähnlichsten Cluster zusammengeführt.
\end{notebox}

\begin{solutionbox}
    Algorithmus:

    \begin{enumerate}
        \item Starte mit jedem Datenpunkt als eigenem Cluster.
        \item Berechne Distanzen zwischen allen Clustern.
        \item Führe die zwei nächstgelegenen Cluster zusammen.
        \item Aktualisiere die Distanzen.
        \item Wiederhole, bis nur noch ein Cluster übrig ist oder bis die
        gewünschte Anzahl von Clustern erreicht ist.
    \end{enumerate}
\end{solutionbox}

\begin{notebox}
    Entscheidend ist, wie man die Distanz zwischen zwei Clustern definiert.

    Häufige Linkage-Kriterien sind:

    \begin{itemize}
        \item single linkage,
        \item complete linkage,
        \item average linkage,
        \item centroid linkage,
        \item Ward linkage.
    \end{itemize}
\end{notebox}

\begin{calculationbox}
    Beispiele für Cluster-Distanzen:

    \textbf{Single linkage:}
    \[
        d(C_a,C_b)
        =
        \min_{x\in C_a,\,y\in C_b}
        \|x-y\|.
    \]

    \textbf{Complete linkage:}
    \[
        d(C_a,C_b)
        =
        \max_{x\in C_a,\,y\in C_b}
        \|x-y\|.
    \]

    \textbf{Average linkage:}
    \[
        d(C_a,C_b)
        =
        \frac{1}{|C_a||C_b|}
        \sum_{x\in C_a}
        \sum_{y\in C_b}
        \|x-y\|.
    \]
\end{calculationbox}

\begin{answerbox}
    Agglomeratives Clustering ist:

    \[
        \boxed{
            \text{bottom-up: einzelne Punkte } \longrightarrow \text{ große Cluster.}
        }
    \]

    Die Wahl des Linkage-Kriteriums beeinflusst stark, welche Cluster entstehen.
\end{answerbox}

%%%%%%%%%%%%%%%%%%%%%%%%%%%%%%%%%%%%%%%%%%%%%%%%%%%%%%%%%%%%%%%%%%%
\subsection{Gaussian Mixture Models}
%%%%%%%%%%%%%%%%%%%%%%%%%%%%%%%%%%%%%%%%%%%%%%%%%%%%%%%%%%%%%%%%%%%

\begin{notebox}
    Gaussian Mixture Models, kurz GMMs, modellieren Daten als Mischung mehrerer
    Gauß-Verteilungen.

    Statt harter Clusterzuordnung wie bei k-means gibt es probabilistische
    Zugehörigkeiten zu Komponenten.
\end{notebox}

\begin{solutionbox}
    Ein GMM mit \(K\) Komponenten hat Parameter:

    \[
        \pi_1,\dots,\pi_K
    \]

    für die Mischgewichte,

    \[
        \mu_1,\dots,\mu_K
    \]

    für die Mittelwerte und

    \[
        \Sigma_1,\dots,\Sigma_K
    \]

    für die Kovarianzmatrizen.

    Die Mischgewichte erfüllen:

    \[
        \pi_k \geq 0,
        \qquad
        \sum_{k=1}^{K}\pi_k=1.
    \]
\end{solutionbox}

\begin{calculationbox}
    Die Dichte eines GMMs lautet:

    \[
        p(x)
        =
        \sum_{k=1}^{K}
        \pi_k
        \mathcal{N}(x\mid \mu_k,\Sigma_k).
    \]

    Jede Komponente beschreibt einen möglichen Cluster.

    Die Gesamtverteilung ist eine gewichtete Summe dieser Komponenten.
\end{calculationbox}

\begin{answerbox}
    Ein Gaussian Mixture Model ist:

    \[
        \boxed{
            p(x)
            =
            \sum_{k=1}^{K}
            \pi_k
            \mathcal{N}(x\mid \mu_k,\Sigma_k)
        }
    \]

    Dabei sind \(\pi_k\) die Mischgewichte, \(\mu_k\) die Mittelwerte und
    \(\Sigma_k\) die Kovarianzmatrizen.
\end{answerbox}

%%%%%%%%%%%%%%%%%%%%%%%%%%%%%%%%%%%%%%%%%%%%%%%%%%%%%%%%%%%%%%%%%%%
\subsection{Latent Variables and Responsibilities in GMMs}
%%%%%%%%%%%%%%%%%%%%%%%%%%%%%%%%%%%%%%%%%%%%%%%%%%%%%%%%%%%%%%%%%%%

\begin{notebox}
    In GMMs kann man sich vorstellen, dass jeder Datenpunkt zuerst eine
    versteckte Clusterkomponente \(z_n\) auswählt.

    Diese latente Variable ist nicht direkt beobachtet.

    Sie gibt an, aus welcher Gauß-Komponente der Datenpunkt entstanden ist.
\end{notebox}

\begin{solutionbox}
    Für einen Datenpunkt \(x_n\) ist die posterior probability, dass er zu
    Komponente \(k\) gehört:

    \[
        \gamma_{nk}
        =
        p(z_n=k \mid x_n).
    \]

    Diese Größe heißt \emph{responsibility}.

    Sie beschreibt, wie stark Komponente \(k\) für Datenpunkt \(x_n\)
    verantwortlich ist.
\end{solutionbox}

\begin{calculationbox}
    Nach Bayes' Regel gilt:

    \[
        \gamma_{nk}
        =
        p(z_n=k \mid x_n)
        =
        \frac{
            \pi_k
            \mathcal{N}(x_n\mid \mu_k,\Sigma_k)
        }{
            \sum_{j=1}^{K}
            \pi_j
            \mathcal{N}(x_n\mid \mu_j,\Sigma_j)
        }.
    \]
\end{calculationbox}

\begin{answerbox}
    Die Responsibility in einem GMM ist:

    \[
        \boxed{
            \gamma_{nk}
            =
            \frac{
                \pi_k
                \mathcal{N}(x_n\mid \mu_k,\Sigma_k)
            }{
                \sum_{j=1}^{K}
                \pi_j
                \mathcal{N}(x_n\mid \mu_j,\Sigma_j)
            }
        }
    \]

    Sie ist eine weiche Clusterzuordnung.
\end{answerbox}

%%%%%%%%%%%%%%%%%%%%%%%%%%%%%%%%%%%%%%%%%%%%%%%%%%%%%%%%%%%%%%%%%%%
\subsection{Maximum Likelihood Parameter Estimation for GMMs}
%%%%%%%%%%%%%%%%%%%%%%%%%%%%%%%%%%%%%%%%%%%%%%%%%%%%%%%%%%%%%%%%%%%

\begin{notebox}
    Die Parameter eines GMMs können mit Maximum Likelihood geschätzt werden.

    Für Daten
    \[
        x_1,\dots,x_N
    \]
    lautet die Likelihood:

    \[
        L(\theta)
        =
        \prod_{n=1}^{N}
        p(x_n\mid \theta).
    \]

    Dabei enthält \(\theta\) alle Mischgewichte, Mittelwerte und Kovarianzen.
\end{notebox}

\begin{calculationbox}
    Für ein GMM ist:

    \[
        p(x_n\mid \theta)
        =
        \sum_{k=1}^{K}
        \pi_k
        \mathcal{N}(x_n\mid \mu_k,\Sigma_k).
    \]

    Damit lautet die Log-Likelihood:

    \[
        \ell(\theta)
        =
        \sum_{n=1}^{N}
        \log
        \left(
            \sum_{k=1}^{K}
            \pi_k
            \mathcal{N}(x_n\mid \mu_k,\Sigma_k)
        \right).
    \]

    Das Problem ist, dass der Logarithmus über einer Summe steht.

    Dadurch kann man die Parameter nicht so einfach direkt durch Ableiten lösen.
\end{calculationbox}

\begin{answerbox}
    Die GMM-Log-Likelihood ist:

    \[
        \boxed{
            \ell(\theta)
            =
            \sum_{n=1}^{N}
            \log
            \left(
                \sum_{k=1}^{K}
                \pi_k
                \mathcal{N}(x_n\mid \mu_k,\Sigma_k)
            \right)
        }
    \]

    Wegen des Logarithmus über der Summe verwendet man typischerweise den
    Expectation-Maximization-Algorithmus.
\end{answerbox}

%%%%%%%%%%%%%%%%%%%%%%%%%%%%%%%%%%%%%%%%%%%%%%%%%%%%%%%%%%%%%%%%%%%
\subsection{Free Energy and ELBO as Auxiliary Objectives}
%%%%%%%%%%%%%%%%%%%%%%%%%%%%%%%%%%%%%%%%%%%%%%%%%%%%%%%%%%%%%%%%%%%

\begin{notebox}
    Die direkte Maximierung der GMM-Log-Likelihood ist schwierig.

    Deshalb führt man ein Hilfsobjective ein: die Free Energy beziehungsweise
    Evidence Lower Bound, kurz ELBO.

    Die Idee ist:

    \[
        \text{Optimiere eine untere Schranke der Log-Likelihood.}
    \]
\end{notebox}

\begin{solutionbox}
    Für latente Variablen \(z\) gilt:

    \[
        \log p(x)
        =
        \log
        \sum_z
        p(x,z).
    \]

    Man führt eine Hilfsverteilung \(q(z)\) ein.

    Dann kann man schreiben:

    \[
        \log p(x)
        =
        \log
        \sum_z
        q(z)
        \frac{p(x,z)}{q(z)}.
    \]

    Mit Jensen's Ungleichung erhält man eine untere Schranke:

    \[
        \log p(x)
        \geq
        \sum_z
        q(z)
        \log
        \frac{p(x,z)}{q(z)}.
    \]

    Diese untere Schranke ist die ELBO.
\end{solutionbox}

\begin{calculationbox}
    Für alle Daten lautet die ELBO:

    \[
        \mathcal{L}(q,\theta)
        =
        \sum_{n=1}^{N}
        \sum_{k=1}^{K}
        q_{nk}
        \log
        \frac{
            \pi_k
            \mathcal{N}(x_n\mid \mu_k,\Sigma_k)
        }{
            q_{nk}
        }.
    \]

    Dabei approximiert \(q_{nk}\) die Zugehörigkeit von Datenpunkt \(x_n\) zu
    Komponente \(k\).

    Im EM-Algorithmus wählt man im E-Schritt:

    \[
        q_{nk}
        =
        \gamma_{nk}
        =
        p(z_n=k\mid x_n).
    \]
\end{calculationbox}

\begin{answerbox}
    Die ELBO ist eine untere Schranke der Log-Likelihood:

    \[
        \boxed{
            \log p(x)
            \geq
            \sum_z
            q(z)
            \log
            \frac{p(x,z)}{q(z)}
        }
    \]

    Für GMMs führt diese Hilfsfunktion zum EM-Algorithmus:

    \[
        \boxed{
            \text{E-Schritt: Verantwortlichkeiten berechnen}
        }
    \]

    \[
        \boxed{
            \text{M-Schritt: Parameter mit diesen Verantwortlichkeiten aktualisieren.}
        }
    \]
\end{answerbox}

%%%%%%%%%%%%%%%%%%%%%%%%%%%%%%%%%%%%%%%%%%%%%%%%%%%%%%%%%%%%%%%%%%%
\subsection{Update of the Mean in GMMs}
%%%%%%%%%%%%%%%%%%%%%%%%%%%%%%%%%%%%%%%%%%%%%%%%%%%%%%%%%%%%%%%%%%%

\begin{notebox}
    Im M-Schritt des EM-Algorithmus werden die Parameter des GMMs aktualisiert.

    Besonders wichtig ist das Update des Mittelwerts \(\mu_k\).

    Die Verantwortlichkeiten \(\gamma_{nk}\) wirken dabei wie weiche Gewichte.
\end{notebox}

\begin{calculationbox}
    Definiere die effektive Anzahl von Punkten in Komponente \(k\):

    \[
        N_k
        =
        \sum_{n=1}^{N}
        \gamma_{nk}.
    \]

    Das Update des Mittelwerts lautet:

    \[
        \mu_k^{\mathrm{new}}
        =
        \frac{1}{N_k}
        \sum_{n=1}^{N}
        \gamma_{nk}x_n.
    \]

    Das ist ein gewichteter Mittelwert der Datenpunkte.
\end{calculationbox}

\begin{solutionbox}
    Interpretation:

    Wenn \(\gamma_{nk}\) groß ist, dann gehört \(x_n\) stark zu Komponente \(k\)
    und beeinflusst den Mittelwert \(\mu_k\) stark.

    Wenn \(\gamma_{nk}\) klein ist, dann hat \(x_n\) nur wenig Einfluss auf
    \(\mu_k\).

    Im Unterschied zu k-means sind die Zuordnungen also nicht hart, sondern
    weich.
\end{solutionbox}

\begin{answerbox}
    Das Mittelwert-Update im GMM ist:

    \[
        \boxed{
            \mu_k^{\mathrm{new}}
            =
            \frac{
                \sum_{n=1}^{N}
                \gamma_{nk}x_n
            }{
                \sum_{n=1}^{N}
                \gamma_{nk}
            }
        }
    \]

    Es ist ein gewichteter Mittelwert mit Responsibilities \(\gamma_{nk}\).
\end{answerbox}

%%%%%%%%%%%%%%%%%%%%%%%%%%%%%%%%%%%%%%%%%%%%%%%%%%%%%%%%%%%%%%%%%%%
\subsection{Relation Between k-means and GMMs}
%%%%%%%%%%%%%%%%%%%%%%%%%%%%%%%%%%%%%%%%%%%%%%%%%%%%%%%%%%%%%%%%%%%

\begin{notebox}
    k-means und GMMs sind eng verwandt.

    k-means verwendet harte Zuordnungen und Clusterzentren.

    GMMs verwenden weiche Zuordnungen und probabilistische Gauß-Komponenten.
\end{notebox}

\begin{solutionbox}
    k-means kann als Grenzfall eines GMMs verstanden werden, wenn:

    \begin{itemize}
        \item alle Kovarianzmatrizen gleich und isotrop sind,
        \item die Varianz sehr klein wird,
        \item die weichen Responsibilities fast zu harten Zuordnungen werden.
    \end{itemize}

    Dann gehört jeder Punkt praktisch nur noch zur nächstgelegenen Komponente,
    ähnlich wie bei k-means.
\end{solutionbox}

\begin{answerbox}
    Vergleich:

    \[
        \boxed{
            \text{k-means: harte Zuordnungen, Zentren, quadratische Abstände}
        }
    \]

    \[
        \boxed{
            \text{GMM: weiche Zuordnungen, Wahrscheinlichkeiten, Kovarianzen}
        }
    \]

    GMMs sind flexibler, aber auch komplexer zu trainieren.
\end{answerbox}

%%%%%%%%%%%%%%%%%%%%%%%%%%%%%%%%%%%%%%%%%%%%%%%%%%%%%%%%%%%%%%%%%%%
\subsection{Discussion}
%%%%%%%%%%%%%%%%%%%%%%%%%%%%%%%%%%%%%%%%%%%%%%%%%%%%%%%%%%%%%%%%%%%

\begin{notebox}
    Clustering ist grundsätzlich schwieriger zu bewerten als supervised learning,
    weil keine Labels gegeben sind.

    Deshalb hängt die Qualität eines Clustering-Ergebnisses stark davon ab,
    welches Ziel man verfolgt.
\end{notebox}

\begin{solutionbox}
    Vorteile von k-means:

    \begin{itemize}
        \item einfach,
        \item schnell,
        \item gut für große Datensätze,
        \item leicht zu interpretieren.
    \end{itemize}

    Nachteile von k-means:

    \begin{itemize}
        \item Anzahl der Cluster \(K\) muss vorgegeben werden,
        \item empfindlich gegenüber Initialisierung,
        \item bevorzugt kugelförmige Cluster,
        \item harte Zuordnungen,
        \item empfindlich gegenüber Ausreißern.
    \end{itemize}
\end{solutionbox}

\begin{solutionbox}
    Vorteile von hierarchischem Clustering:

    \begin{itemize}
        \item liefert eine ganze Hierarchie von Clustern,
        \item Anzahl der Cluster muss nicht sofort festgelegt werden,
        \item Dendrogramm ist interpretierbar.
    \end{itemize}

    Nachteile:

    \begin{itemize}
        \item oft rechenintensiver,
        \item frühe Zusammenführungen können nicht rückgängig gemacht werden,
        \item Ergebnis hängt stark vom Linkage-Kriterium ab.
    \end{itemize}
\end{solutionbox}

\begin{solutionbox}
    Vorteile von GMMs:

    \begin{itemize}
        \item probabilistische Interpretation,
        \item weiche Clusterzuordnungen,
        \item unterschiedliche Clusterformen durch Kovarianzmatrizen,
        \item Unsicherheit der Zuordnung wird sichtbar.
    \end{itemize}

    Nachteile:

    \begin{itemize}
        \item komplexer als k-means,
        \item EM kann in lokalen Optima landen,
        \item Anzahl der Komponenten \(K\) muss gewählt werden,
        \item Kovarianzschätzung kann bei wenig Daten instabil sein.
    \end{itemize}
\end{solutionbox}

\begin{answerbox}
    Die Wahl des Clustering-Verfahrens hängt von der Datenstruktur ab:

    \[
        \boxed{
            \text{k-means: einfache, kompakte Cluster}
        }
    \]

    \[
        \boxed{
            \text{Hierarchisch: Clusterhierarchie und Dendrogramme}
        }
    \]

    \[
        \boxed{
            \text{GMMs: probabilistische, weiche und elliptische Cluster}
        }
    \]
\end{answerbox}

%%%%%%%%%%%%%%%%%%%%%%%%%%%%%%%%%%%%%%%%%%%%%%%%%%%%%%%%%%%%%%%%%%%
\subsection{Prüfungsrelevante Kurzfassung}
%%%%%%%%%%%%%%%%%%%%%%%%%%%%%%%%%%%%%%%%%%%%%%%%%%%%%%%%%%%%%%%%%%%

\begin{answerbox}
    Lecture 8 behandelt Clustering mit k-means, hierarchischem Clustering und
    Gaussian Mixture Models.

    Die wichtigsten Punkte sind:

    \begin{itemize}
        \item Clustering sucht Gruppen in ungelabelten Daten:
        \[
            x_1,\dots,x_N.
        \]

        \item k-means minimiert:
        \[
            J
            =
            \sum_{n=1}^{N}
            \sum_{k=1}^{K}
            r_{nk}
            \|x_n-\mu_k\|^2.
        \]

        \item k-means Zuordnungsupdate:
        \[
            r_{nk}
            =
            \begin{cases}
                1, & k=\operatorname*{arg\,min}_{j}\|x_n-\mu_j\|^2,\\
                0, & \text{sonst}.
            \end{cases}
        \]

        \item k-means Zentrenupdate:
        \[
            \mu_k
            =
            \frac{
                \sum_{n=1}^{N}r_{nk}x_n
            }{
                \sum_{n=1}^{N}r_{nk}
            }.
        \]

        \item Fuzzy-k-means verwendet weiche Zuordnungen:
        \[
            r_{nk}\in[0,1],
            \qquad
            \sum_{k=1}^{K}r_{nk}=1.
        \]

        \item Agglomeratives Clustering startet mit einzelnen Punkten und
        verbindet schrittweise die ähnlichsten Cluster.

        \item Ein Gaussian Mixture Model ist:
        \[
            p(x)
            =
            \sum_{k=1}^{K}
            \pi_k
            \mathcal{N}(x\mid \mu_k,\Sigma_k).
        \]

        \item GMM-Responsibilities:
        \[
            \gamma_{nk}
            =
            \frac{
                \pi_k
                \mathcal{N}(x_n\mid \mu_k,\Sigma_k)
            }{
                \sum_{j=1}^{K}
                \pi_j
                \mathcal{N}(x_n\mid \mu_j,\Sigma_j)
            }.
        \]

        \item GMM-Log-Likelihood:
        \[
            \ell(\theta)
            =
            \sum_{n=1}^{N}
            \log
            \left(
                \sum_{k=1}^{K}
                \pi_k
                \mathcal{N}(x_n\mid \mu_k,\Sigma_k)
            \right).
        \]

        \item ELBO:
        \[
            \log p(x)
            \geq
            \sum_z
            q(z)
            \log
            \frac{p(x,z)}{q(z)}.
        \]

        \item GMM-Mittelwertupdate:
        \[
            \mu_k^{\mathrm{new}}
            =
            \frac{
                \sum_{n=1}^{N}
                \gamma_{nk}x_n
            }{
                \sum_{n=1}^{N}
                \gamma_{nk}
            }.
        \]
    \end{itemize}

    Die Kernaussage lautet:

    \[
        \boxed{
            \text{k-means, hierarchisches Clustering und GMMs sind verschiedene
            Wege, Gruppenstruktur in ungelabelten Daten zu modellieren.}
        }
    \]
\end{answerbox}

%%%%%%%%%%%%%%%%%%%%%%%%%%%%%%%%%%%%%%%%%%%%%%%%%%%%%%%%%%%%%%%%%%%
\section{Lecture 9: Gaussian Mixture Models and the Expectation Maximization Algorithm}
%%%%%%%%%%%%%%%%%%%%%%%%%%%%%%%%%%%%%%%%%%%%%%%%%%%%%%%%%%%%%%%%%%%

\subsection{Overview}

\begin{notebox}
    Lecture 9 setzt die Gaussian Mixture Models aus Lecture 8 fort und führt
    den Expectation-Maximization-Algorithmus, kurz EM-Algorithmus, systematisch
    ein.

    Die zentralen Themen sind:

    \begin{itemize}
        \item Unterschied zwischen Log-Likelihood und Free Energy,
        \item Posterior-Verteilungen als optimale Wahl für die Hilfsverteilungen,
        \item iterative Berechnung optimaler Mittelwerte,
        \item ein erster EM-Algorithmus,
        \item Allgemeinheit des EM-Algorithmus,
        \item EM als Meta-Algorithmus für generative Modelle.
    \end{itemize}
\end{notebox}

\begin{answerbox}
    Die Kernaussage von Lecture 9 ist:

    \[
        \boxed{
            \text{EM optimiert Modelle mit latenten Variablen, indem es
            abwechselnd Posterior-Verteilungen berechnet und Parameter aktualisiert.}
        }
    \]

    Für Gaussian Mixture Models bedeutet das:
    Erst berechnet man weiche Clusterzuordnungen, danach aktualisiert man die
    Modellparameter mit diesen Zuordnungen.
\end{answerbox}

%%%%%%%%%%%%%%%%%%%%%%%%%%%%%%%%%%%%%%%%%%%%%%%%%%%%%%%%%%%%%%%%%%%
\subsection{Gaussian Mixture Models Continued}
%%%%%%%%%%%%%%%%%%%%%%%%%%%%%%%%%%%%%%%%%%%%%%%%%%%%%%%%%%%%%%%%%%%

\begin{notebox}
    Ein Gaussian Mixture Model beschreibt Daten als Mischung mehrerer
    Gauß-Verteilungen.

    Für \(K\) Komponenten lautet das Modell:

    \[
        p(x)
        =
        \sum_{c=1}^{K}
        \pi_c
        \mathcal{N}(x\mid \mu_c,\Sigma_c).
    \]

    Dabei sind:

    \begin{itemize}
        \item \(\pi_c\) die Mischgewichte,
        \item \(\mu_c\) die Mittelwerte,
        \item \(\Sigma_c\) die Kovarianzmatrizen,
        \item \(c\) die latente Komponente beziehungsweise Clusterzugehörigkeit.
    \end{itemize}
\end{notebox}

\begin{solutionbox}
    Die Variable \(c\) ist latent, weil sie nicht direkt beobachtet wird.

    Für jeden Datenpunkt \(x^{(n)}\) wissen wir nicht, aus welcher Komponente er
    erzeugt wurde.

    Man kann sich den generativen Prozess so vorstellen:

    \begin{enumerate}
        \item Wähle eine Komponente:
        \[
            c^{(n)} \sim \operatorname{Categorical}(\pi_1,\dots,\pi_K).
        \]

        \item Ziehe den Datenpunkt aus der gewählten Gauß-Verteilung:
        \[
            x^{(n)} \sim \mathcal{N}(\mu_{c^{(n)}},\Sigma_{c^{(n)}}).
        \]
    \end{enumerate}

    Da \(c^{(n)}\) unbekannt ist, müssen wir über alle möglichen Komponenten
    summieren.
\end{solutionbox}

\begin{answerbox}
    Das GMM hat die Form:

    \[
        \boxed{
            p(x)
            =
            \sum_{c=1}^{K}
            \pi_c
            \mathcal{N}(x\mid \mu_c,\Sigma_c)
        }
    \]

    Die Clusterzugehörigkeit \(c\) ist eine latente Variable.
\end{answerbox}

%%%%%%%%%%%%%%%%%%%%%%%%%%%%%%%%%%%%%%%%%%%%%%%%%%%%%%%%%%%%%%%%%%%
\subsection{Difference Between Log-Likelihood and Free Energy}
%%%%%%%%%%%%%%%%%%%%%%%%%%%%%%%%%%%%%%%%%%%%%%%%%%%%%%%%%%%%%%%%%%%

\begin{notebox}
    Für beobachtete Daten
    \[
        x^{(1)},\dots,x^{(N)}
    \]
    lautet die Log-Likelihood eines GMMs:

    \[
        \log p(X\mid \theta)
        =
        \sum_{n=1}^{N}
        \log p(x^{(n)}\mid \theta).
    \]

    Wegen der latenten Variable \(c\) gilt:

    \[
        p(x^{(n)}\mid \theta)
        =
        \sum_{c}
        p(x^{(n)},c\mid \theta).
    \]

    Also:

    \[
        \log p(X\mid \theta)
        =
        \sum_{n=1}^{N}
        \log
        \sum_{c}
        p(x^{(n)},c\mid \theta).
    \]
\end{notebox}

\begin{solutionbox}
    Das Problem ist der Ausdruck:

    \[
        \log
        \sum_c
        p(x^{(n)},c\mid \theta).
    \]

    Der Logarithmus steht außerhalb der Summe über die latente Variable.

    Dadurch ist die direkte Optimierung der Log-Likelihood schwierig, weil sich
    die Terme für die einzelnen Komponenten nicht einfach trennen lassen.

    Die Free Energy beziehungsweise ELBO ersetzt dieses schwierige Objective
    durch eine untere Schranke, die leichter optimiert werden kann.
\end{solutionbox}

\begin{calculationbox}
    Für jeden Datenpunkt führt man eine Hilfsverteilung
    \[
        q^{(n)}(c)
    \]
    über die latente Komponente ein.

    Dann gilt:

    \[
        \log p(x^{(n)}\mid \theta)
        =
        \log
        \sum_c
        p(x^{(n)},c\mid \theta).
    \]

    Da
    \[
        \sum_c q^{(n)}(c)=1,
    \]
    kann man schreiben:

    \[
        \log p(x^{(n)}\mid \theta)
        =
        \log
        \sum_c
        q^{(n)}(c)
        \frac{
            p(x^{(n)},c\mid \theta)
        }{
            q^{(n)}(c)
        }.
    \]

    Mit Jensen's Ungleichung erhält man:

    \[
        \log p(x^{(n)}\mid \theta)
        \geq
        \sum_c
        q^{(n)}(c)
        \log
        \frac{
            p(x^{(n)},c\mid \theta)
        }{
            q^{(n)}(c)
        }.
    \]
\end{calculationbox}

\begin{calculationbox}
    Die Free Energy beziehungsweise ELBO für alle Daten ist:

    \[
        \mathcal{F}(q,\theta)
        =
        \sum_{n=1}^{N}
        \sum_c
        q^{(n)}(c)
        \log
        \frac{
            p(x^{(n)},c\mid \theta)
        }{
            q^{(n)}(c)
        }.
    \]

    Damit gilt:

    \[
        \log p(X\mid \theta)
        \geq
        \mathcal{F}(q,\theta).
    \]

    Die Free Energy ist also eine untere Schranke an die Log-Likelihood.
\end{calculationbox}

\begin{answerbox}
    Der Unterschied ist:

    \[
        \boxed{
            \log p(X\mid \theta)
            =
            \text{eigentliches Objective}
        }
    \]

    \[
        \boxed{
            \mathcal{F}(q,\theta)
            =
            \text{untere Schranke / Hilfsobjective}
        }
    \]

    mit

    \[
        \boxed{
            \log p(X\mid \theta)
            \geq
            \mathcal{F}(q,\theta).
        }
    \]

    EM optimiert diese untere Schranke abwechselnd bezüglich \(q\) und
    \(\theta\).
\end{answerbox}

%%%%%%%%%%%%%%%%%%%%%%%%%%%%%%%%%%%%%%%%%%%%%%%%%%%%%%%%%%%%%%%%%%%
\subsection{Posterior Distributions as Optimal Choices for \(q^{(n)}(c)\)}
%%%%%%%%%%%%%%%%%%%%%%%%%%%%%%%%%%%%%%%%%%%%%%%%%%%%%%%%%%%%%%%%%%%

\begin{notebox}
    Im E-Schritt des EM-Algorithmus wird die Hilfsverteilung
    \[
        q^{(n)}(c)
    \]
    optimal gewählt.

    Die optimale Wahl ist der Posterior der latenten Variable:

    \[
        q^{(n)}(c)
        =
        p(c\mid x^{(n)},\theta).
    \]
\end{notebox}

\begin{solutionbox}
    Die Free Energy ist eine untere Schranke an die Log-Likelihood.

    Der Abstand zwischen Log-Likelihood und Free Energy ist eine
    Kullback-Leibler-Divergenz:

    \[
        \log p(X\mid \theta)
        -
        \mathcal{F}(q,\theta)
        =
        \sum_{n=1}^{N}
        \operatorname{KL}
        \left(
            q^{(n)}(c)
            \,\|\, 
            p(c\mid x^{(n)},\theta)
        \right).
    \]

    Da eine KL-Divergenz immer nichtnegativ ist, gilt:

    \[
        \mathcal{F}(q,\theta)
        \leq
        \log p(X\mid \theta).
    \]

    Die Schranke wird genau dann eng, wenn

    \[
        q^{(n)}(c)
        =
        p(c\mid x^{(n)},\theta)
    \]

    für alle \(n\).
\end{solutionbox}

\begin{calculationbox}
    Für ein GMM lautet der Posterior:

    \[
        p(c\mid x^{(n)},\theta)
        =
        \frac{
            p(c\mid \theta)p(x^{(n)}\mid c,\theta)
        }{
            p(x^{(n)}\mid \theta)
        }.
    \]

    Mit

    \[
        p(c\mid \theta)=\pi_c
    \]

    und

    \[
        p(x^{(n)}\mid c,\theta)
        =
        \mathcal{N}(x^{(n)}\mid \mu_c,\Sigma_c)
    \]

    folgt:

    \[
        q^{(n)}(c)
        =
        \frac{
            \pi_c
            \mathcal{N}(x^{(n)}\mid \mu_c,\Sigma_c)
        }{
            \sum_{j=1}^{K}
            \pi_j
            \mathcal{N}(x^{(n)}\mid \mu_j,\Sigma_j)
        }.
    \]

    Diese Größe nennt man auch Responsibility:
    \[
        \gamma_{nc}
        =
        q^{(n)}(c).
    \]
\end{calculationbox}

\begin{answerbox}
    Die optimale Wahl der Hilfsverteilung ist:

    \[
        \boxed{
            q^{(n)}(c)
            =
            p(c\mid x^{(n)},\theta)
        }
    \]

    Für GMMs:

    \[
        \boxed{
            q^{(n)}(c)
            =
            \gamma_{nc}
            =
            \frac{
                \pi_c
                \mathcal{N}(x^{(n)}\mid \mu_c,\Sigma_c)
            }{
                \sum_{j=1}^{K}
                \pi_j
                \mathcal{N}(x^{(n)}\mid \mu_j,\Sigma_j)
            }
        }
    \]

    Damit wird die Free-Energy-Schranke an der aktuellen Stelle eng.
\end{answerbox}

%%%%%%%%%%%%%%%%%%%%%%%%%%%%%%%%%%%%%%%%%%%%%%%%%%%%%%%%%%%%%%%%%%%
\subsection{Computation of Optimal Means Using an Iterative Algorithm}
%%%%%%%%%%%%%%%%%%%%%%%%%%%%%%%%%%%%%%%%%%%%%%%%%%%%%%%%%%%%%%%%%%%

\begin{notebox}
    Wenn die Zuordnungen der Datenpunkte zu den Komponenten unbekannt sind,
    können die Mittelwerte \(\mu_c\) nicht direkt wie bei vollständig gelabelten
    Daten berechnet werden.

    EM löst dieses Problem iterativ:

    \begin{itemize}
        \item Berechne weiche Zuordnungen \(q^{(n)}(c)\).
        \item Aktualisiere Mittelwerte mit diesen weichen Zuordnungen.
        \item Wiederhole.
    \end{itemize}
\end{notebox}

\begin{calculationbox}
    Für fixe Responsibilities \(q^{(n)}(c)\) enthält die Free Energy bezüglich
    der Mittelwerte Terme der Form:

    \[
        \sum_{n=1}^{N}
        \sum_c
        q^{(n)}(c)
        \log
        \mathcal{N}(x^{(n)}\mid \mu_c,\Sigma_c).
    \]

    Für eine Gauß-Verteilung ist der relevante Teil:

    \[
        -
        \frac{1}{2}
        \sum_{n=1}^{N}
        q^{(n)}(c)
        (x^{(n)}-\mu_c)^T
        \Sigma_c^{-1}
        (x^{(n)}-\mu_c).
    \]

    Maximierung bezüglich \(\mu_c\) ergibt einen gewichteten Mittelwert.
\end{calculationbox}

\begin{calculationbox}
    Definiere:

    \[
        N_c
        =
        \sum_{n=1}^{N}
        q^{(n)}(c).
    \]

    Dann lautet das Update:

    \[
        \mu_c^{\mathrm{new}}
        =
        \frac{1}{N_c}
        \sum_{n=1}^{N}
        q^{(n)}(c)x^{(n)}.
    \]

    Mit Responsibilities \(\gamma_{nc}\) schreibt man:

    \[
        \mu_c^{\mathrm{new}}
        =
        \frac{
            \sum_{n=1}^{N}
            \gamma_{nc}x^{(n)}
        }{
            \sum_{n=1}^{N}
            \gamma_{nc}
        }.
    \]
\end{calculationbox}

\begin{answerbox}
    Das iterative Mittelwertupdate ist:

    \[
        \boxed{
            \mu_c^{\mathrm{new}}
            =
            \frac{
                \sum_{n=1}^{N}
                q^{(n)}(c)x^{(n)}
            }{
                \sum_{n=1}^{N}
                q^{(n)}(c)
            }
        }
    \]

    beziehungsweise:

    \[
        \boxed{
            \mu_c^{\mathrm{new}}
            =
            \frac{
                \sum_{n=1}^{N}
                \gamma_{nc}x^{(n)}
            }{
                \sum_{n=1}^{N}
                \gamma_{nc}
            }.
        }
    \]

    Das ist ein gewichteter Mittelwert der Datenpunkte.
\end{answerbox}

%%%%%%%%%%%%%%%%%%%%%%%%%%%%%%%%%%%%%%%%%%%%%%%%%%%%%%%%%%%%%%%%%%%
\subsection{A First EM Algorithm for GMM Means}
%%%%%%%%%%%%%%%%%%%%%%%%%%%%%%%%%%%%%%%%%%%%%%%%%%%%%%%%%%%%%%%%%%%

\begin{notebox}
    Ein erster einfacher EM-Algorithmus für GMMs kann formuliert werden, wenn
    man sich zunächst auf die Aktualisierung der Mittelwerte konzentriert.

    Die anderen Parameter, zum Beispiel \(\pi_c\) und \(\Sigma_c\), können dabei
    zunächst als fix angenommen werden.
\end{notebox}

\begin{solutionbox}
    Der Algorithmus läuft iterativ:

    \begin{enumerate}
        \item Initialisiere die Mittelwerte \(\mu_c\).
        \item Berechne für jeden Datenpunkt und jede Komponente die
        Responsibilities:
        \[
            \gamma_{nc}
            =
            p(c\mid x^{(n)},\theta).
        \]

        \item Aktualisiere die Mittelwerte:
        \[
            \mu_c
            \leftarrow
            \frac{
                \sum_{n=1}^{N}
                \gamma_{nc}x^{(n)}
            }{
                \sum_{n=1}^{N}
                \gamma_{nc}
            }.
        \]

        \item Wiederhole E-Schritt und M-Schritt bis zur Konvergenz.
    \end{enumerate}
\end{solutionbox}

\begin{answerbox}
    Ein erster EM-Algorithmus:

    \[
        \boxed{
            \text{E-Schritt: }
            \gamma_{nc}
            =
            p(c\mid x^{(n)},\theta)
        }
    \]

    \[
        \boxed{
            \text{M-Schritt: }
            \mu_c
            \leftarrow
            \frac{
                \sum_{n=1}^{N}
                \gamma_{nc}x^{(n)}
            }{
                \sum_{n=1}^{N}
                \gamma_{nc}
            }
        }
    \]

    EM wechselt also zwischen Schätzen der latenten Zuordnung und Aktualisieren
    der Parameter.
\end{answerbox}

%%%%%%%%%%%%%%%%%%%%%%%%%%%%%%%%%%%%%%%%%%%%%%%%%%%%%%%%%%%%%%%%%%%
\subsection{The Generality of the EM Algorithm}
%%%%%%%%%%%%%%%%%%%%%%%%%%%%%%%%%%%%%%%%%%%%%%%%%%%%%%%%%%%%%%%%%%%

\begin{notebox}
    Der EM-Algorithmus ist nicht auf Gaussian Mixture Models beschränkt.

    Er ist ein allgemeines Verfahren für Modelle mit latenten Variablen.

    Die Struktur ist:

    \[
        \text{beobachtete Daten } X,
        \qquad
        \text{latente Variablen } Z,
        \qquad
        \text{Parameter } \theta.
    \]
\end{notebox}

\begin{solutionbox}
    In vielen Modellen wäre die Maximum-Likelihood-Schätzung einfach, wenn die
    latenten Variablen bekannt wären.

    Da sie aber unbekannt sind, arbeitet EM mit Posterior-Verteilungen über die
    latenten Variablen.

    Der E-Schritt berechnet die aktuelle Unsicherheit über \(Z\).

    Der M-Schritt aktualisiert die Parameter so, als würde man diese weichen
    Zuordnungen als Gewichte verwenden.
\end{solutionbox}

\begin{answerbox}
    EM ist allgemein ein Algorithmus für Modelle der Form:

    \[
        \boxed{
            p(x,z\mid \theta)
        }
    \]

    wobei \(x\) beobachtet und \(z\) latent ist.

    Ziel ist die Maximierung der marginalen Likelihood:

    \[
        \boxed{
            p(x\mid \theta)
            =
            \sum_z p(x,z\mid \theta)
        }
    \]

    beziehungsweise im kontinuierlichen Fall:

    \[
        \boxed{
            p(x\mid \theta)
            =
            \int p(x,z\mid \theta)\,dz.
        }
    \]
\end{answerbox}

%%%%%%%%%%%%%%%%%%%%%%%%%%%%%%%%%%%%%%%%%%%%%%%%%%%%%%%%%%%%%%%%%%%
\subsection{Derivations Do Not Depend on Specific Distribution Forms}
%%%%%%%%%%%%%%%%%%%%%%%%%%%%%%%%%%%%%%%%%%%%%%%%%%%%%%%%%%%%%%%%%%%

\begin{notebox}
    Eine wichtige Eigenschaft von EM ist, dass die grundlegende Herleitung nicht
    von der konkreten Form der Verteilungen abhängt.

    Die Herleitung verwendet nur:

    \begin{itemize}
        \item latente Variablen,
        \item Marginalisierung,
        \item Jensen's Ungleichung,
        \item Posterior-Verteilungen,
        \item Maximierung eines Hilfsobjectives.
    \end{itemize}
\end{notebox}

\begin{solutionbox}
    Für GMMs sind die latenten Variablen diskrete Clusterindikatoren.

    In anderen Modellen können die latenten Variablen aber anders aussehen:

    \begin{itemize}
        \item diskrete Klassen,
        \item kontinuierliche latente Variablen,
        \item fehlende Daten,
        \item Sequenzen verborgener Zustände,
        \item Mischkomponenten,
        \item versteckte Ursachen.
    \end{itemize}

    Die konkrete Form der E- und M-Schritte hängt vom Modell ab.

    Das Prinzip bleibt aber gleich.
\end{solutionbox}

\begin{answerbox}
    Die EM-Herleitung ist allgemein, weil sie auf der Struktur

    \[
        \boxed{
            \log p(x\mid \theta)
            =
            \log \sum_z p(x,z\mid \theta)
        }
    \]

    beziehungsweise

    \[
        \boxed{
            \log p(x\mid \theta)
            =
            \log \int p(x,z\mid \theta)\,dz
        }
    \]

    basiert.

    Sie hängt nicht speziell davon ab, dass die Komponenten Gauß-Verteilungen
    sind.
\end{answerbox}

%%%%%%%%%%%%%%%%%%%%%%%%%%%%%%%%%%%%%%%%%%%%%%%%%%%%%%%%%%%%%%%%%%%
\subsection{Forms for Different Latent Variables}
%%%%%%%%%%%%%%%%%%%%%%%%%%%%%%%%%%%%%%%%%%%%%%%%%%%%%%%%%%%%%%%%%%%

\begin{notebox}
    Je nach Modell kann die latente Variable verschiedene Formen haben.

    Die EM-Idee bleibt dieselbe, aber Summen und Integrale ändern sich abhängig
    vom Typ der latenten Variable.
\end{notebox}

\begin{solutionbox}
    Für diskrete latente Variablen verwendet man Summen:

    \[
        p(x\mid \theta)
        =
        \sum_z
        p(x,z\mid \theta).
    \]

    Für kontinuierliche latente Variablen verwendet man Integrale:

    \[
        p(x\mid \theta)
        =
        \int
        p(x,z\mid \theta)\,dz.
    \]

    Für mehrere latente Variablen können entsprechend mehrfache Summen oder
    Integrale auftreten.
\end{solutionbox}

\begin{answerbox}
    Die Form hängt vom Latenttyp ab:

    \[
        \boxed{
            \text{diskretes } z:
            \quad
            p(x\mid \theta)
            =
            \sum_z p(x,z\mid \theta)
        }
    \]

    \[
        \boxed{
            \text{kontinuierliches } z:
            \quad
            p(x\mid \theta)
            =
            \int p(x,z\mid \theta)\,dz
        }
    \]

    EM bleibt aber strukturell gleich.
\end{answerbox}

%%%%%%%%%%%%%%%%%%%%%%%%%%%%%%%%%%%%%%%%%%%%%%%%%%%%%%%%%%%%%%%%%%%
\subsection{Definition of the EM Meta-Algorithm for Generative Models}
%%%%%%%%%%%%%%%%%%%%%%%%%%%%%%%%%%%%%%%%%%%%%%%%%%%%%%%%%%%%%%%%%%%

\begin{notebox}
    Der EM-Algorithmus kann als Meta-Algorithmus für generative Modelle mit
    latenten Variablen formuliert werden.

    Gegeben ist ein Modell:

    \[
        p(x,z\mid \theta).
    \]

    Beobachtet wird nur \(x\), nicht aber \(z\).
\end{notebox}

\begin{calculationbox}
    Die marginale Log-Likelihood lautet:

    \[
        \log p(X\mid \theta)
        =
        \sum_{n=1}^{N}
        \log
        \sum_z
        p(x^{(n)},z\mid \theta).
    \]

    EM maximiert stattdessen iterativ die Free Energy:

    \[
        \mathcal{F}(q,\theta)
        =
        \sum_{n=1}^{N}
        \sum_z
        q^{(n)}(z)
        \log
        \frac{
            p(x^{(n)},z\mid \theta)
        }{
            q^{(n)}(z)
        }.
    \]
\end{calculationbox}

\begin{solutionbox}
    Der EM-Meta-Algorithmus:

    \begin{enumerate}
        \item Initialisiere die Parameter \(\theta^{old}\).

        \item \textbf{E-Schritt:}
        Berechne die Posterior-Verteilung der latenten Variablen unter den
        aktuellen Parametern:
        \[
            q^{(n)}(z)
            =
            p(z\mid x^{(n)},\theta^{old}).
        \]

        \item \textbf{M-Schritt:}
        Aktualisiere die Parameter durch Maximierung der Free Energy:
        \[
            \theta^{new}
            =
            \operatorname*{arg\,max}_{\theta}
            \mathcal{F}(q,\theta).
        \]

        \item Wiederhole E- und M-Schritt bis zur Konvergenz.
    \end{enumerate}
\end{solutionbox}

\begin{answerbox}
    Der EM-Meta-Algorithmus lautet:

    \[
        \boxed{
            \text{E-Schritt: }
            q^{(n)}(z)
            =
            p(z\mid x^{(n)},\theta^{old})
        }
    \]

    \[
        \boxed{
            \text{M-Schritt: }
            \theta^{new}
            =
            \operatorname*{arg\,max}_{\theta}
            \mathcal{F}(q,\theta)
        }
    \]

    mit

    \[
        \boxed{
            \mathcal{F}(q,\theta)
            =
            \sum_{n}
            \sum_z
            q^{(n)}(z)
            \log
            \frac{
                p(x^{(n)},z\mid \theta)
            }{
                q^{(n)}(z)
            }
        }.
    \]
\end{answerbox}

%%%%%%%%%%%%%%%%%%%%%%%%%%%%%%%%%%%%%%%%%%%%%%%%%%%%%%%%%%%%%%%%%%%
\subsection{Prüfungsrelevante Kurzfassung}
%%%%%%%%%%%%%%%%%%%%%%%%%%%%%%%%%%%%%%%%%%%%%%%%%%%%%%%%%%%%%%%%%%%

\begin{answerbox}
    Lecture 9 behandelt GMMs und die allgemeine Idee des
    Expectation-Maximization-Algorithmus.

    Die wichtigsten Punkte sind:

    \begin{itemize}
        \item Ein GMM hat die Form:
        \[
            p(x)
            =
            \sum_{c=1}^{K}
            \pi_c
            \mathcal{N}(x\mid \mu_c,\Sigma_c).
        \]

        \item Die Log-Likelihood enthält einen schwierigen Term:
        \[
            \log
            \sum_c
            p(x^{(n)},c\mid \theta).
        \]

        \item Die Free Energy ist eine untere Schranke:
        \[
            \log p(X\mid \theta)
            \geq
            \mathcal{F}(q,\theta).
        \]

        \item Die Free Energy lautet:
        \[
            \mathcal{F}(q,\theta)
            =
            \sum_n
            \sum_c
            q^{(n)}(c)
            \log
            \frac{
                p(x^{(n)},c\mid \theta)
            }{
                q^{(n)}(c)
            }.
        \]

        \item Die optimale Wahl von \(q\) ist der Posterior:
        \[
            q^{(n)}(c)
            =
            p(c\mid x^{(n)},\theta).
        \]

        \item Für GMMs sind die Responsibilities:
        \[
            \gamma_{nc}
            =
            \frac{
                \pi_c
                \mathcal{N}(x^{(n)}\mid \mu_c,\Sigma_c)
            }{
                \sum_j
                \pi_j
                \mathcal{N}(x^{(n)}\mid \mu_j,\Sigma_j)
            }.
        \]

        \item Das Mittelwertupdate ist:
        \[
            \mu_c^{new}
            =
            \frac{
                \sum_n
                \gamma_{nc}x^{(n)}
            }{
                \sum_n
                \gamma_{nc}
            }.
        \]

        \item EM ist allgemein für Modelle mit latenten Variablen:
        \[
            p(x,z\mid \theta).
        \]

        \item E-Schritt:
        \[
            q^{(n)}(z)
            =
            p(z\mid x^{(n)},\theta^{old}).
        \]

        \item M-Schritt:
        \[
            \theta^{new}
            =
            \operatorname*{arg\,max}_{\theta}
            \mathcal{F}(q,\theta).
        \]
    \end{itemize}

    Die Kernaussage lautet:

    \[
        \boxed{
            \text{EM optimiert Modelle mit latenten Variablen durch
            abwechselnde Posterior-Berechnung und Parameteraktualisierung.}
        }
    \]
\end{answerbox}

\end{document}